%!TEX TS-program = xelatex
%!TEX root = ../Topo2_WS15/topologie_2.tex
\RequirePackage{fix-cm} 
\documentclass[a4paper, twoside, headsepline, index=totoc,toc=listof, fontsize=10pt, cleardoublepage=empty, headinclude, DIV=12, BCOR=5mm, titlepage,draft]{scrartcl}
\usepackage{scrtime} % KOMA, Uhrzeit ermoeglicht

\usepackage{etoolbox}
\usepackage{letltxmacro}
\usepackage{ifthen}

%--Farbdefinitionen
%-- muss vor tikz geladen werden
\usepackage[usenames, table, x11names]{xcolor}
\definecolor{dark_gray}{gray}{0.45}
\definecolor{light_gray}{gray}{0.6}
\usepackage[final]{graphicx}

%--Zum Zeichnen
%-- muss vor polyglossia bzw. babel geladen werden
\usepackage{tikz}
\usepackage{tikz-cd}
\usetikzlibrary{external}
\tikzset{>=latex}
\usetikzlibrary{shapes,arrows.meta,intersections}
\usetikzlibrary{calc,3d}
\usetikzlibrary{decorations.pathreplacing,decorations.markings, decorations.pathmorphing}
\usetikzlibrary{angles}

%-- Konfiguration von tikzexternalize
\tikzexternalize[prefix=tikz/,up to date check=diff]
\pgfkeys{/pgf/images/include external/.code=\includegraphics{#1}}
\tikzset{external/system call={xelatex \tikzexternalcheckshellescape -halt-on-error -interaction=batchmode --shell-escape -jobname "\image" "\texsource"}}


%-- tikzexternalize fuer tikzcd deaktivieren, da inkompatibel
\AtBeginEnvironment{tikzcd}{\tikzexternaldisable}
\AtEndEnvironment{tikzcd}{\tikzexternalenable}

%-- um Inkompatibilitaeten von quotes und polyglossia bzw. babel zu vermeiden
\tikzset{
  every picture/.append style={
    execute at begin picture={\shorthandoff{"}},
    execute at end picture={\shorthandon{"}}
  }
}
\usetikzlibrary{quotes}
\usepackage{pgfplots}
\usepgfplotslibrary{colormaps}
% \pgfplotsset{compat=1.12}



%-- Mathesymbole etc.
\usepackage{mathtools} % beinhaltet amsmath
\mathtoolsset{showonlyrefs, centercolon,showmanualtags}
\newtagform{brackets}[\textbf]{[}{]}
\usetagform{brackets}
\usepackage{fix-cm}
% \usepackage{wasysym} % zusätzliche Symbole
\usepackage{amssymb} % zusätzliche Symbole
% \usepackage{stmaryrd} % für Blitz
\usepackage{nicefrac} % schräge Brüche
\usepackage{faktor}
\newcommand{\Faktor}[1]{\faktor[\textstyle]{#1}}
\usepackage{xfrac}
\usepackage{cancel}
\usepackage{mathdots} % Verbesserung von Punkten wie zB \ldots
\usepackage[bb=px]{mathalfa}

%-- einzelne Symbole aus dem mathx-package
% \DeclareFontFamily{U}{mathx}{\hyphenchar\font45}
% \DeclareFontShape{U}{mathx}{m}{n}{<-> mathx10}{}
% \DeclareSymbolFont{mathx}{U}{mathx}{m}{n}
% \DeclareMathSymbol{\bigplus}{1}{mathx}{"90}

%-- 
\usepackage{silence}
\WarningFilter{latexfont}{Size substitutions with differences}
\WarningFilter{latexfont}{Font shape `U/bbold/m/n' in size}

\DeclareSymbolFont{bbold}{U}{bbold}{m}{n}
\DeclareSymbolFontAlphabet{\mathbbold}{bbold}
\newcommand{\ind}{\mathbbold{1}} % charakteristische-Funktion-Eins

\def\mathul#1#2{\color{#1}\underline{{\color{black}#2}}\color{black}} %farbiges Untersteichen im Mathe-Modus
\renewcommand{\le}{\leqslant}
\renewcommand{\ge}{\geqslant}

%-- Underbrace als Befehl in LaTeX-Syntax (und ohne Spacingprobleme mit nachfolgenden Operatoren...)
\newcommand{\Underbrace}[2]{{\underbrace{#1}_{#2}}}
\newcommand{\Underbracket}[2]{{\underbracket[0.7pt][2pt]{#1}_{#2}}}

%-- Alles was mit Schrift und XeTeX zu tun hat
\usepackage{palatino}
\usepackage[euler-digits]{eulervm}
\usepackage[no-math]{fontspec}
\usepackage{polyglossia} % moderner babel-ersatz
\setmainlanguage[spelling=new,babelshorthands=true]{german}
\shorthandoff{"}
\setotherlanguage{english}
\defaultfontfeatures{Mapping=tex-text, WordSpace={1.2}, Ligatures={Required,Common,Contextual}} %

%-- Alle hier aufgeführten Schriftarten sind in jeder TeXlive-Installation vorhanden!!
%-- Hauptschriftart festlegen
% \setmainfont{texgyrepagella}[Extension=.otf,UprightFont=*-regular,BoldFont=*-bold,ItalicFont=*-italic,BoldItalicFont=*-bolditalic,ItalicFeatures={Style=Historic}]
\setmainfont{TeXGyrePagellaX}[Extension=.otf,UprightFont=*-Regular,BoldFont=*-Bold,ItalicFont=*-Italic,BoldItalicFont=*-BoldItalic,ItalicFeatures={Style=Historic},Ligatures={Required,Common,Contextual,Historic}]
% \setmainfont[BoldFont={* Bold}, ItalicFont={* Light Italic},SmallCapsFont={Linux Libertine O}, SmallCapsFeatures={Letters=SmallCaps}]{Source Sans Pro}


%-- Serifenlose Schriftart festlegen
% \setsansfont{Roboto}[Scale=MatchUppercase,Extension=.ttf,UprightFont=*-Regular,BoldFont=*-Bold,ItalicFont=*-RegularItalic,BoldItalicFont=*-MediumItalic]
% \setsansfont{RobotoSlab}[Scale=MatchUppercase,Extension=.ttf,UprightFont=*-Regular,BoldFont=*-Bold,ItalicFont=*-Thin,SmallCapsFont=*-Regular]
% \setsansfont{FiraSans}[Scale=MatchUppercase,Extension=.otf, UprightFont=*-Regular, BoldFont=*-SemiBold, ItalicFont=*-RegularItalic, BoldItalicFont=*-MediumItalic]
% \setsansfont{SourceSansPro}[Scale=MatchUppercase,Extension=.otf,UprightFont=*-Regular,BoldFont=*-Semibold,ItalicFont=*-LightIt,BoldItalicFont=*-SemiboldIt]
% \setsansfont{texgyreadventor}[Scale=MatchUppercase, Extension=.otf, UprightFont=*-regular, BoldFont=*-bold, ItalicFont=*-italic,BoldItalicFont=*-bolditalic]
% \setsansfont{Kurier}[Scale=MatchUppercase, Extension=.otf, UprightFont=*Medium-Regular, BoldFont=*Heavy-Regular, ItalicFont=*Medium-Italic,BoldItalicFont=*Heavy-Italic]
% \setsansfont{Iwona}[Scale=MatchUppercase, Extension=.otf, UprightFont=*Medium-Regular, BoldFont=*Heavy-Regular, ItalicFont=*Medium-Italic,BoldItalicFont=*Heavy-Italic]
% \setsansfont{LinBiolinum}[Scale=MatchUppercase, Extension=.otf, UprightFont=*_R, BoldFont=*_RB, ItalicFont=*_RI,BoldItalicFont=*_RBO]
% \setsansfont{texgyrebonum}[Scale=MatchUppercase,Extension=.otf, UprightFont=*-regular, BoldFont=*-bold, ItalicFont=*-italic, BoldItalicFont=*-bolditalic]
\setsansfont{texgyreadventor}[Scale=MatchUppercase,Extension=.otf, UprightFont=*-regular, BoldFont=*-bold, ItalicFont=*-italic, BoldItalicFont=*-bolditalic]



% \newfontfamily{\geometric}{texgyreheros}[Scale=MatchLowercase, Extension=.otf, UprightFont=*-regular, BoldFont=*-bold, ItalicFont=*-italic]
% \addtokomafont{disposition}{\geometric}
\setmonofont{SourceCodePro}[Scale=0.9,Extension=.otf,UprightFont=*-Regular, BoldFont=*-Semibold, ItalicFont=*-Light]
\usepackage{xltxtra}
\usepackage{fontawesome}
\usepackage[final]{microtype}


\usepackage[draft=false]{scrlayer-scrpage} % wie fancyhdr, nur optimiert auf KOMA-Skript, leicht andere Syntax

\flushbottom


%--Mixed
\usepackage[shortlabels,inline]{enumitem}
\setlist[itemize,1]{label=\faCaretRight}
\setlist[enumerate]{font=\bfseries}
\setlist[description]{font=\normalfont\bfseries}
\usepackage[german=quotes]{csquotes}
\usepackage{makeidx}
\usepackage{wrapfig}
\usepackage{float}
\usepackage[margin=10pt, font=small, labelfont={sf, bf}, format=plain, indention=1em]{caption}
\captionsetup[wrapfigure]{name=Abb. }
\usepackage{stackrel}


%--Unterstreichung
\usepackage[normalem]{ulem}
\setlength{\ULdepth}{1.8pt}

%--Indexverarbeitung
\newcommand{\bet}[1]{\emph{#1}}
\newcommand{\Index}[1]{\bet{#1}\index{#1}}
\makeindex
\setindexpreamble{{\noindent\sffamily\small Die \emph{Seitenzahlen} sind mit Hyperlinks versehen und somit anklickbar}  \par \bigskip}
\renewcommand{\indexpagestyle}{scrheadings}


%--Marginnote & todonotes
\deffootnote[1.5em]{1.5em}{1.5em}{\textsuperscript{\thefootnotemark}\ }
\usepackage[fulladjust]{marginnote}
\renewcommand*{\marginfont}{\color{dark_gray} \itshape \footnotesize}
\usepackage[textsize=small]{todonotes}
\usepackage{ragged2e}
\renewcommand*{\raggedleftmarginnote}{\RaggedLeft}
\renewcommand*{\raggedrightmarginnote}{\RaggedRight}

%--Todonotes mit tikz/externalize kompatibel machen
\LetLtxMacro{\oldtodo}{\todo}
\renewcommand{\todo}[2][]{\tikzexternaldisable\oldtodo[#1]{#2}\tikzexternalenable}
\LetLtxMacro{\oldmissingfigure}{\missingfigure}
\renewcommand{\missingfigure}[2][]{\tikzexternaldisable\oldmissingfigure[{#1}]{#2}\tikzexternalenable}

%--Konfiguration von Hyperref pdfstartview=FitH, 
\usepackage[hidelinks, pdfpagelabels,  bookmarksopen=true, bookmarksnumbered=true, linkcolor=black, urlcolor=SkyBlue2, plainpages=false,pagebackref, citecolor=black, hypertexnames=true, pdfauthor={Jannes Bantje}, pdfborderstyle={/S/U}, linkbordercolor=SkyBlue2, colorlinks=false,final]{hyperref}


\newcommand{\appendLink}[1]{#1\,\faExternalLink}
\newcommand{\hrefsym}[2]{\href{#1}{\texttt{\appendLink{#2}}}}
\newcommand{\hrefsymmail}[2]{\href{#1}{\texttt{\faEnvelopeO\,#2}}}
\renewcommand{\url}[1]{\hrefsym{#1}{\nolinkurl{#1}}}



% -- QR-Codes
\usepackage{qrcode} %-- hinter hyperref laden!

%--Römische Zahlen
\newcommand{\RM}[1]{\MakeUppercase{\romannumeral #1{}}}

%%--Abkürzungen etc.
\newcommand{\light}{\text{\Large $\lightning$}}

%-- Definitionen von weiteren Mathe-Befehlen

\DeclareMathOperator{\im}{im} %Bild
\DeclareMathOperator{\id}{id} %identische Abbildung
\DeclareMathOperator{\conj}{conj} %Konjugation
\DeclareMathOperator{\sgn}{sgn} %Signum
\DeclareMathOperator{\End}{End} %Endomorphismen
\DeclareMathOperator{\Hom}{Hom} %Homomorphismen
\DeclareMathOperator{\Iso}{Iso} %Isomorphismen
\DeclareMathOperator{\Span}{span} %Span

\DeclareMathOperator{\grad}{grad} %Gradient
\DeclareMathOperator{\dive}{div} %Gradient
\DeclareMathOperator{\rot}{rot} %Rotation


\DeclareMathOperator{\Char}{char} %Charakteristik
\DeclareMathOperator{\Aut}{Aut} %Automorphismen

% -- Zum Finetuning von Befehlen
\makeatletter
\newcommand{\raisemath}[1]{\mathpalette{\raisem@th{#1}}}
\newcommand{\raisem@th}[3]{\raisebox{#1}{$#2#3$}}
\makeatother

\newcommand{\sing}{{\raisemath{1.1pt}{\scriptscriptstyle\mathrm{sing}}}}
\newcommand{\pt}{\mathrm{pt}}
\newcommand{\op}{\mathrm{op}}
\DeclareMathOperator{\Sp}{Sp}
\DeclareMathOperator{\Keg}{Keg}
\newcommand{\slashedi}{i\hspace{-3.5pt}/}
\newcommand{\cupp}{\smallsmile}
\newcommand{\capp}{\smallfrown}
\DeclareMathOperator*{\colim}{colim}


\newcommand{\D}{\ensuremath{\mathrm{D}\mkern-1.5mu}}
\newcommand{\mathd}{\ensuremath{\mathrm{d}\mkern-1.0mu}}
\newcommand{\Tmap}{\ensuremath{\mathrm{T}\mkern-0.85mu}}
\newcommand{\opL}{\ensuremath{\mathrm{L}\mkern-0.6mu}}


\DeclareMathOperator{\supp}{supp} %Träger


% -- Befehle für FunkAna und Operatoralgebren
\newcommand{\tr}{\mathrm{tr}}
\newcommand{\w}{\mathrm{w}}
\newcommand{\sa}{\mathrm{sa}}
\newcommand{\so}{\mathrm{s.o.}}
\newcommand{\weakT}[1]{\ensuremath{\mathcal{T}_{#1}^{\mkern+1.0mu\text{\raisebox{0.4ex}{$\mathrm{w}$}}}}}
\newcommand{\weakTstar}[1]{\ensuremath{\mathcal{T}_{#1}^{\mkern+1.0mu\text{\raisebox{0.4ex}{$\mathrm{w}$}}^*}}}
\newcommand{\finSub}{\subset\mkern-0.7mu \subset}
\DeclareMathOperator{\Inv}{Inv}
\DeclareMathOperator{\dist}{dist}
\newcommand{\simm}{{\hspace{-1.6pt}\raisemath{0.5pt}{\sim}}}
\DeclareMathOperator{\ev}{ev}
\DeclareMathOperator{\Alg}{Alg}


% -- Kategorien
\newcommand{\TOP}{\textsc{Top}}
\newcommand{\HTOP}{\textsc{HTop}}
\newcommand{\VR}{\textsc{VR}}
\newcommand{\MOD}{\textsc{Mod}}
\newcommand{\MONOIDE}{\textsc{Monoide}}
\newcommand{\MENGEN}{\textsc{Mengen}}
\newcommand{\MAN}{\textsc{Man}}
\newcommand{\GRUPPEN}{\textsc{Gruppen}}
\newcommand{\ABELGRUPPEN}{\textsc{Abel.Gruppen}}
\newcommand{\KAT}{\textsc{Kat}}
\newcommand{\FUN}{\textsc{Fun}}
\newcommand{\SIMP}{\textsc{Simp}}
\newcommand{\VEKT}{\textsc{Vekt}}

\DeclareMathOperator{\cell}{cell}
\DeclarePairedDelimiter{\homologieklasse}{\llbracket}{\rrbracket}
\newcommand{\rand}[1]{\ensuremath{\partial^{\scriptscriptstyle #1}}}


% -- theorem packages
\usepackage{amsthm}
\usepackage{thmtools,thm-restate}
\renewcommand{\listtheoremname}{Übersicht aller Aussagen}

% -- Theoreme als PDF-Lesezeichen
\usepackage{bookmark}
\bookmarksetup{open,numbered}
\makeatletter
\newcommand*{\theorembookmark}{%
  \bookmark[
    dest=\@currentHref,
    rellevel=1,
    keeplevel,
  ]{%
    \thmt@thmname\space\csname the\thmt@envname\endcsname
    \ifx\thmt@shortoptarg\@empty
    \else
      \space(\thmt@shortoptarg)%
    \fi
  }%
}   
\makeatother

% -- Definition der einzelnen Umgebungen
\declaretheoremstyle[%
	headfont=\sffamily\bfseries,
	notefont=\normalfont\sffamily\scshape,
	bodyfont=\normalfont,
	headformat=\NUMBER\ \NAME\NOTE,
	headpunct=.,
	postheadspace=1em,
	spaceabove=15pt,spacebelow=10pt,
	postheadhook=\theorembookmark]%
{mainstyle}
\declaretheoremstyle[%
	headfont=\sffamily\bfseries,
	notefont=\normalfont\sffamily\scshape,
	bodyfont=\normalfont,
	headformat=\NUMBER\NAME\NOTE,
	headpunct=.,
	postheadspace=1em,
	spaceabove=15pt,spacebelow=10pt,
	postheadhook=\theorembookmark]%
{mainstyle_unnumbered}
\declaretheoremstyle[%
	headfont=\sffamily\bfseries,
	notefont=\normalfont\sffamily\scshape,
	bodyfont=\normalfont,
	headformat=swapnumber,
	headpunct=.,
	postheadspace=1em,
	spaceabove=15pt,spacebelow=10pt,
	postheadhook=\theorembookmark,
	qed=\qedsymbol]%
{mainstyleB}
\declaretheorem[name=Definition,parent=section,style=mainstyle]{definition}
\declaretheorem[name=Definition,numbered=no,style=mainstyle_unnumbered]{definition*}
\declaretheorem[name=Theorem,sharenumber=definition,style=mainstyle]{theorem}
\declaretheorem[name=Theorem,numbered=no,style=mainstyle_unnumbered]{theorem*}
\declaretheorem[name=Proposition,sharenumber=definition,style=mainstyle]{proposition}
\declaretheorem[name=Lemma,sharenumber=definition,style=mainstyle]{lemma}
\declaretheorem[name=Satz,sharenumber=definition,style=mainstyle]{satz}
\declaretheorem[name=Korollar,sharenumber=definition,style=mainstyle]{korollar}
\declaretheorem[name=Korollar,sharenumber=definition,style=mainstyleB]{korollarB}
\declaretheorem[name=Frage,numbered=no,style=mainstyle_unnumbered]{frage}
\declaretheorem[name=Erinnerung,sharenumber=definition,style=mainstyle]{erinnerungA}


% -- Beweise
\declaretheoremstyle[headfont=\bfseries\scshape,bodyfont=\normalfont,headpunct=:,postheadspace=1em,spacebelow=12pt,spaceabove=2pt,qed=\qedsymbol]{beweise}
\declaretheoremstyle[headfont=\sffamily\bfseries,bodyfont=\normalfont,headpunct=:,postheadspace=1em,spacebelow=10pt,spaceabove=10pt]{bemerkungen}
\declaretheorem[name=Beweis,numbered=no,style=beweise]{beweis}
\declaretheorem[name=Bemerkung,sharenumber=definition,style=mainstyle]{bemerkung}
\declaretheorem[name=Beispiel,sharenumber=definition,style=mainstyle]{beispiel}
\declaretheorem[name=Übung,numbered=no,style=bemerkungen]{uebung}
\declaretheorem[name=Erinnerung,numbered=no,style=bemerkungen]{erinnerung}


%--Skalarprodukt
\DeclarePairedDelimiterX\sprod[2]{\langle}{\rangle}{#1\,\delimsize\vert\,#2}
\DeclarePairedDelimiterX\skal[2]{\langle}{\rangle}{#1\,,\,#2}

%--Betrag, Gaußklammer
\DeclarePairedDelimiter{\abs}{\lvert}{\rvert}
\DeclarePairedDelimiter{\floor}{\lfloor}{\rfloor}
\DeclarePairedDelimiter{\ceil}{\lceil}{\rceil}

%--Norm
\DeclarePairedDelimiter\norm{\Vert}{\Vert}


%--Umklammern
\DeclarePairedDelimiter\enbrace{(}{)}
\DeclarePairedDelimiter\benbrace{[}{]}
\DeclarePairedDelimiter\lenbrace{<}{>}
\newcommand{\ssbrace}[1]{{\scriptscriptstyle\enbrace{#1}}}

%--Mengen
\newcommand\SetSymbol[1][]{\nonscript\:#1\vert\allowbreak\nonscript\:\mathopen{}}
\providecommand\given{} % to make it exist
\DeclarePairedDelimiterX\set[1]\{\}{\renewcommand\given{\SetSymbol[\delimsize]}#1}

%--Abbildungsdefinition
\newcommand{\mapdef}[5]{%
	\[
		\begin{array}{rcl}
			\textstyle #1 &\xrightarrow{\minwidthbox{#5}{2em}} & \textstyle #2 \\[0.5ex]
			\textstyle #3 &\xmapsto{\minwidthbox{\mbox{ }}{2em}} & \textstyle #4
		\end{array}
	\]
}

%--modifiziertes Stackrel
\newcommand{\StackText}[2]{\stackrel{\mbox{\scriptsize #1}}{#2}}

%--Befehl für eine Box mit Mindestweite
\DeclareRobustCommand{\minwidthbox}[2]{%
  \ifmmode
    \expandafter\mathmakebox
  \else
    \expandafter\makebox
  \fi
  [\ifdim#2<\width\width\else#2\fi]{#1}%
}

%--Befehl zum Einkreisen
\newcommand*\circled[1]{\tikzexternaldisable\tikz[baseline=(char.base)]{\node[shape=circle,draw,inner sep=2pt] (char) {#1};}\tikzexternalenable}


%--Differential
\newcommand{\diff}[2]{\ensuremath{\frac{{\partial #1}}{{\partial #2}} }}
\newcommand{\diffd}[2]{\ensuremath{\frac{\mathd #1}{\mathd #2} }}

%--Inhaltsverzeichnis
\usepackage[tocindentauto]{tocstyle}
\usetocstyle{KOMAlike}	
\shorthandon{"}
		
%--Punkte (als letztes laden)
\usepackage{ellipsis}

\newcommand{\fach}{Operatoralgebren}
\newcommand{\semester}{WiSe 2015}
\newcommand{\homepage}{https://wwwmath.uni-muenster.de/u/wilhelm.winter/wwinter/operatoralgebren.html}

\newcommand{\prof}{Prof.\ Dr.\ Wilhelm Winter}
\publishers{\scalebox{11}{\Huge$C^*$}}
\providecommand{\verfasser}{Jannes Bantje}

%--Konfiguration von scrheadings
\setheadsepline{1pt}[\color{light_gray}]
\pagestyle{scrheadings}
%\fancyhf[H, F]{}
\clearscrheadfoot

\providecommand{\shortFach}{\fach}
\lehead{ \includegraphics[height=0.5 cm,keepaspectratio]{../!config/Bilder/Logo_WWU_Muenster_light_gray.pdf}}
\rehead{\rule{0cm}{0.5cm}\footnotesize \sffamily \color{light_gray} \verfasser -- Mitschrift \shortFach}
\lohead{\rule{0cm}{0.5cm} \footnotesize \sffamily \color{light_gray} Stand: \today \; \thistime[:]}
\rohead{\includegraphics[height=0.5 cm,keepaspectratio]{../!config/Bilder/fb10logo_gray.pdf}}


\ofoot[{ \color{dark_gray} \LARGE \sffamily \thepage}]{{ \color{dark_gray} \LARGE \sffamily \thepage}} %hier wir auch der plain Stil bearbeitet!
\automark{section}
\ifoot{ \color{dark_gray} \small \leftmark}

%--Metadaten
\author{\verfasser}
\titlehead{\includegraphics[height=1.5cm, keepaspectratio]{../!config/Bilder/Logo_WWU_Muenster.pdf}%
\hfill \includegraphics[height=1.3cm, keepaspectratio]{../!config/Bilder/fb10logo.pdf}}
\title{Skript \fach}
\subtitle{Mitschrift der Vorlesung  \enquote{\fach} von \prof}

\publishers{
Aktuelle Version verfügbar bei: \bigskip\\
\normalsize
\begin{minipage}[t]{0.45\textwidth}
	\centering{
	\qrcode[height=3.5cm, version=4]{https://github.com/JaMeZ-B/latex-wwu}
	\medskip\\
	\includegraphics[height=0.4cm, keepaspectratio]{../!config/Bilder/github_octo.pdf} 
	\includegraphics[height=0.4cm, keepaspectratio]{../!config/Bilder/GitHub_Logo.pdf} (inklusive Sourcecode)\\
	{\small\url{https://github.com/JaMeZ-B/latex-wwu}}}
\end{minipage}
\hfill
\begin{minipage}[t]{0.45\textwidth}
	\centering{
	\qrcode[height=3.5cm, version=4]{btsync://B6WH2DISQ5QVYIRYIEZSF4ZR2IDVKPN3I?n=latex_share}
	\medskip\\
	\includegraphics[height=0.4cm, keepaspectratio]{../!config/Bilder/bt_sync_logo.pdf} 
	{\large \textbf{Bittorrent} Sync} \\
	\texttt{B6WH2DISQ5QVYIRYIEZSF4ZR2IDVKPN3I}
	}
\end{minipage}
}


\begin{document}
\pagenumbering{Roman}
\maketitle
\begin{abstract}
\section*{Vorwort --- Mitarbeit am Skript}
Dieses Dokument ist eine Mitschrift aus der Vorlesung \enquote{\fach, \semester}, gelesen von \prof. Der Inhalt entspricht weitestgehend dem Tafelanschrieb. Für die
Korrektheit des Inhalts übernehme ich keinerlei Garantie! Für Bemerkungen und Korrekturen -- und seien es nur Rechtschreibfehler -- bin ich sehr dankbar. 
Korrekturen lassen sich prinzipiell auf drei Wegen einreichen: 
\begin{itemize}
	\item Persönliches Ansprechen in der Uni, Mails an \hrefsymmail{mailto:\mail}{\mail} (gerne auch mit annotieren PDFs) 
	\item \emph{Direktes} Mitarbeiten am Skript: Den Quellcode poste ich auf GitHub (siehe oben), also stehen vielfältige Möglichkeiten der Zusammenarbeit zur Verfügung:
	Zum Beispiel durch Kommentare am Code über die Website und die Kombination Fork + Pull Request. Wer sich verdient macht oder ein Skript zu einer Vorlesung, die 
	ich nicht besuche, beisteuern will, dem gewähre ich gerne auch Schreibzugriff.
	
	Beachten sollte man dabei, dass dazu ein Account bei \url{github.com} notwendig ist, der allerdings ohne Angabe von persönlichen Daten angelegt werden kann. 
	Wer bei GitHub (bzw. dem zugrunde liegenden Open-Source-Programm \enquote{\texttt{git}}) -- verständlicherweise -- Hilfe beim Einstieg braucht, dem helfe ich gerne 
	weiter. Es gibt aber auch zahlreiche empfehlenswerte Tutorials im Internet.\footnote{zB. \url{https://try.github.io/levels/1/challenges/1}, ist auf Englisch, aber dafür 
	interaktives LearningByDoing}
	\item \emph{Indirektes} Mitarbeiten: \TeX-Dateien per Mail verschicken. 
	
	Dies ist nur dann sinnvoll, wenn man einen ganzen Abschnitt ändern möchte (zB. einen alternativen Beweis geben), da ich die Änderungen dann per Hand einbauen muss! Ich freue mich aber auch über solche Beiträge!
\end{itemize}
\section*{Vorlesungshomepage}
{\centering 
\begin{minipage}[c][][c]{\textwidth}
	\centering \qrcode[height=3cm]{\homepage} \medskip\\
	\footnotesize \url{\homepage}
\end{minipage}
\par}

\section*{Anmerkung}
Innerhalb dieser Mitschrift wird man öfter den Ausdruck \enquote{Warum?} finden. Dies sind vom Dozenten bewusst weggelassene Details, die zu verstärktem Mitdenken beim Lesen 
animieren sollen. Oftmals sind dies schon aus vorherigen Semestern bekannte Sachverhalte. Nur an wenigen Stellen habe ich die fehlenden Details hinzugefügt.
\end{abstract}

\tableofcontents
\cleardoubleoddemptypage

\pagenumbering{arabic}
\setcounter{page}{1}
\setcounter{footnote}{0}

\section{Spektraltheorie in Banachalgebren} % (fold)
\label{sec:1}

\begin{definition}[{name=[Banachalgebra]}]
	Eine \Index{Banachalgebra} ist eine vollständige normierte Algebra. Für die Multiplikation gilt also $\norm{a b} \le \norm{a} \cdot \norm{b}$
	Wir betrachten in dieser Vorlesung fast ausschließlich Algebren über $\mathbb{C}$. Eine Banachalgebra heißt \Index{unital}, falls sie ein Einselement $\ind$ mit 
	$\norm{\ind} =1$ enthält. 
\end{definition}

\begin{definition}[{name=[Menge der invertierbaren Elemente und Spektrum]}]
	Für eine unitale $\mathbb{C}$-Algebra $A$ setzen wir 
	\begin{align}
		\Inv(A) &:= \set{a \in A  \text{ invertierbar} } \\
		\sigma(a) &:=\sigma_A(a) := \set[\big]{\lambda \in \mathbb{C} \given \lambda \cdot \ind - a \notin \Inv(A)} \marginnote{das Vorzeichen der Differenz spielt für die Definition der Menge keine Rolle}
	\end{align}
	$\sigma(a)$ heißt \Index{Spektrum} von $a \in A$. 
\end{definition}

\begin{bemerkung}[{name=[Invertierbarkeit von $1 -ab$]},label=bem:13]
	Sei $A$ wie oben und $a,b \in A$. Dann gilt
	\[
		\ind -ab \in \Inv(A) \iff \ind -ba \in \Inv(A)
	\]
\end{bemerkung}
\begin{beweis}
	Durch simples Nachrechnen verifiziert man, dass $(\ind-ba)^{-1} = \ind+ b(\ind- ab)^{-1} a$ gilt.
\end{beweis}

\begin{satz}[label=satz14_poly,{name=[Verträglichkeit des Spektrums mit Polynomen]}]
	Sei $A$ eine unitale $\mathbb{C}$-Algebra und $a \in A$ mit $\sigma(a)\neq \emptyset$. Sei weiter $p \in \mathbb{C}[Z]$ ein Polynom. Dann gilt 
	\[
		\sigma\enbrace[\big]{p(a) }= p \enbrace[\big]{\sigma(a)}
	\]
\end{satz}
\begin{beweis}
	Die Aussage ist trivial, falls $p$ konstant ist. Sei also $p$ nicht konstant. Zu $\mu \in \mathbb{C}$ existieren $\lambda_0, \ldots , \lambda_n \in \mathbb{C}$ mit 
	$\lambda_0 \neq 0$ und $p(z) -\mu = \lambda_0 \cdot (z- \lambda_1) \cdot  \ldots \cdot (z- \lambda_n)$, also
	\[
		p(a) - \mu = \lambda_0 \cdot (a-\lambda_1) \cdot  \ldots \cdot (a- \lambda_n)
	\]
	Damit gilt: $p(a) -\mu$ ist genau dann invertierbar, wenn $a -\lambda_1, \ldots , a-\lambda_n $ invertierbar sind. Also 
	\begin{align}
		\mu \in \sigma \enbrace[\big]{p(a)} &\iff \lambda_i \in \sigma(a) \text{ ist für ein } i \in \set{1,\ldots ,n} \\
		&\iff p(\lambda) -\mu =0 \text{ für ein } \lambda \in \sigma(a) \\
		&\iff p(\lambda)=\mu \text{ für ein } \lambda \in \sigma(a)   
	\end{align}
	Somit folgt dann $\sigma \enbrace*{p(a)} = p \enbrace*{\sigma(a)}$.
\end{beweis}

\begin{satz}[label=satz15_neumann,{name=[Neumannsche Reihe]}]
	Sei $A$ eine unitale Banachalgebra und $a \in A$. Falls $\norm{a}<1$, so ist $\ind-a$ invertierbar und 
	\[
		(\ind-a)^{-1} = \sum_{n=0}^{\infty} a^n \tag*{(\Index{Neumannsche Reihe})}
	\]
\end{satz}
\begin{beweis}
	Die Folge der Partialsummen $\enbrace[\big]{\sum_{n=0}^{k} a^n}_{k \in \mathbb{N}} \subset A$ ist eine Cauchyfolge (warum?), sodass wegen Vollständigkeit die unendliche 
	Summe existiert. Es gilt dann
	\begin{align}
		(\ind-a) \enbrace*{\sum_{n=0}^{\infty} a^n} \xleftarrow{k \to \infty}(\ind-a) \enbrace*{\sum_{n=0}^{k} a^n} = \ind-a^{k+1} \xrightarrow{k \to \infty} \ind
	\end{align}
	Die Summe ist also ein Rechtsinverses. Genauso folgt, dass sie auch ein Linksinverses ist.
\end{beweis}

\begin{korollar}[{name=[Invertierbare Elemente bilden offene Menge]},label=inv_offen]
	Für jede unitale Banachalgebra $A$ ist $\Inv(A) \subset A$ offen.
\end{korollar}
\begin{beweis}
	Falls $a \in \Inv(A)$ und $\norm{b-a} < \norm{a^{-1}}^{-1}$, so gilt $\norm{b a^{-1} -\ind} \le \norm{b-a} \cdot \norm{a^{-1}}<1$.  Also ist auch
	$-\enbrace*{\ind- \enbrace*{ba^{-1} -\ind}} = ba^{-1}$ invertierbar nach \autoref{satz15_neumann} und damit auch $b \in \Inv(A)$, denn $b$ ist ein Linksinverses für 
	$a^{-1}(ba^{-1})^{-1} \in \Inv(A)$.
\end{beweis}

\begin{definition}[{name=[Spektralradius]}]
	Sei $A$ eine unitale Banachalgebra. Für $a \in A$ definieren wir den \Index{Spektralradius} durch
	\[
		r(a) := \sup_{\lambda \in \sigma(a)} \abs{\lambda} 
	\] 
\end{definition}

\begin{proposition}[{name=[Spektrum ist kompakt und Spektralradius durch Norm beschränkt]},label=spec_kompakt]
	Sei $a \in A$ wie oben. Dann ist das Spektrum $\sigma(a)$ kompakt und $r(a) \le \norm{a}$.
\end{proposition}
\begin{beweis}
	Angenommen $\abs{\lambda} > \norm{a}$, dann gilt $\norm{\lambda^{-1} a} < 1$, also ist $\ind- \lambda^{-1} a$ invertierbar nach \autoref{satz15_neumann} und 
	damit auch $\lambda -a $. Also kann $\lambda$ nicht im Spektrum von $a$ enthalten sein. 
	Es folgt, dass $r(a) \le \norm{a}$ ist. 	
	Weiter ist $\mathbb{C} \setminus \sigma(a) \subset \mathbb{C}$ offen, da $\Inv(A) \subset A$ offen ist. Damit muss $\sigma(a)$ nach Heine-Borel kompakt sein.
\end{beweis}

\begin{satz}[{name=[In unitalen Banachalgebren sind die Spektren nie leer]},label=spec_nichtleer]
	Sei $A$ eine unitale Banachalgebra. Für jedes $a \in A$ ist $\sigma(a)$ nichtleer. 
\end{satz}
\begin{beweis}
	Falls $\abs{\lambda}> \norm{a}$, so existiert $\enbrace*{\ind- \lambda^{-1} a}^{-1}$ und $(\lambda-a)^{-1}$ mit der gleichen Begründung wie eben. Es gilt dann
	\begin{align}
		(\lambda-a)^{-1} = \lambda^{-1} \enbrace*{1- \lambda^{-1} a}^{-1} \stackrel{\mbox{\scriptsize\ref{satz15_neumann}}}{=} \lambda^{-1} \enbrace*{\sum_{n=0}^{\infty} 
		\lambda^{-n}a^n}
	\end{align}
	und weiter
	\[
		\norm{(\lambda-a)^{-1}} \le \abs*{\lambda^{-1}} \cdot \sum_{n=0}^{\infty} \enbrace*{\abs*{\lambda^{-1}} \cdot \norm{a}}^n  \xrightarrow{\abs{\lambda} \to \infty} 0 
		\label{19:eq:1} \tag{$*$}
	\]
	da man die Summe als geometrische Reihe mit $2$ abschätzen kann, falls $\abs{\lambda}> 2 \cdot \norm{a}$ ist.
	Sei nun $\lambda_0 \in \mathbb{C}\setminus \sigma(a)$ und $\abs{\lambda-\lambda_0}< \norm{(a- \lambda_0)^{-1}}^{-1}$. Dann gilt 
	\[
		\norm{(\lambda-\lambda_0)(a - \lambda_0)^{-1}} \le \abs{\lambda-\lambda_0} \cdot \norm{(a-\lambda_0)^{-1}} <1   
	\]
	Also ist wieder nach \autoref{satz15_neumann} $\enbrace*{\ind- (\lambda-\lambda_0)(a-\lambda_0)^{-1}}^{-1} = \sum_{n=0}^{\infty} (\lambda-\lambda_0)^n (a-\lambda_0)^{-n}$ und 
	\begin{align}
		-\sum_{n=0}^{\infty} (\lambda-\lambda_0)^n(a-\lambda_0)^{-(n+1)} &= (a- \lambda_0)^{-1} \enbrace*{(\lambda-\lambda_0)(a-\lambda_0)^{-1}-\ind}^{-1} \\
		&= \enbrace[\Big]{(\lambda-\lambda_0) \Underbrace{(a-\lambda_0)^{-1}(a-\lambda_0)}{=\ind}-(a- \lambda_0)}^{-1} = (\lambda-a)^{-1}
	\end{align}
	Für ein stetiges Funktional $\varphi \in A^*$ erhalten wir 
	\[
		\varphi\enbrace*{(\lambda-a)^{-1}} = \sum_{n=0}^{\infty} - \varphi \enbrace*{(a-\lambda_0)^{-(n+1)}} \cdot (\lambda-\lambda_0)^n
	\]
	falls $\abs{\lambda-\lambda_0} < \norm{(a- \lambda_0)^{-1}}^{-1}$. Daher ist $\lambda \mapsto \varphi \enbrace*{(\lambda-a)^{-1}}$ \Index{holomorph} auf 
	$\mathbb{C}\setminus \sigma(a)$. Falls $\sigma(a)= \emptyset$, so ist $\lambda \mapsto \varphi \enbrace*{(\lambda-a)^{-1}}$ \bet{ganz}\index{ganze Funktion}, d.h. holomorph 
	auf ganz $\mathbb{C}$. Wegen \eqref{19:eq:1} und dem Satz von Liouville\footnote{siehe \url{https://de.wikipedia.org/wiki/Satz_von_Liouville_\%28Funktionentheorie\%29}} ist 
	dann $\varphi \enbrace*{(\lambda-a)^{-1}}=0$ für alle $\lambda\in \mathbb{C}$. Aber $\varphi \in A^*$ war beliebig, und wegen Hahn-Banach gilt $(\lambda-a)^{-1}=0$. Dies 
	ist ein Widerspruch. 
\end{beweis}

\begin{satz}[{name=[Formel für den Spektralradius]},label=satz1.10]
	Sei $A$ eine unitale Banachalgebra. Für $a \in A$ gilt dann 
	\[
		r(a)=\lim_{n \to \infty}\norm{a^n}^{\frac{1}{n}}
	\]
\end{satz}
\begin{beweis}
	Sei $\varphi \in A^*$ ein stetiges Funktional. Wir haben im Beweis von \autoref{spec_nichtleer} gesehen, dass die Zuordnung 
	$\lambda \mapsto \varphi \enbrace*{(\lambda-a)^{-1}}$ analytisch ist für $\lambda \notin \sigma(a)$. Insbesondere gilt dies für $\abs{\lambda}> r(a)$. Dann ist auch
	$\lambda \mapsto \varphi \enbrace*{(\lambda^{-1}-a)^{-1}}$ analytisch für $0< \abs{\lambda}< \sfrac{1}{r(a)}$, ebenso 
	$\lambda \mapsto \varphi \enbrace*{(\ind- \lambda a)^{-1}}$. Die Singularität bei $0$ ist hebbar, also definiert $f(\lambda) := \varphi \enbrace*{(\ind-\lambda a)^{-1}}$
	eine analytische Funktion für $\abs{\lambda} < \sfrac{1}{r(a)}$. Dann existiert eine Folge $(\alpha_n)_{n \in N} \subset \mathbb{C}$ mit 
	$f(\lambda) = \sum_{n=0}^{\infty} \alpha_n \cdot \lambda^n$ für $\abs{\lambda}<\sfrac{1}{r(a)}$. Die Konvergenz ist absolut und gleichmäßig, also gilt insbesondere 
	$\abs{\alpha_n \cdot \lambda^n} \xrightarrow{\raisemath{-1pt}{n \to \infty}}0$. Falls andererseits $\abs{\lambda} < \sfrac{1}{\norm{a}} \le \sfrac{1}{r(a)}$ 
	ist, so gilt
	\[
		f(\lambda) = \varphi \enbrace*{(\ind- \lambda \cdot a)^{-1}} \StackText{\ref{satz15_neumann}}{=} \varphi \enbrace*{\sum_{n=0}^{\infty} (\lambda \cdot a)^n} = 
		\sum_{n=0}^{\infty} \varphi \enbrace[\big]{(\lambda \cdot a)^n} = \sum_{n=0}^{\infty} \varphi(a^n) \cdot \lambda^n
	\]
	also $\alpha_n =\varphi(a^n)$ für $n \in N$. Es gilt $\varphi(a^n) \cdot \lambda^n \xrightarrow{\raisemath{-1pt}{n \to\infty}} 0$. Aber $\varphi \in A^*$ war beliebig, 
	also ist für jedes $\varphi \in A^*$ die Folge $\enbrace*{\varphi(\lambda^n a^n)}_\mathbb{N} = \enbrace*{\varphi(a^n) \cdot \lambda^n}_{n \in \mathbb{N}}$ beschränkt für 
	$\abs{\lambda}<\sfrac{1}{r(a)}$.

	Nach dem Prinzip der gleichmäßigen Beschränktheit ist dann $(\lambda^n a^n)_{n \in \mathbb{N}}$ beschränkt für jedes $\lambda$ mit $\abs{\lambda}< \sfrac{1}{r(a)}$. Damit 
	existiert für $0 < \abs{\lambda}< \sfrac{1}{r(a)}$ ein $M \ge 0$ mit $\norm{\lambda^n a^n} \le M$, $n \in \mathbb{N}$, also 
	\[
		\norm{a^n}^{\sfrac{1}{n}}\le\frac{M^{\sfrac{1}{n}}}{\abs{\lambda}} \xrightarrow{n \to \infty} \frac{1}{\abs{\lambda}} 
	\]
	Es folgt $\limsup_n \norm{a^n}^{\sfrac{1}{n}} \le r(a)$. Falls $\lambda \in \sigma(a)$, so gilt nach \autoref{satz14_poly} $\lambda^n \in \sigma(a^n)$ und wegen 
	\autoref{spec_kompakt} $\abs{\lambda^n} \le \norm{a^n}$, also $\abs{\lambda} \le \norm{a^n}^{\sfrac{1}{n}}$ und es folgt insgesamt
	\begin{align}
		r(a) \le \inf_n \norm{a^n}^{\sfrac{1}{n}} \le \liminf_n \norm{a^n}^{\sfrac{1}{n}} \le \limsup_n \norm{a^n}^{\sfrac{1}{n}} \le r(a)
	\end{align}
	und damit $\lim_n \norm{a^n}^{\sfrac{1}{n}} = r(a)$, woraus die Behauptung folgt. 
\end{beweis}

\begin{proposition}[label=prop1.11_einfach,{name=[Einfache abelsche unitale Banachalgebra]}]
	Sei $A$ eine abelsche unitale Banachalgebra. Falls $A$ einfach\marginnote{einfach = besitzt keine nichttrivialen Ideale} ist, so ist $A \cong \mathbb{C}$.\index{einfach} 
\end{proposition}
\begin{beweis}
	Sei $a \in A$, $\alpha \in \sigma(a)$. Dann ist $J := \overline{(a-\alpha \cdot \ind)\cdot A} \lhd A$ ein abgeschlossenes Ideal. Für jedes $b \in A$ ist $(a-\alpha)b$ nicht 
	invertierbar. Damit folgt dann $\norm{(a-\alpha)b - \ind}\ge 1$ für jedes $b \in A$ nach \autoref{satz15_neumann}. Damit ist auch 
	$\dist\enbrace*{(a-\alpha \cdot \ind)\cdot A,\ind}\ge 1$, also $\ind \notin J$ und damit $J \neq A$. Da $A$ einfach ist, folgt $J=\set{0}$, also muss 
	$a-\alpha \cdot \ind=0$ und weiter $a \in \mathbb{C} \cdot \ind$ gelten. Damit ist $A = \mathbb{C} \cdot \ind \cong \mathbb{C}$.
\end{beweis}

\begin{definition}[{name=[Charakter]},label=def:charakter]
	Sei $A$ eine abelsche unitale Banachalgebra. Eine lineare, multiplikative, nicht verschwindende Abbildung $\chi \colon A \to \mathbb{C}$ heißt \Index{Charakter} auf $A$.  
	Es gilt dann $\chi(\ind)=1$. Wir setzen $\Omega(A) := \set[\big]{\chi \given \chi \text{ Charakter auf }A}$ (\Index{Charakterspektrum} oder einfach nur \Index{Spektrum} von 
	$A$). 
\end{definition}

\begin{satz}[label=satz:1.13,{name=[Stetigkeit von Charakteren, Zusammenhang mit maximalen Idealen]}]
	Sei $A$ eine abelsche unitale Banachalgebra. Dann gilt
	\begin{enumerate}[(i),itemsep=0pt]
		\item Jeder Charakter auf $A$ ist stetig mit Norm $1$.
		\item $\chi \mapsto \ker \chi$ definiert eine Bijektion  
		\(
			\Omega(A) \to J_{\max} := \set[\big]{J \given J \subsetneq A \text{ ist maximales Ideal} }
		\)
	\end{enumerate}
\end{satz}
\begin{beweis} \leavevmode
	\begin{enumerate}[(i)]
		\item Sei $a \in A$ mit $\norm{a}<1$ und $\chi(a)=1$. 
		Setze $b= \sum_{n=1}^{\infty} a^n$, dann ist $\ind + b = \sum_{n=0}^{\infty} a^n = (\ind-a)^{-1}$. 
		Also ist $\ind = (\ind+b)(\ind-a) = \ind + b -a -ab$ und damit $ab+a=b$. 
		Es folgt $\chi(a)\chi(b) + \chi(a)= \chi(b)$. Dies ist ein Widerspruch. Also muss $\norm{\chi}\le 1$ sein. 
		Mit $\chi(\ind)=1$ folgt $\norm{\chi}=1$. 
		\item Für jedes $\chi \in \Omega(A)$ ist $\ker \chi \lhd A$ ein Ideal. 
		Weiter ist der Kern abgeschlossen in $A$, da $\chi$ stetig ist. 
		Wir zeigen, dass $\ker \chi$ maximal ist. 
		Es ist $\ker \chi \neq A$, denn es gilt $\chi \neq 0$. 
		Falls $c \in A \setminus \ker \chi$, so gilt 
		\(
			\chi(c)^{-1} \cdot c-\ind \in \ker \chi
		\).
		Also ist $\ker \chi + \mathbb{C} \cdot c = \ker \chi + \mathbb{C} \cdot \ind=A$, da für $a \in A$ gilt $a- \chi(a) \cdot \ind \in \ker \chi$ und damit 
		$a= \enbrace*{a- \chi(a) \cdot \ind}+ \chi(a) \cdot \ind \in \ker \chi + \mathbb{C} \cdot \ind$. 
		Damit ist der Kern von $\chi$ maximal.
		
		Zur Injektivität: Es seien $\chi_1, \chi_2 \in \Omega(A)$ Charaktere mit $\ker \chi_1 = \ker \chi_2$. Dann gilt für $a \in A$
		\begin{align}
			\chi_1(a) - \chi_2(a) = \chi_1 \enbrace[\Big]{\Underbrace{a - \chi_2(a) \cdot \ind}{\in \ker \chi_2=\ker \chi_1}} = 0
		\end{align}
		Also ist $\chi_1=\chi_2$. Für die Surjektivität betrachten wir ein maximales, nichttriviales Ideal $J \subsetneq A$. 
		Dann ist $\dist \enbrace*{J,\ind}=1$, denn für den offenen Ball mit Radius $1$ um $\ind$ gilt $\mathring{B}_A(\ind,1) \subset \Inv(A)$ nach \autoref{satz15_neumann}. 
		Dann ist $\ind \notin \overline{J} \lhd A$. 
		Da $J$ maximal ist, folgt $J=\overline{J}$ und damit ist 
		$\sfrac{A}{J}$ eine einfache, abelsche, unitale Banachalgebra und mit \autoref{prop1.11_einfach} folgt $\sfrac{A}{J} \cong \mathbb{C}$. Also ist 
		$\chi \colon A \to \sfrac{A}{J} \cong \mathbb{C}$ ein Charakter mit $\ker \chi= J$. \qedhere
	\end{enumerate}
\end{beweis}

\begin{satz}[label=satz:1.14,{name=[Spektrum entspricht Menge der $\chi(a)$]}]
	Sei $A$ eine abelsche unitale Banachalgebra. Für $a \in A$ gilt $\sigma(a)=\set[\big]{\chi(a) \given \chi \in \Omega(A)}$.
\end{satz}
\begin{beweis}
	Für $\lambda \in \sigma(a)$ ist $J := (a -\lambda)A \lhd A$ ein echtes Ideal, da $(a-\lambda)$ nicht invertierbar ist. 
	Nach dem Lemma von Zorn ist $J$ in einem maximalen echten Ideal enthalten, also ist $J \subset \ker \chi$ für ein $\chi \in \Omega(A)$ nach \autoref{satz:1.13}. 
	Dann ist $\chi(a- \lambda)=0$ für einen Charakter $\chi$ und somit auch $\chi(a)=\lambda$. 
	Es gilt also $\sigma(a) \subseteq \set{\chi(a) \given \chi \in \Omega(A)}$.	
	Für $\chi \in \Omega(A)$ gilt $a- \chi(a) \cdot \ind \in \ker \chi$. Damit kann $a- \chi(a) \cdot \ind$ nicht invertierbar sein. Also ist $\chi(a) \in \sigma(a)$ und es folgt 
	die zweite Inklusion.
\end{beweis}

\begin{satz}[label=satz:1.15,{name=[Charakterspektrum ist kompakt]}]
	Sei $A$ eine abelsche unitale Banachalgebra. Dann ist $\Omega(A)$ ein kompakter Hausdorffraum bezüglich der $\w^*$-Topologie.
\end{satz}
\begin{beweis}
	Es ist $\Omega(A) \subset B^1(A^*)$, wobei die Einheitskugel kompakt und Hausdorffsch ist bezüglich der $\w^*$-Topologie nach Banach-Alaoglu und da der Raum lokalkonvex 
	ist. Sei $\chi_\gamma \raisemath{-2pt}{\grenzw{\raisebox{-4pt}{\scriptsize$\w^*$}}}\chi$ ein Netz. $\chi_\gamma$ ist multiplikativ für jedes $\gamma$ und $\chi_\gamma(\ind)=1$. Damit gelten diese Eigenschaften auch für 
	$\chi$.
\end{beweis}

\begin{definition}[{name=[Charakterspektrum und Gelfandtransformation]}]
	Sei $A$ eine abelsche unitale Banachalgebra. $\Omega(A)$ mit der $\w^*$-Topologie heißt \bet{Spektrum}\index{Spektrum!Charakterspektrum} von $A$. Für $a \in A$ definieren wir
	die \Index{Gelfandtransformation} $\widehat{a} \colon \Omega(A) \to \mathbb{C}$ durch $\widehat{a}(\chi)=\chi(a)$. 
\end{definition}

\begin{satz}[label=satz:117,{name=[Gelfandtransformation ist normvermindernd]}]
	Sei $A$ eine abelsche unitale Banachalgebra. Dann ist $\widehat{\hspace{1ex}} \colon A \to C \enbrace*{\Omega(A)}$, $a \mapsto \widehat{a}$ ein normvermindernder Homomorphismus und 
	$r(a) = \norm{\widehat{a}}_{\infty}$ für $a \in A$. Weiter gilt $\sigma(a)= \widehat{a}\enbrace*{\Omega(A)}$.
\end{satz}
\begin{beweis}
	$\widehat{a} \colon \Omega(A) \to \mathbb{C}$ ist stetig nach Definition der $\w^*$-Topologie, denn die $\w^*$-Topologie ist per Definition die gröbste Topologie auf $A^*$, sodass die Einsetzungshomomorphismen stetig sind.
	% Für $U \subset \mathbb{C}$ offen ist
% 	\[
% 		\widehat{a}^{-1}(U) = \set[\big]{\chi \in \Omega(A) \given  \widehat{a}(\chi) \in U} = \set[\big]{\chi \in \Omega(A) \given \chi(a) \in U} =
% 		\set[\big]{\chi \in A^* \given \chi(a) \in U} \cap \Omega(A)
% 	\]
% 	Die linke Menge des Schnitts ist gleich
% 	\[
% 		\smashoperator[r]{\bigcup_{(\mu,\varepsilon), B(\mu,\varepsilon) \subset U}} \,\,\set[\big]{\varphi \in A^* \given \abs{\varphi(a)-\mu}<\varepsilon }
% 	\]
% 	und somit offen in $A^*$.
	% Damit ist $\widehat{a}^{-1}(U)$ offen in der Relativtopologie.
	Wir haben weiter
	\[
		\sigma(a) \StackText{\ref{satz:1.14}}{=} \set[\big]{\chi(a) \given \chi \in \Omega(A)}  = \set[\big]{\widehat{a}(\chi) \given \chi \in \Omega(A)} = \widehat{a} 
		\enbrace[\big]{\Omega(A)}
	\]
	Also gilt außerdem
	\[
		\norm{a} \StackText{\ref{spec_kompakt}}{\le} r(a) = \sup_{\lambda \in \sigma(a)} \abs*{\lambda} \StackText{\ref{satz:1.14}}{=} \sup_{\chi \in \Omega(A)} \abs{\widehat{a}(\chi)} 
		= \norm{\widehat{a}}_{\infty} \qedhere
	\]
\end{beweis}
% section 1 (end)
\newpage

\section{$C^*$-Algebren und Funktionalkalkül} % (fold)
\label{sec:2}

\begin{definition}[{name=[Involutive Banachalgebra]}]
	Eine Banachalgebra heißt \Index{involutiv}, falls eine antilineare Abbildung ${}^* \colon A \to A$ existiert, für die folgenden Identitäten gelten:
	\[
		a^{**}=a \qquad \quad (ab)^*= b^* a^* \qquad \quad \norm{a^*} = \norm{a}  
	\] 
\end{definition}

\begin{definition}[{name=[$C^*$-Algebra]},label=def:22]
	Eine \bet{$C^*$-Algebra}\index{C-Algebra@$C^*$-Algebra} ist eine involutive Banachalgebra mit $\norm{a^*a}= \norm{a}^2$ für $a \in A$.
\end{definition}

\begin{definition}[{name=[{selbstadjungiert, normal, unitär, Projektion, Isometrie}]}]
	Sei $A$ ein $C^*$-Algebra. $a \in A$ heißt 
	\[
		\begin{cases}
			\text{ \Index{selbstadjungiert} }  &\text{ falls } a=a^*\\
			\text{ \Index{normal} }  &\text{ falls }  aa^*=a^*a \\
			\text{ \Index{unitär} }  &\text{ falls }  aa^*=a^*a=\ind \\
			\text{ \Index{Projektion} }  &\text{ falls }   a=a^*=a^2 \\
			\text{ \Index{Isometrie} }  &\text{ falls } a^*a= \ind  
		\end{cases}
	\]
	Wo die Existenz eines Einselements nötig ist, fordern wir zusätzlich, dass $A$ unital ist.
\end{definition}

\begin{definition}[{name=[${ }^*$-Homomorphismus]}]
	Es seien $A,B$ zwei $C^*$-Algebren. Eine Abbildung $\pi \colon A \to B$ heißt \bet{${}^*$"=Homomorphismus}\index{Stern-Homomorphismus@Sten-, ${}^*$-Homomorphismus}, falls $\pi$ 
	linear, multiplikativ und ${}^*$-erhaltend ist. $\pi$ heißt
	\bet{${}^*$-Einbettung}\index{Stern-Einbettung@Stern-, ${}^*$-Einbettung}, falls $\pi$ zusätzlich isometrisch ist. 
\end{definition}

\begin{bemerkung}[{name=[Direkte Summe von $A$ mit den komplexen Zahlen]}]
	Falls $A$ eine (involutive) Banachalgebra ist, so ist $A \oplus \mathbb{C}$ mit offensichtlichen Vektorraumoperationen, der naheliegenden Multiplikation
	$(a,\lambda) \cdot (b,\mu) = \enbrace*{ab + \mu a + \lambda b, \lambda \mu}$ sowie $(a,\lambda)^* = (a^*,\overline{\lambda})$ im involutiven Fall und Norm definiert wir durch
	$\norm{(a,\lambda)}:= \norm{a}_{A} + \abs{\lambda}$ eine unitale (involutive) Banachalgebra mit $\ind=(0,1)$. $A \oplus \mathbb{C}$ ist jedoch keine $C^*$-Algebra, auch wenn
	$A$ eine $C^*$-Algebra ist.
\end{bemerkung}

\begin{satz}[label=satz:26,{name=[Unitalisierung]}]
	Sei $A$ eine $C^*$-Algebra. Dann existiert eine unitale $C^*$-Algebra $A^\simm$ mit einer isometrischen ${}^*$-Einbettung $\iota \colon A \hookrightarrow A^\simm$, sodass
	$\iota(A) \lhd A^\simm$ ein Ideal ist mit 
	\[
		\Faktor{A^\simm}{\iota(A)} \cong \begin{cases}
			0, &\text{ falls }A \text{ unital} \\
			\mathbb{C} , &\text{ sonst }  
		\end{cases}
	\]
\end{satz}
\begin{beweis}
	Für $a \in A$ definieren wir $\iota(a) \in \mathcal{B}(A)$, wobei $\mathcal{B}(A)$ die beschränkten Operatoren auf $A$ aufgefasst als Banachraum sind, durch
	$\iota(a)(b)=ab$. Definiere $A^\simm := \iota(A) + \mathbb{C} \cdot \ind_{\mathcal{B}(A)} \subset \mathcal{B}(A)$. Dann ist $A^\simm$ eine unitale Algebra mit Involution
	\[
		\enbrace*{\iota(a)+\lambda \cdot \ind_{\mathcal{B}(A)}}^* := \iota(a^*) + \overline{\lambda} \cdot \ind_{\mathcal{B}(A)}
	\]
	Dass $\iota(A) \lhd A^\simm$ ein Ideal ist, ist klar. $A^\simm$ ist normiert mittels der Operatornorm. Für $0 \neq a \in A$ gilt  
	\begin{align}
		\norm{a}_A = \frac{1}{\norm{a}_A}\norm{aa^*}_A = \norm*{a \enbrace*{\frac{1}{\norm{a^*} } \cdot a^*}}_A &\le \norm*{\iota(a) \enbrace*{\frac{1}{\norm{a^*}}\cdot a^* }}_A \\
		&\le \sup_{\norm{b}_A \le 1} \norm[\big]{\iota(a)(b)}_A = \norm[\big]{\iota(a)}_{\mathcal{B}(A)} \\
		&= \sup_{\norm{b}_A \le 1} \norm{ab}_A \\
		&\le \norm{a}_A   
	\end{align}
	Also ist $\iota$ isometrisch. Es folgt, dass $\iota(A) \subset \mathcal{B}(A)$ abgeschlossen ist, da $A$ vollständig ist. Weiter folgt, dass auch 
	$A^\simm = \iota(A) + \mathbb{C} \cdot \ind_{\mathcal{B}(A)}$ vollständig ist bezüglich $\norm{\cdot}_{\mathcal{B}(A)}$. Wir zeigen nun, dass $\norm{\cdot}_{\mathcal{B}(A)}$ 
	eine $C^*$-Norm ist auf $A^\simm$: 
	\begin{align}
		\norm*{\iota(a)+ \lambda \cdot \ind_{\mathcal{B}(A)}}_{\mathcal{B}(A)}^2 &= \sup_{b \in A, \norm{b} \le 1 } \norm*{ab + \lambda b}^2_A \\
		&= \sup_{b \in A, \norm{b} \le 1} \norm[\big]{\enbrace*{ab+\lambda b}^* (ab +\lambda b)}_A \\
		&= \sup_{b \in A, \norm{b} \le 1} \norm[\big]{b^* a^* ab + \lambda b^* a^* b + \overline{\lambda} b^* ab + \abs*{\lambda}^2 b^* b}_{A} \\
		&\le \sup_{b \in A, \norm{b} \le 1} \norm[\big]{a^*ab + \lambda a^* b + \overline{\lambda} a b + \abs{\lambda}^2 b }_A \tag{$*$} \label{eq:26:1}\\
		&= \norm[\big]{\iota \enbrace*{a^* a + \lambda a^* + \overline{\lambda} a} + \abs{\lambda}^2 \cdot \ind_{\mathcal{B}(A)} }_{\mathcal{B}(A)} \\
		&= \norm*{ \enbrace*{\iota(a) + \lambda \cdot \ind_{\mathcal{B}(A)}}^* \enbrace*{\iota(a) + \lambda \cdot \ind_{\mathcal{B}(A)}}}_{\mathcal{B}(A)} \\
		&\le \norm*{\enbrace*{\iota(a) + \lambda \cdot \ind_{\mathcal{B}(A)}}^*}_{\mathcal{B}(A)} \norm*{\iota(a) + \lambda \cdot \ind_{\mathcal{B}(A)}}_{\mathcal{B}(A)}   
		\tag{$**$} \label{eq:26:2}
	\end{align}
	Es folgt, dass 
	\begin{align}
		\norm*{\iota(a) + \lambda \cdot \ind_{\mathcal{B}(A)}}_{\mathcal{B}(A)} \le \norm*{\enbrace*{\iota(a) + \lambda \cdot \ind_{\mathcal{B}(A)}}^*}_{\mathcal{B}(A)}
		\le \norm*{\iota(a) + \lambda \cdot \ind_{\mathcal{B}(A)}}_{\mathcal{B}(A)} 
	\end{align}
	Damit folgt, dass \eqref{eq:26:1} und \eqref{eq:26:2} Gleichungen sind und $\norm*{\cdot }_{\mathcal{B}(A)}$ eine $C^*$-Norm auf $A^\simm$ ist. Falls $A$ unital ist, gilt
	$\iota(\ind_A) = \ind_{\mathcal{B}(A)}$ und somit $\iota(A)=A^\simm$. Ist $A$ nicht unital, so ist $A^\simm \cong \iota(A) \oplus \mathbb{C}$ als Vektorräume und es muss
	$\sfrac{A^\simm}{\iota(A)}\cong \mathbb{C}$ gelten.
\end{beweis}

\begin{bemerkung}[{name=[Alternative Unitalisierung]},label=bem:27]
	Man definiert oft auch eine unitale $C^*$-Algebra 
	\[
		A^+ := \begin{cases}
			A^\simm, &\text{ falls }A \text{ nicht unital} \\
			A  \oplus \mathbb{C} , &\text{ falls } A \text{ unital} 
		\end{cases}
	\]
	mit $\norm*{(a,\lambda)} := \max \set*{\norm*{a}_A, \abs{\lambda}}$. Dann ist in jedem Fall $A^+ \cong A \oplus \mathbb{C}$ als Vektorraum. $A^+$ lässt sich charakterisieren 
	durch die universelle Eigenschaft:\marginnote{Blatt3 Aufgabe 2}
	\[
		\begin{tikzcd}
			\mbox{ } & B \\
			A \rar[hook] \urar& A^+ \uar[dashed,"\exists !"']
		\end{tikzcd}
	\]
	wobei $B$ eine unitale $C^*$-Algebra ist und $A \to B$ ein ${}^*$-Homomorphismus ist. Eine ähnliche Beschreibung für $A^\simm$ ist möglich, aber etwas \enquote{subtiler}.
\end{bemerkung}

\begin{definition}[{name=[Spektrum in nichtunitalen $C^*$-Algebren]}]
	Sei $A$ eine nichtunitale $C^*$-Algebra. Für $a \in A$ setzen wir 
	\[
		\sigma(a) := \sigma_{A^\simm} \enbrace[\big]{\iota(a)} \quad \text{ und } \quad r(a) := r_{A^\simm} \enbrace[\big]{\iota(a)} 
	\]
	Wieder schreiben wir $\Omega(A) := \set[\big]{\chi \colon A \to \mathbb{C} \text{ Charaktere}}$ und \marginnote{Für $A$ unital sind maximale Ideale automatisch abgeschlossen}
	\[
		J_{\max}(A)= \set[\big]{J \given J \lhd A \text{ echtes maximales abgeschlossenes Ideal}}
	\]
\end{definition}

\begin{satz}[label=satz:29,{name=[Spektrum einer nichtunitalen $C^*$-Algebra ist lokalkompakt]}]
	Sei $A$ eine nichtunitale abelsche $C^*$-Algebra. Dann ist $\Omega(A)$ ein lokalkompakter Hausdorffraum bezüglich der $\w^*$-Topologie. Für $a \in A$ gilt 
	$\sigma(a) = \set[\big]{\chi(a) \given \chi \in \Omega(A)} \cup \set*{0}$.
\end{satz}
\begin{beweis}
	Für $\chi \in \Omega(A)$ sei $\chi^\simm \colon A^\simm \to \mathbb{C}$ die eindeutig bestimmte, lineare, unitale Fortsetzung. Es gilt $\chi^\simm \in \Omega(A^\simm)$ und
	\[
		\Omega(A^\simm) = \set[\big]{\chi^\simm \given \chi \in \Omega(A)} \cup \set[\big]{\chi_\infty}
	\]
	wo $\chi_\infty \colon A^\simm \to \sfrac{A^\simm}{\iota(A)} \cong \mathbb{C}$. $\Omega(A^\simm)$ ist ein kompakter Hausdorffraum nach \autoref{satz:1.15}, also muss der linke 
	Teil der Vereinigung ein lokalkompakter Hausdorffraum sein.  Man überprüft\marginnote{Blatt 2 Aufgabe 2}, dass $\chi \mapsto \chi^\simm$ ein Homöomorphismus 
	$\Omega(A)\to \Omega(A^\simm) \setminus \set*{\chi_\infty}$ bezüglich der $\w^*$-Topologie ist. Es gilt
	\[
		\sigma(a) := \sigma_{A^\simm} \enbrace[\big]{\iota(a)} \StackText{\ref{satz:117}}{=} \set[\big]{\chi \enbrace*{\iota(a)} \given \chi \in \Omega(A^\simm)} = 
		\Underbrace{\set[\big]{\chi^\simm \enbrace*{\iota(a)}\given\chi \in \Omega(A)}}{= \set[\big]{\chi(a) \given \chi \in \Omega(A)}} \cup \Underbrace{\set[\big]{\chi_\infty 
		\enbrace*{\iota(a)}}}{=\set*{0}} \qedhere
	\]
\end{beweis}

\begin{proposition}[label=prop:210,{name=[Eigenschaften von Projektionen, unitären und selbstadjungierten Elementen]}]
	Sei $A$ eine $C^*$-Algebra. Dann gilt
	\begin{enumerate}[(i),itemsep=0pt]
		\item Für eine Projektion $0 \neq p  \in A$ gilt $\norm*{p} =1$.
		\item Ist $\ind \in A$ und $u \in A$ unitär, dann ist $\norm*{u}=1$ und $\sigma(u) \subset \mathbb{T}$.
		\item \label{210:3} Für $a \in A_{\sa}$ ist $r(a)= \norm*{a}$ und $\sigma(a) \subset \mathbb{R}$.
		\item Für $a \in A$ beliebig gilt $\norm*{a} = \norm*{a^* a}^{\sfrac{1}{2}}= r(a^* a)^{\sfrac{1}{2}}$.
	\end{enumerate}
\end{proposition}
\begin{beweis}
	\leavevmode
	\begin{enumerate}[(i),itemsep=0pt]
		\item Es gilt $\norm*{p} = \norm*{p^*p}= \norm*{p}^2$ und somit muss $\norm*{p}=1$ gelten.
		\item Es gilt $\norm*{u}^2= \norm*{u^* u}= \norm*{\ind}=1$, also $\norm*{u}=1=\norm*{u^*}$. Wenn $\lambda \in \sigma(u)$, dann gilt $\lambda \neq 0$ und es gilt
		$\lambda^{-1} \in \sigma(u^*)$. Da die Spektralradien von $u$ und $u^*$ kleiner gleich $1$ sind, muss $\abs{\lambda},\abs{\lambda^{-1}}\le 1$ gelten, also $\abs{\lambda}=1$
		und damit $\lambda \in \mathbb{T}$.
		\item Es gilt $\norm*{a^2}= \norm*{a^* a}= \norm*{a}^2$. Eine Induktion liefert $\norm*{a^{2^n}}= \norm*{a}^{2^n}$, also nach \autoref{satz1.10} 
		\[
			r(a) = \lim_{n \to \infty} \norm*{a^n}^{\sfrac{1}{n}} = \lim_{n \to \infty} \norm*{a^{2^n}}^{\sfrac{1}{2^n}} = \norm*{a}
		\]
		Weiter existiert\footnote{$(ia)^0 = \ind_{A^\simm}$ falls $A$ nicht unital} $e^{ia} := \sum_{n=0}^{\infty} \frac{1}{n!}(ia)^n \in A^\simm$ und ist unitär mit 
		\(
			\enbrace{e^{ia}}^* = e^{(ia)^*} = e^{-ia} = \enbrace{e^{ia}}^{-1}
		\), 
		also ist $\sigma(e^{ia}) \subset \mathbb{T}$ nach (ii). Falls $\lambda \in \sigma(a)$, dann setze $b := \sum_{n=1}^{\infty} \frac{1}{n!} i^n (a-\lambda)^{n-1}$. Dann gilt
		\[
			(a-\lambda) b e^{i \lambda} = \enbrace[\big]{e^{i(a-\lambda)}- \ind}e^{i \lambda} = e^{i a} - e^{i \lambda}
		\]
		Also ist $e^{i \lambda} \in \sigma(e^{ia}) \subset \mathbb{T}$ und somit $\lambda \in \mathbb{R}$.
		\item Folgt direkt aus \ref{210:3}. \qedhere
	\end{enumerate}
\end{beweis}

\begin{korollar}[label=kor:211,{name=[*-Homomorphismen sind kontraktiv]}]
	Es seien $A$ und $B$ zwei $C^*$-Algebren und $\pi \colon A \to B$ ein ${}^*$-Homomorphismus. Dann ist $\pi$ kontraktiv.
	Die Aussage bleibt richtig, falls $A$ nur eine involutive Banachalgebra ist. 
\end{korollar}
\begin{beweis}
	Wir dürfen $\pi$ als unital annehmen (betrachte sonst $\pi^{\hspace{-0.5pt}+}\colon A^{\hspace{-0.9pt}+} \to B^{\hspace{-0.5pt}\sim}$). 
	Für $d \in A$ gilt $\sigma_B \enbrace[\big]{\pi(d)} \subset \sigma_A(d)$, also $r_B \enbrace[\big]{\pi(d)} \le r_A(d)$. 
	Für $a \in A$ gilt nach \autoref{prop:210} \ref{210:3}
	\[
		\norm*{\pi(a)}^2 = \norm*{\pi(a)^* \pi(a)} = \norm*{\pi(a^*a)} = r_B \enbrace[\big]{\pi(a^*a)} \le r_A(a^*a) = \norm*{a}^2 \qedhere
	\]
\end{beweis}

\begin{korollar}[{name=[höchstens eine $C^*$-Norm]},label=korr:212]
	Auf einer ${}^*$-Algebra gibt es höchstens eine $C^*$-Norm.
\end{korollar}
\begin{beweis}
	Sei $A$ eine ${}^*$-Algebra mit $C^*$-Normen $\norm*{\cdot}_1$ und $\norm*{\cdot}_2$. 
	Wir bezeichnen die $C^*$-Algebren entsprechend mit $A_1$ und $A_2$. 
	Es seien $\pi_1 \colon A_1 \to A_2$ und $\pi_2 \colon A_2 \to A_1$ die Identitäten. Dann gilt
	\[
		\norm*{a}_1 = \norm*{\pi_2(a)}_1 \le \norm*{a}_2 = \norm*{\pi_1(a)}_2 \le \norm*{a}_1 \qedhere
	\]
\end{beweis}

\begin{korollar}[label=kor:213,{name=[Charaktere sind *-Homomorphismen]}]
	Sei $A$ eine $C^*$-Algebra und $\chi$ ein Charakter auf $A$. Dann gilt $\chi(a^*) = \overline{\chi(a)}$ für $a \in A$, das heißt $\chi$ ist ein ${}^*$-Homomorphismus.
\end{korollar}
\begin{beweis}
	Sei $a \in A$, dann ist $a=\frac{1}{2}(a+a^*) + i \cdot \frac{1}{2i}(a-a^*) =: b + i \cdot c$ und $b,c$ sind selbstadjungiert. Es gilt 
	\[
		\chi(b) \StackText{\ref{satz:29}}{\in} \sigma(b) \StackText{\ref{prop:210}\ref{210:3}}{\subset} \mathbb{R} \quad \text{ und } \quad \chi(c) \in \sigma(c) \subset \mathbb{R} 
	\]
	Damit folgt $\chi(a^*) = \chi(b - i \cdot c) = \chi(b) - i \cdot \chi(c) = \overline{\chi(b) + i \cdot \chi(c)} = \overline{\chi(a)}$.
\end{beweis}

Aus dem bisher gezeigten können wir nun eine Version von \autoref{satz:1.13} für nicht notwendigerweise unitale $C^*$-Algebren ableiten: 

\begin{korollar}[{name=[Stetigkeit von Charakteren und Zusammenhang mit maximalen Idealen]}]
	Sei $A$ eine abelsche $C^*$-Algebra. Dann gilt
	\begin{enumerate}[(i),itemsep=0pt]
		\item Jeder Charakter $\chi \colon A \to \mathbb{C}$ ist stetig mit Norm $\le 1$.
		\item Die Abbildung $\chi \mapsto \ker \chi$ definiert eine Bijektion 
		\[
			\Omega(A) \to J_{\max} = \set[\big]{J \given J \subsetneq A \text{ maximales, abgeschlossenes Ideal}}
		\]
	\end{enumerate}
\end{korollar}
\begin{beweis}
	\leavevmode
	\begin{enumerate}[(i),itemsep=0pt]
		\item folgt aus \autoref{kor:211} und \autoref{kor:213} oder auch aus \autoref{satz:1.13} und \autoref{satz:26}.
		\item Wir zeigen nur die Surjektivität: Sie $J \subsetneq A$ ein abgeschlossenes, maximales Ideal. Dann ist $J \subsetneq A^\simm$ ein abgeschlossenes Ideal und es gilt
		\begin{align}
			J \subset A \cap \bigcap_{\substack{I \lhd A^\simm \text{ max} \\ J \subset I} } \hspace{-1ex} I  
			\StackText{\ref{satz:1.13}}{=} A \cap \bigcap_{\substack{\chi \in \Omega(A^\simm) \\ J \subset \ker \chi}} \hspace{-1ex}\ker \chi 
			&\StackText{\ref{satz:29}}{=} A \cap \ker \chi_\infty \cap \bigcap_{\substack{\chi \in \Omega(A) \\ J \subset \ker \chi^\simm}} \ker \chi^\simm  
		\end{align}
		Damit ist $J$ enthalten im Schnitt über alle Kerne der $\chi \in \Omega(A)$, wobei $J \subset \ker \chi$. Damit ist $J \subset \ker \chi$ für ein $\chi \in \Omega(A)$ und
		da $\ker \chi \lhd A$ abgeschlossen und maximal ist, folgt $J=\ker \chi$. \qedhere
	\end{enumerate}
\end{beweis}

\begin{beispiel}[{name=[für C*-Algebren]}]
	Wir haben nun schon einiges über $C^*$-Algebren gelernt, es dabei aber total versäumt konkrete Beispiele anzugeben. Dies holen wir nun nach:
	\begin{enumerate}[(i)]
		\item Sei $\mathcal{H}$ ein Hilbertraum. Dann ist $\mathcal{B}(\mathcal{H})$ eine $C^*$-Algebra mit Operatornorm und Adjunktion: Für $T \in \mathcal{B}(\mathcal{H})$ gilt
		\(
			\norm*{T} = \sup \set[\big]{\abs*{\skal*{\eta}{T \xi}} \given \norm*{\xi}, \norm*{\eta} \le 1}
		\)
		also
		\begin{align}
			\norm*{T^* T} = \sup \set[\big]{\abs*{\skal*{\eta}{T^* T \xi}} \given \norm*{\xi}, \norm*{\eta} \le 1} &\ge \sup \set[\big]{ \abs*{\skal*{T \eta}{T \eta}} \given 
			\norm*{\eta} \le 1}\\
			&= \enbrace*{\sup \set[\big]{\norm*{T \eta} \given \norm*{\eta} \le 1}}^2 = \norm*{T}^2 \ge \norm*{T^*T} \qedhere
		\end{align}
		\item Insbesondere ist $M_n := M_n(\mathbb{C}) = \mathcal{B}(\mathbb{C}^n)$ eine $C^*$-Algebra für $n \in \mathbb{N}$.
		\item Jede abgeschlossene ${}^*$-Unteralgebra von $\mathcal{B}(\mathcal{H})$ ist eine $C^*$-Algebra, zum Beispiel $\mathcal{K}(\mathcal{H})$.
		\item Sei $X$ ein lokalkompakter Hausdorffraum. Dann ist $C_0(X)$ eine $C^*$-Algebra mit punktweisen Operationen, Involution $f^*= \overline{f}$ und der Supremumsnorm.
		Tatsächlich ist jede abelsche $C^*$-Algebra von dieser Form: 
	\end{enumerate}
\end{beispiel}

\begin{satz}[name={Gelfand},label=satz16:gelfand]
	Sei $A \neq 0$ eine abelsche $C^*$-Algebra.\marginnote{$C_0 \enbrace*{\Omega(A)} = C \enbrace*{\Omega(A)}$, falls $A$ unital nach \autoref{satz:1.15}} 
	Dann ist die Gelfandtransformation $\widehat{\hspace{5pt}} \colon A \to C_0 \enbrace*{\Omega(A)}$ ein isometrischer ${}^*$-Isomorphismus.
\end{satz}
\begin{beweis}
	Betrachte das folgende Diagramm
	\[
		\begin{tikzcd}
			\color{dark_gray} a \rar[mapsto,dark_gray] \dar[phantom,"\in",sloped,dark_gray] & \color{dark_gray} \enbrace[\big]{\chi \mapsto \chi(a)} \dar[phantom,"\in",sloped,dark_gray] \\[-2pt]
			A \dar[hook,"\iota"] \rar["\widehat{\hspace{1ex}}"] 
			& \set[\big]{f \in C \enbrace[\big]{\Omega(A) \cup\set*{\chi_\infty}}\given f(\chi_\infty)=0} \cong C_0\enbrace*{\Omega(A)} \dar[hook] \\
			A^\simm \rar["\widehat{\hspace{1ex}}"] & C \enbrace[\big]{\Omega(A^\simm)} = C \enbrace[\big]{ \set[\big]{\chi^\simm \given \chi \in \Omega(A)} \cup \set*{\chi_\infty}} 
			\StackText{\ref{satz:29}}{\cong} C \enbrace[\big]{\Omega(A) \cup \set*{\chi_\infty}} 
		\end{tikzcd}
	\]
	Die Gelfandtransformation ist ein ${}^*$-Homomorphismus nach \autoref{kor:213} und \autoref{satz:117}. Weiter gilt 
	\[
		\norm*{a}^2 = \norm*{a^*a} \StackText{\ref{prop:210}}{=} r(a^*a) \StackText{\ref{satz:117}}{=} \norm[\big]{\widehat{a^*a}}_\infty = \norm*{\widehat{a^*} \widehat{a}}_\infty
		= \norm[\big]{\overline{\widehat{a}} \,\widehat{a}}_\infty = \norm[\big]{\abs{\widehat{a}}^2 }_\infty = \norm*{\widehat{a}}^2_\infty
	\]
	Damit ist $a \mapsto \widehat{a}$ isometrisch. 
	Damit folgt, dass $C := \set*{\widehat{a} \given a \in A} \subset C_0 \enbrace*{\Omega(A)}$ abgeschlossen ist. 
	Außerdem ist $C$ eine ${}^*$-Unteralgebra, also insbesondere invariant unter komplexer Konjugation. 
	Weiter trennt $C$ die Punkte von $\Omega(A)$ und zu jedem $\chi \in \Omega(A)$ existiert ein $a \in A$ mit $\widehat{a}(\chi)= \chi(a) \neq 0$, das heißt $\ev_\chi(C) \neq \set*{0}$ für alle $\chi \in \Omega(A)$. 
	Mit Stone-Weierstraß folgt, dass $C \subset C_0 \enbrace*{\Omega(A)}$ dicht ist. 
	Da $C$ abgeschlossen ist, folgt also $C=C_0 \enbrace*{\Omega(A)}$.
\end{beweis}

\begin{bemerkung}[label=bem217,{name=[Spektrum von $f \in C(X)$ für $X$ kompakt]}]
	Sei $X$ ein kompakter Hausdorffraum und $f \in C(X)$, dann ist $\sigma_{C(X)}(f)=f(X)$.
\end{bemerkung}
\begin{beweis}
	Die Abbildung $\delta \colon X \to \Omega\enbrace*{C(X)}$ definiert durch $x \mapsto \enbrace*{f \mapsto f(x)}$ ist ein Homöomorphismus wegen $C(X) \cong C \enbrace[\big]{\Omega(C(X))}$. 
	Es gilt
	\begin{align}
		\sigma_{C(X)}(f) \StackText{\ref{satz:117}}{=} \widehat{f} \enbrace[\big]{\Omega(C(X))} = \set*{\widehat{f}(\chi) \given \chi \in \Omega \enbrace[\big]{C(X)}} 
		&= \set*{\chi(f) \given \chi \in \Omega \enbrace[\big]{C(X)}} \\
		&= \set[\big]{\delta(x)(f) \given x \in X} = f(X) \qedhere
	\end{align}
\end{beweis}

\begin{definition}[{name=[erzeuge C*-Algebra]}]
	Sei $S \subset A$ Teilmenge einer $C^*$-Algebra $A$.\index{C-Algebra@$C^*$-Algebra!erzeugte $C^*$-Algebra}
	Dann bezeichnet $C^*(S) \subset A$ die kleinste $C^*$-Unteralgebra von $A$, die $S$ enthält; also
	\[
		C^*(S) := \smashoperator[r]{\bigcap_{\substack{B \subset A \text{ $C^*$-Unteralg.}\\ \text{mit } S \subset B}}} \enspace B
	\]
\end{definition}

\begin{satz}[name={Funktionalkalkül},label=funkkal]
	Sei $A$ eine $C^*$-Algebra und $a \in A$ normal.\index{Funktionalkalkül}
	Dann existiert ein eindeutig bestimmter ${}^*$-Isomorphismus 
	\[
		C^*(a) \cong C_0 \enbrace[\big]{\sigma(a)\setminus \set*{0}} \qquad \text{ mit }\quad  a \longmapsto \id_{\sigma(a)\setminus \set*{0}}
	\]
	Falls $A$ unital ist, so existiert ein eindeutig bestimmter unitaler ${}^*$-Isomorphismus 
	\[
		C^*(a,\ind) \cong C \enbrace[\big]{\sigma(a)} \qquad \text{ mit }\quad  a \longmapsto \id_{\sigma(a)}
	\]
	Für $f \in C_0\enbrace*{\sigma(a) \setminus \set*{0}}$ bzw. $f \in C \enbrace*{\sigma(a)}$ bezeichnen wir mit $f(a) \in A$ das Urbild von $f$ unter diesem Isomorphismus.
	Es gilt $f(a) \in C^*(a,\ind_{A^\simm}) \subset A^\simm$, falls $f(0)\neq 0$. 
	Weiter gilt 
	\[
		\sigma \enbrace[\big]{f(a)} = f \enbrace[\big]{\sigma(a)}
	\]
	Für $g \in C \enbrace[\big]{f(\sigma(a))}$ gilt dann $g \enbrace[\big]{f(a)}= (g \circ f)(a)$.
\end{satz}
\begin{beweis}
	Sei zunächst $A$ unital. Da $a \in A$ normal ist, ist $C^*(a,\ind) = \overline{\Alg(a,a^*,\ind)} \subset A$ abelsch. Die Abbildung
	\mapdef{\alpha \colon\Omega \enbrace[\big]{C^*(a,\ind)}}{\widehat{a} \enbrace[\big]{\Omega \enbrace*{C^*(a,\ind)}}\StackText{\ref{satz:117}}{=} \sigma_{C^*(A,\ind)}(a)\subset \mathbb{C}}%
	{\chi}{\widehat{a}(\chi)=\chi(a)}{}
	ist ein Homöomorphismus: Die Stetigkeit ist klar, denn $\Omega \enbrace*{C^*(a,\ind)}$ trägt die $\w^*$-Topologie. 
	Die Surjektivität ist trivial. 
	Zur Injektivität: Falls $\chi_1(a)= \widehat{a}(\chi_1) = \widehat{a}(\chi_2) = \chi_2(a)$ ist, so ist $\chi_1|_{\Alg(a,a^*,\ind)} = \chi_2|_{\Alg(a,a^*,\ind)}$ und mit der Stetigkeit von $\chi_1$ und $\chi_2$ folgt, dass $\chi_1=\chi_2$ sein muss.
	Die Inverse ist stetig, da $\Omega \enbrace*{C^*(a,\ind)}$ kompakt ist.
	
	Wir erhalten nach \autoref{satz16:gelfand} $C^*(a,\ind) \cong C \enbrace[\big]{\Omega(C^*(a,\ind))}$ mittels $a \mapsto \widehat{a}$ und $C \enbrace[\big]{\Omega(C^*(a,\ind))} \cong C \enbrace*{\sigma_{C^*(a,\ind)}(a)}$ mittels $\widehat{a} \mapsto \widehat{a} \circ \alpha^{-1}$. 
	Für $\alpha(\chi)=t \in \sigma_{C^*(a,\ind)}(a)$ haben wir 
	\[
		\enbrace*{\widehat{a} \circ \alpha^{-1}}(t) = \enbrace*{\widehat{a} \circ \alpha^{-1}} \enbrace[\big]{\alpha(\chi)} = \widehat{a}(\chi) = \alpha(\chi)=t
	\]
	Also ist $\widehat{a} \circ \alpha^{-1}=\id_{\sigma_{C^*(a,\ind)}(a)}$. 
	Wir nenen die Komposition der beiden Isomorphismen $\beta$ und zeigen, dass $\beta$ tatsächlich die in der Behauptung geforderten Eigenschaften hat:
	Für $f \in C \enbrace*{\sigma_{C^*(a,\ind)}(a)}$ gilt 
	\[
		\sigma_{C^*(a,\ind)} \enbrace[\big]{f(a)} = \sigma_{C^*(a,\ind)} \enbrace*{\beta^{-1}(f)} = \sigma_{\beta \enbrace*{C^*(a,\ind)}}(f) = \sigma_{C \enbrace[\big]{\sigma_{C^*(a,\ind)}(a)}} (f) \StackText{\ref{bem217}}{=} f \enbrace*{\sigma_{C^*(a,\ind)}(a)}
	\]
	Falls $A$ nicht unital ist, so haben wir
	\[
		\begin{tikzcd}[ampersand replacement=\&]
			C^*(a,\ind_{A^\simm}) \rar["\StackText{\ref{satz16:gelfand}}{\cong}"] \& C \enbrace[\big]{\Omega \enbrace*{C^*(a,\ind_{A^\simm})}} \rar["\StackText{s.o.}{\cong}"]
			\& C \enbrace*{\sigma_{C^*(a,\ind)}(a)} \\
			C^*(a) \rar["\StackText{\ref{satz16:gelfand}}{\cong}"]  \uar[phantom,"\subset", sloped]
			\& \begin{array}{c}
				C_0 \enbrace[\big]{\Omega \enbrace*{C^*(a)}} \StackText{\ref{satz:29}}{=} \\
				C_0 \enbrace[\big]{\Omega \enbrace*{C^*(a,\ind_{A^\simm})} \setminus \set*{\chi_\infty}}
			\end{array}
			  \rar["\cong"] \uar[phantom,"\subset", sloped]
			\& \begin{array}{c}
				C_0 \enbrace*{\sigma_{C^*(a,\ind_{A^\simm})}(a) \setminus \set*{0}} = \\
				\set*{f \in C \enbrace*{\sigma_{C^*(a,\ind_{A^\simm})}(a)} \given f(0)=0}
			\end{array}  \uar[phantom,"\subset", sloped]\\
			a \rar[mapsto] \uar[phantom,"\in", sloped] \& \widehat{a} \rar[mapsto] \uar[phantom,"\in", sloped]\& \id_{\sigma_{C^*(a,\ind_{A^\simm})}(a)} \uar[phantom,"\in", sloped]
		\end{tikzcd}
	\]
	Der Funktionalkalkül $\beta^{-1}$ ist ein ${}^*$-Homomorphismus, daher gilt für ${}^*$-Polynome $g$, dass 
	\[
		g \enbrace[\big]{f(a)} = g \enbrace*{\beta^{-1}(f)}= \beta^{-1}(g \circ f) = (g \circ f) (a)
	\]
	Für eine beliebige Funktion $g\in C \enbrace*{\sigma_{C^*(a,\ind)}(a)}$ approximieren wir $g$ gleichmäßig auf $\sigma_{C^*(a,\ind)}(a)$ durch $^*$-Polynome und benutzen die Stetigkeit von $\beta^{-1}$.
	
	Es bleibt zu zeigen, dass $\sigma_{C^*(a,\ind)}(a)= \sigma_A(a)$ ist (der unitale Fall genügt). 
	Offensichtlich gilt $\Inv \enbrace*{C^*(a,\ind_A)} \subseteq  \Inv(A)$ und somit $\sigma_A(a) \subseteq \sigma_{C^*(a,\ind_A)}(a)$. 
	Um zu zeigen, dass auch die andere Inklusion gilt, sei $b \in C^*(a,\ind_A)$ mit $\norm*{b}\le 1$ und $b \in \Inv(A)$.
	Dann ist $b$ normal, denn $b$ und $b^*$ sind in $C^*(a,\ind_A)$ und diese Algebra ist abelsch.
	$C^*(b,b^{-1})$ ist auch abelsch, da $(b^{-1})^* = (b^*)^{-1}$ und  
	\[
		b(b^*)^{-1} = b(b^*)^{-1} b^{-1} b = b \enbrace*{bb^*}^{-1} b = b(b^*b)^{-1}b = b b^{-1} (b^*)^{-1}b=(b^*)^{-1}b
	\]
	Weiter ist $\norm*{b^*b} \le 1$ nach der $C^*$-Gleichung und $\ind - b^*b$ ist selbstadjungiert. 
	Ferner gilt $0<\chi(b^*b) \le 1$ für alle Charaktere $\chi \in \Omega \enbrace*{C^*(b,b^{-1})}$. 
	Damit ist
	\[
		\sigma_{C^*(b,b^{-1})} \enbrace*{\ind - b^*b} \StackText{\ref{satz:117}}{=} \set*{\chi \enbrace*{\ind - b^*b} \given \chi \in \Omega \enbrace*{C^*(b,b^{-1})}} \subseteq 
		[0,1)
	\]
	Da das Spektrum $\sigma_{C^*(b,b^{-1})}(\ind-b^*b)$ kompakt ist, folgt $1 > r(\ind - b^*b)= \norm*{\ind -b^*b}$. Damit existiert die Neumannsche Reihe
	\[
		(b^*b)^{-1} = \sum_{n=0}^{\infty} (\ind-b^*b)^n \in C^*(b,\ind)
	\]
	und es folgt $b^{-1}=(b^*b)^{-1}b^* \in C^*(b,\ind) \subseteq C^*(a,\ind)$.
	Wenn $b \in C^*(a,\ind)$ invertierbar in $A$ ist, so ist also tatsächlich auch $b^{-1} \in C^*(a,\ind)$, womit die andere Inklusion $\sigma_{C^*(a,\ind)}(a) \subseteq \sigma_A(a)$ folgt. 
	Ebenso folgt $\sigma_{C^*(a,\ind)} \enbrace[\big]{f(a)}=\sigma_A \enbrace[\big]{f(a)}$ für $f \in C \enbrace[\big]{\sigma(a)}$.
\end{beweis}

\begin{korollar}[label=korr:220,{name=[Spektrum in C*-Unteralgebren]}]
	Es seien $A \subseteq B$  $C^*$-Algebren und $a \in A$. Dann gilt
	\[
		\sigma_A(a) \cup \set*{0} = \sigma_B(a) \cup \set*{0}
	\]
	\begin{enumerate}[(i),itemsep=0pt]
		\item Wenn $A$ und $B$ unital sind und $\ind_A=\ind_B$ ist, dann gilt $\sigma_A(a) = \sigma_B(a)$.
		\item Wenn $A$ nicht unital ist und $B$ unital ist, so ist $0 \in \sigma_A(a) = \sigma_B(a)$. 
		\item Wenn beide unital sind, die Einselemente aber nicht übereinstimmen, so gilt $\sigma_A(a) \cup \set*{0} = \sigma_B(a)$.
	\end{enumerate}
\end{korollar}
\begin{beweis} \leavevmode
	\begin{enumerate}[(i)]
		\item Es gilt wieder $\Inv(A) \subset \Inv(B)$, also $\sigma_B(a) \subset \sigma_A(a)$.
		Für $b \in A \cap \Inv(B)$ normal gilt dann $b \in \Inv(A)$ wie im Beweis von \autoref{funkkal}. Damit folgt
		\[
			\sigma_A(a) \subset \sigma_B(a) \implies \sigma_A(a) = \sigma_B(a)
		\]
		Falls $b$ nicht normal ist, betrachten wir $b^*b \in \Inv(B)$. Dann gilt
		\[
			(b^*b)^{-1} b^* = b^{-1} \in A
		\]
		\item Betrachte $A^\simm \cong C^*(a,\ind_B) \subset B$ (algebraisch klar, als $C^*$-Algebren nach \autoref{kor:211}). Dann gilt
		\[
			\overline{\sigma}_A(a) := \sigma_{A^\simm}(a) \StackText{(i)}{=}\sigma_B(a)
		\]
		\item Wenn $\ind_A \neq \ind_B$, so gilt für $a \in \Inv(B)$ auch $a^{-1}\in A$. 
		Dann folgt aber $\ind_B = a^{-1}a\in A$ (Beweis \ref{funkkal}). 
		Dies ist ein Widerspruch, also muss $0 \in \sigma_B(a)$ gelten. 
		Wenn $\lambda \neq 0$, dann gilt $\lambda \cdot \ind_A - a \in \Inv(A)$ genau dann, wenn $\lambda \cdot \ind_B -a \in \Inv(B)$. 
		Dies sieht man wie folgt:
		\begin{align}
			\enbrace*{\lambda \cdot \ind_B- a}^{-1} = \enbrace*{\lambda \cdot \ind_A -a}^{-1} + \lambda^{-1}\enbrace*{\ind_B-\ind_A}
		\end{align}
		Damit ist $\sigma_A(a) \setminus \set*{0} = \sigma_B(a) \setminus \set*{0}$. \qedhere
	\end{enumerate}
\end{beweis}
% section 2 (end)
\newpage

\section{Positivität} % (fold)
\label{sec:3}

\begin{definition}[{name=[Positivität]}]
	Sei $A$ eine $C^*$-Algebra. Ein Element $a \in A$ heißt \Index{positiv} ($a \ge 0$), wenn $a=a^*$ und $\sigma(a) \subseteq [0,\infty)$ ist. Mit $A_+$ bezeichnen wir die 
	Menge der positiven Elemente. Wenn $b,c \in A_{\sa}$ sind, dann definieren wir 
	\[
		b \le c \enspace :\Longleftrightarrow \enspace c-b \ge 0
	\]
\end{definition}

\begin{bemerkung}[{name=[positive Elemente in Unteralgebren und in $C_0(X)$]},label=bem:32] \leavevmode
	\begin{enumerate}[(i),itemsep=2pt]
		\item Für $a \in A \subseteq B$ gilt $a \ge 0$ in $A$ genau dann, wenn $a \ge 0$ in $B$ gilt (nach \autoref{korr:220}).
		\item Für $a \in A = C_0(X)$ ist $a \ge 0$ genau dann, wenn $a(x) \ge 0$ für alle $x \in X$ gilt (nach \ref{bem217}).
	\end{enumerate}
\end{bemerkung}

\begin{proposition}[label=prop:33,{name=[Existenz von Quadratwurzeln]}]
	Sei $A$ eine $C^*$-Algebra und $a \in A_+$. Dann gibt es ein eindeutiges $b \in  A_+$, sodass $b^2=a$. Wir schreiben $a^{\sfrac{1}{2}}$ für dieses Element $b$.
\end{proposition}
\begin{beweis}
	Ein solches $b \in C^*(a) \subseteq A$ existiert nach dem Funktionalkalkül, da $C^*(a) \cong C_0 \enbrace*{\sigma(a) \setminus \set*{0}}$ via $a \mapsto \id_{\sigma(a)} \ge 0$.
	Sei $c \in A_+$ ein weiteres Element mit $c^2=a$. Dann gilt $ca=c^3=ac$ und damit ist $C^*(a,c)$ abelsch. 
	Sei $\Omega := \Omega \enbrace*{C^*(a,c)}$ das Charakterspektrum und $\Phi \colon C^*(a,c) \to C_0(\Omega)$ die Gelfandtransformation. 
	Dann sind $\Phi(b)$ und $\Phi(c)$ Quadratwurzeln von $\Phi(a)$. 
	Anders ausgedrückt gilt für alle $\omega \in \Omega$
	\[
		\Phi(b)(\omega)^2 = \Phi(a)(\omega) = \Phi(c)(\omega)^2
	\]
	Damit folgt $\Phi(b)(\omega)=\Phi(c)(\omega)$ für alle $\omega$ und somit $\Phi(b)=\Phi(c)$, also $b=c$.
\end{beweis}

\begin{proposition}[label=prop:34,{name=[Einfache Abschätzungen für positive Elemente]}]
	Sei $a$ eine unitale $C^*$-Algebra, $a \in A_{\sa}$, $t \in \mathbb{R}_+ = [0,\infty)$. Dann gilt
	\begin{enumerate}[(i),itemsep=2pt]
		\item Aus $\norm*{a-t}\le t$ folgt $a \ge 0$
		\item Aus $a \ge 0$ und $\norm*{a}\le t$ folgt $\norm*{a-t}\le t$ 
	\end{enumerate}
\end{proposition}
\begin{beweis}
	Wir können annehmen, dass $A=C^*(a,\ind)=C \enbrace*{\sigma(a)}$ gilt. (Warum?)
	\begin{enumerate}[(i),itemsep=2pt]
		\item Es gilt $\sup_{x \in \sigma(a)} \abs*{x-t} = \norm*{a-t} \le t$, woraus $\abs*{x-t} \le t $ für alle $x \in \sigma(a)\subseteq \mathbb{R}$ folgt. Damit folgt
		$x \ge 0$ für alle $x \in \sigma(a)$, also $a \ge 0$.
		\item Aus $\norm*{a}\le t$ und $a \ge 0$ folgt $0 \le x \le t$ für alle $x \in \sigma(a)$. 
		Subtrahieren mit $t$ liefert $-t \le x-t \le 0$ für alle $x$, also $\abs*{x-t}\le t $. 
		Damit folgt wieder $\norm*{a-t}\le t$. \qedhere
	\end{enumerate}
\end{beweis}

\begin{korollar}[label=korr:35,{name=[Die positiven Elemente bilden einen Kegel]}]
	Sei $A$ eine $C^*$-Algebra. 
	Dann ist $A_+$ ein Kegel, d.h. für $a,b \in A_+$ und $\lambda, \mu \in \mathbb{R}_+$ gilt 
	\[
		\lambda \cdot a + \mu \cdot b \in A_+ 
	\]
\end{korollar}
\begin{beweis}
	Wir dürfen annehmen, dass $A$ unital ist. Dass $\lambda a \ge 0$ und $\mu b  \ge 0$ gilt, ist klar. Wir müssen also nur zeigen, dass aus $a,b\in A_+$ $a+b\in A_+$ folgt. Es gilt
	\begin{align}
		\norm[\big]{a- \norm*{a}\cdot \ind_A} \le \norm*{a} \quad \text{ und } \quad \norm[\big]{b- \norm*{b} \cdot \ind_A} \le \norm*{b}
	\end{align}
	nach voriger \autoref{prop:34}(ii). Damit folgt weiter aus der Dreiecksungleichung
	\[
		\norm[\big]{a+b - \enbrace*{\norm*{a} + \norm*{b}}} \le \norm[\big]{a - \norm*{a}} + \norm[\big]{b - \norm*{b}} \le \norm*{a} + \norm*{b}
	\]
	Also ist $a+b \ge 0$ nach \autoref{prop:34}(i). 
\end{beweis}

\begin{satz}[label=satz:36,{name=[$a^*a$ ist positiv]}]
	Sei $A$ eine $C^*$-Algebra und $a \in A$. Dann gilt $a^*a \ge 0$.
\end{satz}
\begin{beweis}
	Definiere Funktionen $f,g \colon \mathbb{R} \to \mathbb{R}$ wie in \cref{fig:id_zerlegen} und $h \colon \mathbb{R} \to \mathbb{R}$ durch $h(t)=t^{\sfrac{1}{3}}$.
	\begin{figure}[t]
		\centering{
		\begin{tikzpicture}[scale=0.5]
			\draw[->,thick] (-3,0) -- (3,0);
			\draw[->,thick] (0,-2) -- (0,3);
			\draw[Firebrick2,ultra thick,dashed] (-3,0) -- (0,0) -- (2.8,2.8) node[below]{$f$};
			\draw[SeaGreen4,ultra thick,dashed] (-2.8,2.8) -- (0,0) -- (2.8,0) node[below left]{$g$};
		\end{tikzpicture}
		\caption{Zerlegen der Identität in zwei Funktionen $f$ und $g$, sodass $\id=f-g$} \label{fig:id_zerlegen}
		}
	\end{figure}
	Es gilt $a^*a \in A_{\sa}$, $a^*a = f(a^*a)- g(a^*a)$ und $f(a^*a) \cdot g(a^*a)=0$. 
	Für $c=a \,g(a^*a)$ ergibt sich
	\begin{align}
		-c^*c = -g(a^*a) a^*a g(a^*a) = -g(a^*a) \enbrace[\big]{f(a^*a) - g(a^*a)} g(a^*a) = g(a^*a)^3 \in A_+
	\end{align}
	Nach \autoref{bem:13} gilt $\sigma(-cc^*) \setminus \set*{0} = \sigma(-c^*c) \setminus \set*{0}$, also $-cc^* \in A_+$. 
	Schreibe nun $c=d + i e$ mit $d,e \in A_\sa$. Dann gilt
	\begin{align}
		c^*c+cc^* = 2d^2+ 2e^2
	\end{align}
	Also ist $c^*c= 2d^2+2e^2- cc^* \ge 0$. Damit ist $\sigma(c^*c) \subseteq \mathbb{R}_+ \cap \mathbb{R}_- = \set*{0}$. 
	Also ist $\norm*{c}^2= \norm*{c^*c}= r(c^*c)=0$. 
	Insgesamt ist $g(a^*a)^3=0$, also $g(a^*a)=h \enbrace*{g(a^*a)^3}=0$ und $a^*a=f(a^*a)\ge 0$.
\end{beweis}

\begin{satz}[label=satz:37,{name=[Eigenschaften positiver Elemente]}]
	Sei $A$ eine $C^*$-Algebra. Dann gilt
	\begin{enumerate}[(i),itemsep=2pt]
		\item Es gilt $A_+ = \set*{a^*a \given a \in A}$.
		\item Es gilt $A = \Span A_+$.
		\item $A_+$ ist ein abgeschlossener Kegel in $A$.
		\item $\le$ ist eine partielle Ordnung auf $A_\sa$.
		\item Für $a,b \in A_\sa$ und $c \in A$ mit $a \le b$ gilt $c a c^* \le c b c^*$.
		\item Für $0 \le a \le b$ gilt $\norm*{a} \le \norm*{b}$.
		\item Wenn $A$ unital ist und $a \in A$ positiv und invertierbar, dann folgt aus $a \le b$ stets, dass auch $b$ invertierbar ist und $0 \le b^{-1} \le a^{-1}$ gilt.
	\end{enumerate}
\end{satz}
\begin{beweis}
	\leavevmode
	\begin{enumerate}[(i)]
		\item Folgt aus \autoref{satz:36} und \autoref{prop:33}.
		\item Wenn $a \in A$ ist, dann können wir $a=b+ic$ mit $b,c \in A_\sa$ schreiben. Definiere $b_+ := f(b)$, $b_- := g(b)$, $c_+ := f(c)$, $c_- := g(c)$. Diese liegen alle in 
		$A_+$. Dann gilt $b=b_+ - b_-$, $c=c_+ -c_-$ und $a= b_+ -b_- + i (c_+ -c_-) \in \Span A_+$.
		\item Die Kegeleigenschaft ist klar nach \autoref{korr:35}. 
		Zur Abgeschlossenheit: Sei $a \in \overline{A_+}$. 
		Dann gilt $a=a^*$, da $A_\sa$ abgeschlossen ist.
		Für jedes $\varepsilon>0$ gibt es $b \in A_+$ mit $\norm*{b-a} < \varepsilon$. Dann gilt
		\begin{align}
			\norm*{(a+\varepsilon-b) - \varepsilon} \le \varepsilon
		\end{align}
		und somit $a+ \varepsilon-b \ge 0$ nach \autoref{prop:34}. 
		Dann ist $a+\varepsilon= a+ \varepsilon -b +b \ge 0$ nach \autoref{korr:35} und es folgt $\sigma(a+\varepsilon) \subseteq [0,\infty)$. 
		Damit ist $\sigma(a) \subseteq [\varepsilon,\infty)$. 
		Lassen wir $\varepsilon$ gegen $0$ gehen, so folgt $\sigma(a) \subseteq [0,\infty)$. Insgesamt folgt $a \ge 0$ und damit $\overline{A_+} =A_+$.
		\item \begin{description}[font=\normalfont]
			\item[Reflexivität:] Es gilt offensichtlich $a \le a$.
			\item[Antisymmetrie:] Aus $a \le b$ und $b \le a$ folgt $\sigma(b-a) \subset \mathbb{R}_+$ sowie $\sigma(a-b) \subset \mathbb{R}_+$ und damit folgt $\sigma(a-b)=\sigma(b-a)=\set*{0}$. 
			Da $a-b$ selbstadjungiert ist, folgt aus \autoref{prop:210} $a-b=0$.
			\item[Transitivität:] Gilt $a \le b$ und $b \le c$, so ist 
			\(
				c-a = c-b+b-a \StackText{\ref{korr:35}}{\ge} 0 \implies c \ge a
			\).
		\end{description}
		\item Es gilt 
		\(
			c^*bc- c^*ac = c^*(b-a) c = c^*(b-a)^{\sfrac{1}{2}} (b-a)^{\sfrac{1}{2}} c = d^* d \ge 0
		\) mit $d := (b-a)^{\sfrac{1}{2}} c$.
		\item Sei $0 \le a \le b$. Wir haben $\norm*{b} \cdot \ind_{A^\simm} \ge b$ nach Gelfand für $C^*(b,\ind_{A^\simm})$. 
		Damit ist $a \le \norm*{b} \cdot \ind_{A^\simm}$. 
		Gelfand für $C^*(a,\ind_{A^\simm})$ liefert $\norm*{b} \cdot \ind_{A^\simm} \ge \norm*{a} \cdot \ind_{A^\simm}$. 
		Damit muss dann bereits $\norm*{b} \ge \norm*{a}$ sein.
		\item \emph{Übungsaufgabe 1 b) von Blatt 5! Siehe \cref{sub:beweis_37_vii}.} \qedhere
	\end{enumerate}
\end{beweis}

\begin{satz}[{name=[Die Wurzel ist monoton]}]
	Sei $A$ eine $C^*$-Algebra. Für $0 \le a \le b \in A$ gilt $a^{\sfrac{1}{2}} \le b^{\sfrac{1}{2}}$.
\end{satz}
\begin{beweis}
	Es seien eine reelle Zahl $t >0$ und $c,d \in A_\sa^\simm$ gegeben mit 
	\[
		\enbrace*{t \cdot \ind + b^{\sfrac{1}{2}} + a^{\sfrac{1}{2}}} \enbrace*{t \cdot \ind + b^{\sfrac{1}{2}} - a^{\sfrac{1}{2}}} = c + i d
	\]
	Dann gilt 
	\begin{align}
		c &=\frac{1}{2} \cdot \enbrace*{\enbrace*{t \cdot \ind + b^{\sfrac{1}{2}} + a^{\sfrac{1}{2}}}\enbrace*{t \cdot \ind + b^{\sfrac{1}{2}} -a^{\sfrac{1}{2}}} 
		+ \enbrace*{t \cdot \ind + b^{\sfrac{1}{2}}-a^{\sfrac{1}{2}}} \enbrace*{t \cdot \ind + b^{\sfrac{1}{2}} + a^{\sfrac{1}{2}}} } \\
		&= t^2 \cdot \ind + 2 t \cdot \Underbrace{b^{\sfrac{1}{2}}}{\ge 0} + \Underbrace{b -a}{\ge 0} \ge t^2 \cdot \ind
	\end{align}
	Da $t\neq 0$ ist, folgt $c \in \Inv(A^\simm) \cap A_+$ nach \autoref{satz:37}(vii). Da $c^{-\sfrac{1}{2}}d c^{-\sfrac{1}{2}}$ positiv ist, ist das Spektrum von 
	$\ind + i \cdot c^{-\sfrac{1}{2}}d c^{-\sfrac{1}{2}}$ in $1+ i \cdot \mathbb{R}$ enthalten. Weiter gilt
	\[
		\ind + i \cdot c^{\sfrac{1}{2}}d c^{\sfrac{1}{2}} = c^{-\sfrac{1}{2}}(c + i \cdot d) c^{-\sfrac{1}{2}} \in \Inv(A^\simm)
	\]
	Also ist $c+ i \cdot d$ invertierbar und damit $\enbrace*{t \cdot \ind + b^{\sfrac{1}{2}} - a^{\sfrac{1}{2}}}$ linksinvertierbar. 
	Dieses Element ist aber selbstadjungiert und muss damit dann auch rechtsinvertierbar, also invertierbar sein. 
	Damit folgt nun $-t \notin \sigma \enbrace*{b^{\sfrac{1}{2}} - a^{\sfrac{1}{2}}}$ und somit ist $\sigma \enbrace*{b^{\sfrac{1}{2}} - a^{\sfrac{1}{2}}} \subset \mathbb{R}_+$ und $b^{\sfrac{1}{2}} \ge a^{\sfrac{1}{2}}$.
\end{beweis}

\begin{beispiel}[{name=[Das Quadrat ist nicht monoton]}]
	Im Allgemeinen folgt aus $0 \le a \le b$ \emph{nicht} $a^2 \le b^2$. 
	Betrachte dafür die $C^*$-Algebra $A=M_2$ und die beiden Projektionen
	\[
		p =\begin{pmatrix}
			1 & 0 \\
			0 & 0
		\end{pmatrix} \quad q = \frac{1}{2} \begin{pmatrix}
			1 & 1 \\ 1 & 1
		\end{pmatrix}
	\]
	Dann gilt $p \le p+q$, aber $p^2=p \nleqslant (p+q)^2 = p + q+ pq + qp$, denn
	\[
		q + pq + qp = \frac{1}{2} \enbrace*{\begin{pmatrix}
			1 & 1 \\ 1 & 1
		\end{pmatrix} + \begin{pmatrix}
			1 & 1 \\ 0 & 0
		\end{pmatrix} + \begin{pmatrix}
			1 & 0 \\ 1 &0
		\end{pmatrix}} = \frac{1}{2} \begin{pmatrix}
			3 & 2 \\ 2 & 1
		\end{pmatrix}  
	\]
	Die Matrix hat einen negativen Eigenwert, da die Determinante $-1$ ist.
\end{beispiel}

\begin{bemerkung}[{name=[{wenn Quadrat monoton, dann $A$ abelsch}]}]
	Falls für alle $a,b \in A$ mit $0 \le a \le b$ schon $a^2 \le b^2$ gilt, so ist $A$ bereits abelsch.
\end{bemerkung}

\begin{lemma}[label=lem:311,{name=[gerichtete Menge in einer $C^*$-Algebra]}]
	Sei $A$ eine $C^*$-Algebra. Dann ist $\Lambda := \set*{a \in A_+ \given \norm*{a} < 1}$ gerichtet bezüglich \enquote{$\le$}.
\end{lemma}
\begin{beweis}
	\emph{Im \cref{sub:alternativer_beweis_lemma_311} findet man auch den Beweis von \citeauthor{Murphy} mit zusätzlichen Erklärungen, der meiner Meinung nach etwas \enquote{flüssiger} zu lesen ist.}\\
	Für $a,b \in \Lambda$ setze $a' := a \enbrace*{\ind_{A^\simm}-a}^{-1}$ und $b' := b \enbrace*{\ind_{A^\simm} -b}^{-1}$.
	Durch Betrachtung der \hyperref[satz15_neumann]{von-Neumann-Reihe} sieht man sofort, dass $a'$ und $b'$ positiv sind.
	Dann gilt 
	\[
		c := (a' +b') \enbrace*{\ind_{A^\simm} + a' +b'}^{-1} \in \Lambda
	\]
	 nach Gelfand. Es ist nun
	 \begin{align}
	 	a=a(\ind-a)^{-1} \enbrace*{(\ind-a)^{-1}}^{-1} &\StackText{\ref{bem:13}}{=} a (\ind -a)^{-1} \enbrace*{\ind + a (\ind -a)^{-1}}^{-1} = a' (\ind+a')^{-1} \\
		&= \enbrace[\big]{(\ind +a')-\ind} \enbrace*{\ind +a'}^{-1} = \ind - (\ind +a')^{-1} \\
		\intertext{Da $0\le \ind +a' \le \ind + a' +b'$ jeweils invertierbar sind, ist $(\ind+ a'+b')^{-1} \le (\ind+a')^{-1}$ nach \autoref{satz:37}}
		&\le \ind - (\ind +a' +b')^{-1}\StackText{\ref{bem:13}}{=} \ind - \enbrace*{\ind + (a'+b')(\ind +a' +b')^{-1}} \\
		&= (a'+b')(\ind +a' +b')^{-1}=c
	 \end{align}
	 Damit ist $a \le c$. Die Relation $b \le c$ folgt ebenso.
\end{beweis}

\begin{definition}[{name=[approximative Eins]}]
	Eine \Index{approximative Eins} in einer $C^*$-Algebra $A$ ist (aufsteigendes) Netz $(e_\lambda)_\Lambda \subset A_+^1$ mit $a=\lim_\lambda a e_\lambda$ beziehungsweise 
	$a = \lim_\lambda e_\lambda a$ für jedes $a \in A$.
\end{definition}

\begin{satz}[{name=[Existenz von approximativen Einsen]},label=satz:313]
	Jede $C^*$-Algebra $A$ besitzt eine approximative Eins.
\end{satz}
\begin{beweis}
	Sei $\Lambda := \set*{a \in A_+ \given \norm*{a} < 1}$ wie in \autoref{lem:311} und setze $e_\lambda := \lambda$ für $\lambda \in \Lambda$. 
	Dann ist $(e_\lambda)_\Lambda \subset A_+^1$ ein aufsteigendes Netz (aufsteigend bezüglich \enquote{$\le$}, gerichtet nach \autoref{lem:311}). 
	Wir wollen die geforderte Eigenschaft nun zunächst für $a \in \Lambda$ beweisen.
	
	Sei dazu $1 \ge \varepsilon>0$, $\Omega := \sigma(a) \setminus \set*{0} \subset (0,1]$ und $\Phi \colon C^*(a) \to C_0(\Omega)$ der Gelfandisomorphismus. 
	Betrachte $f := \Phi(a) = \id_\Omega$ und die kompakte Menge $K:= \set*{\omega \in \Omega \given f(\omega)\ge \varepsilon} \subset \Omega$. 
	Nach Urysohn existiert ein $g \colon \Omega \to [0,1]$ stetig mit kompaktem Träger und $g(\omega)=1$ für $\omega \in K$. 
	Wähle $\delta>0$ mit $\delta<1$ und $1-\delta< \varepsilon$. Dann gilt 
	\[
		\norm[\big]{f- \delta \cdot g \,f}_\infty = \norm[\big]{(\ind-\delta \cdot g) f}_\infty \le \varepsilon
	\]
	Mit $\overline{\lambda} := \Phi^{-1}(\delta \cdot g) \in \Lambda$ ergibt sich $\norm*{a- \overline{\lambda} a}_A \le \varepsilon$. 
	Falls nun $\lambda \in \Lambda$ ist mit $\lambda \ge \overline{\lambda}$, so gilt $0 \le a(\ind-e_\lambda)a \le a(\ind- e_{\overline{\lambda}})a$ nach \autoref{satz:37}(v), da $0 \le \ind - e_\lambda \le \ind -e_{\overline{\lambda}}$.
	Dann folgt 
	\begin{align}
		\norm*{a- e_\lambda a}^2 = \norm*{(\ind-e_\lambda)^{\sfrac{1}{2}} (\ind - e_\lambda)^{\sfrac{1}{2}} a}^2 \le \norm*{(\ind- e_\lambda)^{\sfrac{1}{2}}a}^2
		 &= \norm*{a^* (\ind -e_\lambda)^{\sfrac{1}{2}}(\ind-e_\lambda)^{\sfrac{1}{2}}a} \\
		 &\le \norm*{a^* (\ind-e_{\overline{\lambda}})a} \\
		 &\le \norm*{(\ind-e_{\overline{\lambda}})a} \le \varepsilon
	\end{align}
	Damit folgt $a = \lim_\lambda e_\lambda a$ für $a \in \Lambda$. Da $A= \Span \Lambda$ gilt, sind wir fertig.
\end{beweis}

\begin{proposition}[{name=[approximative Einsen in separablen C*-Algbren]}]
	Sei $A$ eine separable $C^*$-Algebra -- das heißt A besitzt eine abzählbare dichte Teilmenge -- dann besitzt $A$ eine abzählbare idempotente approximative Eins $(e_n)_\mathbb{N}$, das heißt es gilt $e_n e_{n+1}=e_n$.
\end{proposition}
\begin{beweis}
	\emph{Übung! (Blatt 5 Aufgaben 3 und 4)}
\end{beweis}
% section 3 (end)
\newpage

\section{Ideale, Quotienten und Limiten} % (fold)
\label{sec:4}

\begin{proposition}[{name=[aus abgeschlossenem Ideal eine $C^*$-Unteralgebra bauen]},label=prop:41]
	Sei $L$ ein abgeschlossenes Linksideal in der $C^*$-Algebra $A$.
	Dann existiert ein aufsteigendes Netz $(e_\lambda)_\Lambda \subset L \cap A_+^1$ mit $a =\lim_\lambda a e_\lambda$ für $a \in L$ und $L \cap L^* \subset A$ ist eine $C^*$-Unteralgebra.
\end{proposition}
\begin{beweis}
	Sei $B := L \cap L^*$, dann ist $B \subset A$ eine $C^*$-Unteralgebra:
	\begin{itemize}[itemsep=0pt]
		\item $L^* \subset A$ abgeschlossen, da $^*$ isometrisch ist. Damit ist auch $L \cap L^*$ abgeschlossen in $A$
		\item $L \cap L^* \subset$ ist ein Untervektorraum.
		\item $L^* \subset A$ ist ein Rechtsideal: $l^* a = (a^*l)^* \in L^*$. Damit ist $L \cap L^*$ bereits eine Unteralgebra: 
		\[
			a,b \in L \cap L^* \implies \Underbracket{\minwidthbox{a}{1em}}{\in L^*} \, \Underbracket{\minwidthbox{b}{1em}}{\in L} \in L \cap L^*
		\]
		\item $L \cap L^* = (L \cap L^*)^*$ ist klar.
	\end{itemize}
	Mit Satz \autoref{satz:313} folgt, dass $B$ eine approximative Eins $(e_\lambda)_\Lambda$ besitzt. 
	Für $a \in L$ ist $a^*a \in B$ und damit 
	\[
		\lim_\lambda \norm*{a-a e_\lambda}  = \lim_\lambda \norm[\big]{(\ind-e_\lambda)a^* a(\ind-e_\lambda)}^{\sfrac{1}{2}} =0 \qedhere
	\]
\end{beweis}

\begin{satz}[{name=[Quotientennorm durch approximative Eins]},label=satz:42]
	Sei $A$ eine $C^*$-Algebra und $J \lhd A$ ein (abgeschlossenes, zweiseitiges) Ideal. 
	Dann ist $J \subset A$ eine $C^*$-Unteralgebra. 
	Falls $(e_\lambda)_\Lambda \subset J$ eine approximative Eins ist, so gilt für $a \in A$ 
	\[
		\norm*{a+J}_{\sfrac{A}{J}} = \lim_\lambda \norm*{a-a e_\lambda} = \lim_\lambda \norm*{a- e_\lambda a}
	\]
	wobei $\norm*{a +J}_{\sfrac{A}{J}} = \inf_{b \in J} \norm*{a+b}$ die Quotientennorm bezeichnet. 
\end{satz}
\begin{beweis}
	Nach \autoref{prop:41} existiert $(u_\gamma)_\Gamma \subset J \cap A_+^1$ mit $a= \lim_\gamma a u_\gamma$ für $a \in J$. Dann gilt
	\[
		a^* = \lim_\gamma (a u_\gamma)^* =\lim_\gamma u_\gamma a^* \in J
	\]
	für $a \in J$, also ist $J=J^*$. Mit \autoref{prop:41} folgt dann wiederum, dass $J=J\cap J^*$ eine $C^*$-Unteralgebra von $A$ ist. 
	Sei nun $(e_\lambda)_\Lambda \subset J$ eine beliebige approximative Eins für $J$. 
	Gegeben $a \in A$ und $\varepsilon>0$, so wähle $b \in J$ mit $\norm*{a+b} \le \norm*{a+ J}_{\sfrac{A}{J}} + \frac{\varepsilon}{2}$. 
	Wähle dann $\lambda_0 \in \Lambda$ mit $\norm*{b- \varepsilon_\lambda b} \le \frac{\varepsilon}{2}$ für alle $\lambda \ge \lambda_0$.
	Für $\lambda \ge \lambda_0$ gilt dann 
	\begin{align}
		\norm*{a+J}_{\sfrac{A}{J}} \le \norm*{a- e_\lambda a} \le \norm[\big]{(\ind-e_\lambda)(a+b)} + \norm[\big]{(\ind-e_\lambda)b} 
		&\le \norm*{a+b} + \norm*{(\ind -e_\lambda) b} \\
		&\le \norm*{a+J}_{\sfrac{A}{J}}+ \frac{\varepsilon}{2} + \frac{\varepsilon}{2}  
	\end{align}
	Damit folgt $\norm*{a-e_\lambda a} \xrightarrow{\minwidthbox{\lambda}{1em}} \norm*{a+J}_{\sfrac{A}{J}}$. 
	$\norm*{a- a e_\lambda } \xrightarrow{\minwidthbox{\lambda}{1em}} \norm*{a+J}_{\sfrac{A}{J}}$ folgt ebenso.
\end{beweis}

\begin{bemerkung}[{name=[Ideal eines Ideals ist ein Ideal]}]
	Für Ideale $J \lhd A$ und $I \lhd J$ gilt $I \lhd A$.
\end{bemerkung}
\begin{beweis}
	Es sei $a \in A$ und $b \in I_+$. Dann gilt
	\(
		ab= \Underbracket{a b^{\sfrac{1}{2}}}{\in J} b^{\sfrac{1}{2}} \in I
	\).
\end{beweis}

\begin{satz}[label=satz:44,{name=[Quotient einer C*-Algebra]}]
	Sei $A$ eine $C^*$-Algebra und $J \lhd A$ ein Ideal. Dann ist $\sfrac{A}{J}$ eine $C^*$-Algebra (mit induzierten Operationen und der Quotientennorm). 
\end{satz}
\begin{beweis}
	Fixiere eine approximative Eins $(e_\lambda)_\Lambda$ für $J$. Sei $a \in A$ und $b \in J$. Dann gilt
	\begin{align}
		\norm*{a+J}^2_{\sfrac{A}{J}} \StackText{\ref{satz:42}}{=} \lim_\lambda \norm*{a-a e_\lambda}^2 &= \lim_\lambda \norm*{(\ind-e_\lambda)a^*a(\ind-e_\lambda)} \\
		&\le \sup_\lambda \norm[\big]{(\ind- e_\lambda)(a^*a +b)(\ind- e_\lambda)} + \Underbracket{\lim_\lambda \norm[\big]{(\ind-e_\lambda)b(\ind-e_\lambda)}}{=0} \\[-.8em]
		&\le \norm*{a^*a+b}
	\end{align}
	Wir erhalten 
	\begin{align}
		\norm*{a+J}_{\sfrac{A}{J}}^2 \le \inf_{b \in J} \norm*{a^* a+b} = \norm*{a^*a +J}_{\sfrac{A}{J}} \le \norm*{a^* +J}_{\sfrac{A}{J}} \cdot \norm*{a+J}_{\sfrac{A}{J}} = \norm*{a+J}^2_{\sfrac{A}{J}}
	\end{align}
	Damit ist $\norm*{(a^* +J)(a+J)}_{\sfrac{A}{J}} = \norm*{a+J}^2_{\sfrac{A}{J}}$ und somit ist die \hyperref[def:22]{$C^*$-Gleichung} erfüllt.
	Dass $\sfrac{A}{J}$ ein Banachraum ist, haben wir bereits in Funktionalanalysis, Proposition 2.16 gesehen. 
\end{beweis}

\begin{proposition}[label=prop:45,{name=[injektive *-Homomorphismen sind isometrisch]}]
	Es seien $A,B$ $C^*$-Algebren, $\pi \colon A \to B$ ein injektiver $^*$-Homomorphismus. Dann ist $\pi$ bereits isometrisch.
\end{proposition}
\begin{beweis}
	Zu zeigen: $\norm*{\pi(a^*a)}=\norm*{\pi(a)}^2= \norm*{a}^2= \norm*{a^*a}$ für $a \in A$. 
	Wir dürfen also annehmen, dass $A= C^*(a^*a)$ und $B=\overline{\pi(A)}$ abelsch sind.\marginnote{$a^*a$ normal $\implies C^*(a^*a)$ abelsch} 
	Weiter können wir $\pi$ durch $\pi^\simm \colon A^\simm  \to B^\sim$ ersetzen, falls $A$ nicht unital ist.
	Wir dürfen also auch $A$ und $B$ als unital annehmen.
	$\pi$ induziert eine stetige Abbildung $\pi^* \colon \Omega(B) \to \Omega(A)$ durch $\chi \mapsto \chi \circ \pi$, also ist $\pi^* \enbrace*{\Omega(B)} \subset \Omega(A)$ kompakt nach \autoref{satz:1.15} und insbesondere auch abgeschlossen.
	Falls $\pi^*$ nicht surjektiv wäre, so gäbe es ein $0 \neq f \in C \enbrace*{\Omega(A)}_+ \cong A_+$ mit $f|_{\pi^*(\Omega(B))} =0$. Dann gilt \marginnote{Gelfand"=Isomorphismus in der Notation unterdrückt}
	\begin{align}
		0=f \enbrace[\big]{\pi^*(\chi)} = \chi \circ \pi(f) 
	\end{align}
	für alle Charaktere $\chi \in \Omega(B)$. 
	Also ist $\pi(f)=0$, woraus mit der Injektivität wieder $f=0$ folgt. 
	Dies ist ein Widerspruch und somit ist $\pi^*$ surjektiv.
	Für $a \in A$ gilt dann
	\[
		\norm*{a}_A \stackrel[\text{Isom.}]{\text{Gelfand-}}{=} \sup_{\rho \in \Omega(A)} \abs*{\rho(a)}= \sup_{\chi \in \Omega(B)} \abs[\big]{\chi \circ \pi(a)} = \norm*{\pi(a)}_B \qedhere
	\]
\end{beweis}

\begin{satz}[{name=[Bild ist $C^*$-Unteralgebra]}]
	Seien $A,B$ $C^*$-Algebren, $\pi \colon A \to B$ ein $^*$-Homomorphismus. Dann ist $\pi(A) \subset B$ eine $C^*$-Unteralgebra.
\end{satz}
\begin{beweis}
	Betrachte die $C^*$-Algebra $\sfrac{A}{\ker \pi}$ (siehe \ref{satz:44}) und den wohldefinierten $^*$-Homomorphismus
	\mapdef{\dot{\pi}\colon\sfrac{A}{\ker \pi_n}}{B}{a + \ker \pi}{\pi(a)}{}
	$\dot{\pi}$ ist injektiv, also isometrisch nach \autoref{prop:45}.
	Damit ist $\pi(A)$ vollständig und $\pi(A) \subset B$ abgeschlossen.
	Also ist $\pi(A)$ eine $C^*$-Unteralgebra von $B$.
\end{beweis}

\begin{proposition}[{name=[Summe von C*-Unteralgebra mit Ideal]}]
	Sei $A$ eine $C^*$-Algebra, $B \subset A$ eine $C^*$-Unteralgebra und $J \lhd A$ ein Ideal. Dann ist $B+J \subset A$ eine $C^*$-Unteralgebra.
\end{proposition}
\begin{beweis}
	\emph{Übung! (Blatt 7, Aufgabe 1)}
\end{beweis}

\begin{definition}[{name=[hereditäre C*-Unteralgebra]}]
	Eine $C^*$-Unteralgebra $B \subset A$ heißt \Index{hereditär}, $B \subher A$, falls gilt
	\[
		a \in A_+, b \in B_+, \enspace a \le b \implies a \in B
	\]
\end{definition}

\begin{satz}[label=satz:49,{name=[Bijektion abgeschlossene Linksideale und hereditäre C*-Algebren]}]
	Für eine $C^*$-Algebra $A$ definiert
	\[
		\begin{array}{rcl}
			\set[\big]{L \subset A \given L \text{ ist abg.Linksideal} } &\xrightarrow{\minwidthbox{ }{2em}} & \set[\big]{B \subher A} \\[0.5ex]
			\textstyle L &\xmapsto{\minwidthbox{\mbox{ }}{2em}} & \textstyle L \cap L^* \\[0.5ex]
			\textstyle \set*{a \in A \given a^*a \in B} &\reflectbox{\ensuremath{\xmapsto{\minwidthbox{\mbox{ }}{2em}}}} & \textstyle B 
		\end{array}
	\]
	eine natürliche Bijektion.
\end{satz}
\begin{beweis}
	Sei $B \subher A$. Dann ist $L_B := \set*{a \in A \given a^*a \in B} \subset A$ ein abgeschlossenes Linksideal: Wenn $a,b \in L_B$, dann gilt
	\begin{align}
		(a^*+b^*)(a+b) \le (a^*+b^*)(a+b)+ (a-b)^*(a-b) = 2a^*a + 2b^*b \in B
	\end{align}
	also $(a^*+b^*)(a+b) \in B$, da $B$ hereditär ist, und somit $a+b \in L_B$. 
	Für $a \in A$ und $b \in L_B$ gilt
	\begin{align}
		(ab)^* ab = b^*a^* ab \le b^* \enbrace*{\norm*{a^*a}\cdot \ind_{A^\simm}} b = \norm*{a^*a} b^*b \in B
	\end{align}
	Da $B$ hereditär ist, folgt $(ab)^*ab \in B$ und somit $ab \in L_B$. 
	Also ist $L_B$ wirklich ein Linksideal und $L_B \subset A$ ist abgeschlossen, da $L_B$ das Urbild von $B$ unter $a \mapsto a^*a$ ist.
	
	Mit \autoref{prop:41} folgt, dass $L_B \cap L_B^* \subset A$ eine $C^*$-Algebra ist. Es gilt
	\[
		0 \le b \in L_B \cap L_B^* \implies 0 \le b^{\sfrac{1}{2}} \in L_B \cap L_B^*
	\]  
	Also ist $b \in B$ und somit $L_B \cap L_B^* \subset B$. 
	Weiter ist klar, dass $B \subset L_B \cap L_B^*$ ist, also insgesamt $B=L_B \cap L_B^*$ gilt. 
	Die Verknüpfung $B \mapsto L_B \mapsto L_B \cap L_B^* $ entspricht also der Identität.
	
	Für die zweite Verknüpfung verifizieren wir zunächst, dass $L \cap L^*$ eine hereditäre $C^*$-Unteralgebra ist: 
	Nach \autoref{prop:41} müssen wir dafür nur noch zeigen, dass $L \cap L^*$ hereditär ist.
	Sei dazu $a \in A_+$ und $0 \le a \le b$ für ein $b \in L \cap L^*$. 
	Sei $(e_\lambda)_\Lambda$ ein Netz wie in \autoref{prop:41}. Es gilt
	\(
		0 \le (\ind-e_\lambda)a(\ind - e_\lambda) \le (\ind -e_\lambda) b (\ind -e_\lambda)
	\)
	für alle $\lambda \in \Lambda$. Damit erhalten wir
	\[
		\norm[\big]{a^{\sfrac{1}{2}}(\ind -e_\lambda)}^2= \norm[\big]{(\ind-e_\lambda)a(\ind -e_\lambda)} \le \norm[\big]{(\ind-e_\lambda)b(\ind -e_\lambda)} \le \norm[\big]{b(\ind -e_\lambda)}
	\]
	Damit ist $a^{\sfrac{1}{2}}=\lim_\lambda a^{\sfrac{1}{2}} e_\lambda$ und damit $a^{\sfrac{1}{2}} \in L$, da $L$ ein Linksideal ist. 
	Dann ist aber auch $a \in L$. Es bleibt nun $L= L_{L \cap L^*}$ zu zeigen.
	$L \subseteq L_{L \cap L^*}$ folgt sofort, da $L$ ein Linksideal und $L^*$ ein Rechtsideal ist.
	Sei nun $a \in L_{L \cap L^*}$ positiv, also $a \in A$ mit $a^*a \in L \cap L^*$.
	Dann gilt $a^2=a^*a \in L \cap L^*$ und somit $a \in L \cap L^* \subset L$, da $L \cap L^*$ eine $C^*$-Algebra ist.
	Wir haben nun $\enbrace*{L_{L \cap L^*}}_+ \subseteq L_+$ und es folgt $L_{L \cap L^*} \subseteq L$.
	Also entspricht auch $L \mapsto L \cap L \mapsto L_{L \cap L^*}$ der Identität.
\end{beweis}

\begin{satz}[{name=[Charakterisierung hereditär]}]
	Sei $B \subset A$ eine $C^*$-Unteralgebra. Dann gilt: 
	\[
		B \subher A \iff bac \in B \text{ für alle } b,c \in B, a \in A 
	\]
\end{satz}
\begin{beweis}
	Für die erste Implikation sei $B = L \cap L^*$ für ein abgeschlossenes Linksideal $L$ nach \autoref{satz:49}. 
	Für $b,c \in B$ und $a \in A$ gilt $bac \in L$ und
	\[
		(bac)^* = c^*a^*b^* \in L
	\]
	Also ist $bac \in L \cap L^*=B$ wie gewünscht. 
	Für die zweite Implikation seien $a \in A_+$, $b \in B_+$ mit $0 \le a \le b$.
	Sei weiter $(e_\lambda)_\Lambda \subset B$ eine approximative Eins für $B$.
	Dann gilt $e_\lambda a e_\lambda \in B$ für $\lambda \in \Lambda$ nach Vorraussetzung.
	Wir haben nun wie im vorigen Beweis
	\[
		0 \le (\ind-e_\lambda)a(\ind - e_\lambda) \le (\ind -e_\lambda) b (\ind -e_\lambda) \grenzw{\lambda} 0
	\]
	Wie eben gilt auch $a^{\sfrac{1}{2}}(\ind -e_\lambda) \grenzw{\lambda} 0$ und  es folgt $B \ni e_\lambda a e_\lambda \to a \in B$.
	Also ist $B$ hereditär.
\end{beweis}

\begin{korollarB}[{name=[Ideale sind hereditär]}]
	Sei $J \lhd A$ ein Ideal, dann gilt $J \subher A$. 
\end{korollarB}

\begin{korollar}[{name=[hereditäre Unteralgebren der Form $\overline{aAa}$]}]
	Sei $A$ eine $C^*$-Algebra, $a \in A_+$ ein positives Element. 
	Dann gilt $\overline{aAa} \subher A$. 
	Falls $A$ separabel ist und $B \subher A$, so gilt $B=\overline{hAh}$ für ein $h \in A_+$.
\end{korollar}
\begin{beweis}
	\emph{Übung! (Blatt5, Aufgaben 3 und 4 zusammen mit Blatt 7, Aufgabe 3)}
\end{beweis}

\begin{satz}[{name=[Limes von C*-Algebren]},label=satz:ger_lim]
	Sei $(A_n,\varphi_n)_{n \in \mathbb{N}}$ ein \Index{induktives System von $C^*$-Algebren}, das heißt die $A_n$ sind $C^*$-Algebren und die $\varphi_n \colon A_n \to A_{n+1}$ sind $^*$-Homomorphismen.
	Dann existiert eine (bis auf Isomorphie) eindeutig bestimmte $C^*$-Algebra $\Dlim(A_n,\varphi_n)$ mit folgender universeller Eigenschaft:
	\begin{itemize}
		\item Es existieren $^*$-Homomorphismen $\varphi_{n,\infty} \colon A_n \to \Dlim(A_n,\varphi_n)$ mit $\varphi_{n,\infty} = \varphi_{n+1,\infty} \circ \varphi_{n}$.
		\item Falls $\psi_n \colon A_n \to B$ für $n \in \mathbb{N}$ $^*$-Homomorphismen sind mit $\psi_n = \psi_{n+1} \circ \varphi_n$, so existiert ein eindeutig bestimmter $^*$-Homomorphismus $\psi \colon \Dlim(A_n,\varphi_n) \to B$ mit $\psi_n = \psi \circ \varphi_{n,\infty}$ für $n \in \mathbb{N}$.
	\end{itemize}
	Kürzer: Folgendes Diagramm kommutiert
	\begin{equation}
		\begin{tikzcd}[column sep=4em,row sep=5.5em,arrows={sloped,pos=.37}]
			& & & & \displaystyle \Dlim(A_n,\varphi_n) \ar[dd,"\exists! \psi",dashed,sloped=false,pos=.5, bend left] \\
			A_0 \rar["\varphi_0",pos=.5]  & \ldots  \rar["\varphi_{n-1}",pos=.5] %\ar[urrrr,"\varphi_{0,\infty}",gray] \ar[drrrr,"\psi_0"',gray]
			& A_n \rar["\varphi_n",pos=.5] \ar[urr,"\varphi_{n,\infty}"] \ar[drr,"\psi_n"',end anchor={[yshift=-0.3ex]}] 
			& A_{n+1} \rar \ar[ur,"\varphi_{n+1,\infty}"'] \ar[dr,"\psi_{n+1}",end anchor={[yshift=0.3ex]}] 
			& \ldots \\
			& & & &  B
		\end{tikzcd} \label{eq:413} \tag{$*$}
	\end{equation}
\end{satz}
\begin{beweis}
	Die Menge
	\[
		\textstyle\prod_\mathbb{N} A_n := \set*{(a_n)_\mathbb{N} \given a_n \in A_n, \sup_n \norm*{a_n} < \infty}
	\]
	ist eine $C^*$-Algebra mit komponentenweisen Operationen und Supremumsnorm, denn sie ist eine Banachalgebra nach Blatt 1, Aufgabe 1 c) und $\norm*{\cdot}_\infty$ ist offensichtlich eine $C^*$-Norm.
	Es ist nun
	\[
		\textstyle\bigoplus_\mathbb{N} A_n := \set[\big]{(a_n)_\mathbb{N} \given a_n \in A_n , \lim_{n \to \infty} \norm*{a_n}=0} \lhd \prod_\mathbb{N} A_n \marginnote{warum?}
	\]
	Sei $q \colon \prod_\mathbb{N} A_n \to \sfrac{\prod_\mathbb{N} A_n}{\bigoplus_\mathbb{N} A_n}$ die Quotientenabbildung und weiter
	\[
		\textstyle C := \set[\big]{(a_n)_\mathbb{N} \in \prod\nolimits_\mathbb{N} A_n \given  \exists n_0 \in \mathbb{N} \text{ mit } a_{n+1}= \varphi_n(a_n) \text{ für } n \ge n_0}
	\]
	$C \subset \prod_\mathbb{N} A_n$ ist eine $^*$-Unteralgebra und folglich ist $\overline{C} \subset \prod_\mathbb{N} A_n$ eine $C^*$-Unteralgebra.
	Wir definieren 
	\[
		\Dlim(A_n,\varphi_n) := q \enbrace*{\overline{C}} \subset \Faktor{\prod\nolimits_\mathbb{N} A_n}{\bigoplus\nolimits_\mathbb{N} A_n}
	\]
	Der Übergang zum Quotienten ist notwendig, damit $\varphi_{n+1,\infty} \circ \varphi_{n} = \varphi_{n,\infty}$ gilt.
	Wir definieren nun $^*$-Homomorphismen $\varphi_n' \colon A_n \to C$ durch 
	\[
		\varphi_n'(a) := \enbrace[\big]{0, \ldots ,0,a,\varphi_n(a), \varphi_{n+1}(\varphi_n(a)),\ldots } 
	\]
	und $\varphi_{n,\infty} := q \circ \varphi_n' \colon A_n \to \Dlim(A_n,\varphi_n)$. 
	Gegeben $B$ und $\psi_n \colon A_n \to B$, so ist $\psi$ auf $q(C)$ eindeutig bestimmt durch \eqref{eq:413}, das heißt $\psi_n = \psi \circ \varphi_{n,\infty} \colon A_n \to B$.
	Wie man dies explizit durchführt, wird in \cref{sub:konstruktion_limes} erläutert.
	Aber $q(C) \subset q(\overline{C})= \Dlim(A_n,\varphi_n)$ ist dicht und $\psi$ ist normvermindernd auf $q(C)$, setzt sich also eindeutig auf $q(\overline{C})$ fort.
	Eindeutigkeit von $\Dlim (A_n,\varphi_n)$ ist klar.
\end{beweis}

\begin{beispiel}[{name=[gerichtete Limites von $C^*$-Algebren]}]
	Beispiele für den gerichteten Limes
	\begin{enumerate}[(i)]
		\item Betrachte
		\[
			\begin{tikzcd}
				M_1 \rar & \ldots  \rar & M_n \rar["\varphi_n"] & M_{n+1} \rar & \ldots \rar & \displaystyle\Dlim (M_n,\varphi_n) \cong \mathcal{K} \enbrace*{\ell^2(\mathbb{N})}
			\end{tikzcd}
		\]
		wobei die $\varphi_n$ eine Matrix rechts unten mit Nullen auffüllt. Es gilt
		\[
			\mathcal{K} \enbrace*{\ell^2(\mathbb{N})} = \overline{\bigcup_\mathbb{N} \Underbrace{\mathcal{B} \enbrace[\big]{\Underbrace{\Span \set*{\eta_0, \ldots , \eta_n}}{\cong \mathbb{C}^{n+1}}}}{\cong M_{n+1}}} \subset \mathcal{B} \enbrace*{\ell^2(\mathbb{N})}
		\]
		nach Korollar 8.11 aus Funktionalanalysis.
		\item Betrachtet man den folgenden Limes
		\[
			\begin{tikzcd}
				M_2 \rar & \ldots  \rar & M_{2^n} \rar["\varphi_n"] & M_{2^{n+1}} \rar & \ldots \rar & \displaystyle\Dlim (M_{2^n},\varphi_n) =: M_{2^\infty}
			\end{tikzcd}
		\]
		wobei die verbindenden Abbildungen die Matrizen zweimal auf der Diagonalen platzieren, also $X \mapsto \begin{psmallmatrix}
			X & 0 \\
			0 & X 
		\end{psmallmatrix}$, so erhält man die sogenannte \Index{CAR-Algebra} $M_{2^\infty}$. 
		Wikipedia zum Namen:
		
		\blockquote[\url{https://de.wikipedia.org/wiki/CAR-Algebra}]{Es handelt sich um eine $C^*$-Algebra, die eng mit den in der Quantenmechanik untersuchten kanonischen Antivertauschungsrelationen (engl. \textbf{c}anonical \textbf{a}nticommutation \textbf{r}elation, daher der Name CAR) verbunden ist und daher auch Fermionenalgebra genannt wird.}
	\end{enumerate}
\end{beispiel}
% section 4 (end)
\newpage

\section{Zustände und die GNS-Konstruktion} % (fold)
\label{sec:5}

\begin{definition}[{name=[positive Funktionale und Zustände]}]
	Sei $A$ eine $C^*$-Algebra. 
	Eine Abbildung $\varphi \colon A \to \mathbb{C}$ heißt \Index{positives Funktional}, falls $\varphi$ linear ist und $\varphi(A_+) \subset \mathbb{R}_+$ gilt. 
	$\varphi$ heißt \Index{Zustand}, falls außerdem $\norm*{\varphi}= 1$ gilt. 
	Wir schreiben $S(A) := \set*{\text{Zustände auf }A}$.\marginnote{engl. \emph{\textbf{s}tates}}
\end{definition}

\begin{proposition}[{name=[positive Sesquilinearformen aus positiven Funktionalen]}]
	Sei $\varphi \colon A \to \mathbb{C}$ ein positives Funktional. Dann ist $(a,b) \mapsto \varphi(a^*b)$ eine positiv semidefinite Sesquilinearform\index{Sesquilinearform} auf $A$.
	% \mapdef{A \times A}{\mathbb{C}}{(a,b)}{\varphi(a^*b)}{}
	% eine positive Sesquilinearform\index{Sesquilinearform}.
	Insbesondere gilt die Cauchy-Schwarz-Ungleichung 
	\[
		\abs[\big]{\varphi(a^*b)} \le \varphi(a^*a)^{\sfrac{1}{2}}\varphi(b^*b)^{\sfrac{1}{2}}
	\]
\end{proposition}
\begin{beweis}
	Das einzige, was man sich kurz überlegen muss, ist die Eigenschaft von $(a,b) \mapsto \varphi(a^*b)$ hermitesch zu sein: Dies folgt mit (ii) aus folgendem \autoref{satz:53}
	\[
		\overline{\varphi (a^* b)} = \varphi \enbrace[\big]{(a^*b)^*} = \varphi(b^*a) \mathrel{\reflectbox{\ensuremath{\longmapsto}}} (b,a) \qedhere
	\]
\end{beweis}

\begin{satz}[label=satz:53,{name=[Eigenschaften positiver Funktionale]}]
	Sei $\varphi$ ein positives Funktional auf $A$. Dann gilt 
	\begin{enumerate}[(i),itemsep=0pt]
		\item $\varphi$ ist beschränkt,
		\item $\varphi(a^*)=\overline{\varphi(a)}$ für $a \in A$ und
		\item $\abs*{\varphi(a)}^2 \le \norm*{\varphi} \cdot \varphi(a^*a)$ für $a \in A$.
	\end{enumerate}
\end{satz}
\begin{beweis}
	\leavevmode
	\begin{enumerate}[(i)]
		\item Falls $\varphi$ unbeschränkt ist, so existiert $(a_n)_\mathbb{N} \subset A_+^1$ mit $\varphi(a_n) >2^n$ für $n \in \mathbb{N}$.
		Setze $h := \sum_{n \in \mathbb{N}} \frac{1}{2^n} \cdot a_n \subset A_+$.
		Dann gilt für alle $k \in \mathbb{N}$
		\begin{align}
			\varphi(h) \ge \varphi \enbrace*{\sum_{n=0}^{k} \frac{1}{2^n} \cdot a_n} > k+1,
		\end{align}
		was ein Widerspruch ist.
		\item Es sei $a=b-c+ i \cdot d - i \cdot e$ mit $b,c,d,e \in A_+$. Dann gilt
		\begin{align}
			\varphi(a^*) = \varphi \enbrace*{b-c- i \cdot d + i \cdot e} &= \varphi(b) - \varphi(c) - i \cdot \varphi(d) + i \cdot \varphi(e) \\
			&= \overline{\varphi(b) - \varphi(c) + i \cdot \varphi(d) - i \cdot \varphi(e)} = \overline{\varphi(a)}
		\end{align}
		\item Sei $(e_\lambda)_\Lambda$ eine approximative Eins. Dann gilt 
		\[
			\abs*{\varphi(a)}^2 = \lim_\lambda \abs*{\varphi(e_\lambda \cdot a)}^2  \StackText{C.S.}{\le}\lim_\lambda \varphi(e_\lambda^2) \cdot \varphi(a^*a) \StackText{\ref{satz:54}}{\le} \norm*{\varphi} \cdot \varphi(a^*a) \qedhere
		\]
	\end{enumerate}
\end{beweis}

\begin{satz}[label=satz:54,{name=[Äquivalenzen zur Positivität von Funktionalen]}]
	Sei $A$ eine $C^*$-Algebra und $\varphi \in A^*$ ein (beschränktes) Funktional. Dann sind äquivalent:
	\begin{enumerate}[(i),itemsep=0pt]
		\item $\varphi$ ist positiv.
		\item $\norm*{\varphi} = \lim_\lambda\varphi(e_\lambda)$ für jede approximative Eins $(e_\lambda)_\Lambda$ in $A$.
		\item $\norm*{\varphi} = \lim_\lambda\varphi(e_\lambda)$ für \emph{eine} approximative Eins $(e_\lambda)_\Lambda$ in $A$.
	\end{enumerate}
	Falls $A$ unital ist, so ist $\varphi$ genau dann positiv, wenn $\varphi(1)=\norm*{\varphi}$ ist.
\end{satz}
\begin{beweis}
	Für die Implikation (i) $\Rightarrow$ (ii) sei $\varphi$ positiv und $(e_\lambda)_\Lambda$ eine approximative Eins. 
	Dann ist $\enbrace*{\varphi(e_\lambda)}_\Lambda \subset \benbrace*{0,\norm*{\varphi}}$ ein aufsteigendes Netz, da $\varphi$ positiv und linear ist, also existiert der Limes und es gilt $\lim_\lambda \varphi(e_\lambda) \le \norm*{\varphi}$. 
	Sei $a \in A^1$. Dann gilt
	\[
		\abs*{\varphi(a)}^2 \grenzwl{\lambda\to\infty} \abs*{\varphi(e_\lambda a)}^2\, \StackText{C.S.}{\le}\, \varphi \enbrace*{e_\lambda^2} \cdot \varphi(a^*a) \stackrel{e_\lambda^2 \le e_\lambda}{\le} \varphi(e_\lambda) \cdot \varphi(a^*a) \le \lim_\lambda \varphi(e_\lambda) \cdot \norm*{\varphi}
	\]
	Damit folgt $\norm*{\varphi}^2 \le \lim_\lambda \varphi(e_\lambda) \cdot \norm*{\varphi} \le \norm*{\varphi}^2$, also $\norm*{\varphi}=\lim_\lambda \varphi(e_\lambda)$.
	
	Die Implikation (ii) $\Rightarrow$ (iii) ist trivial. 
	Für (iii) $\Rightarrow$ (i) sei $(e_\lambda)_\Lambda$ die approximative Eins mit $\norm*{\varphi}= \lim_\lambda \varphi(e_\lambda)$.
	Behauptung: $\varphi(A_\sa) \subset \mathbb{R}$. Sei $a \in A_\sa^1$ mit $\varphi(a) = \alpha + i \cdot \beta$, $\alpha,\beta \in \mathbb{R}$. 
	Wir dürfen $\beta \le 0$ annehmen (betrachte sonst $a^*$). 
	Sei $n \in \mathbb{N}$. Dann gilt 
	\begin{align}
		\norm*{a- i \cdot n \cdot e_\lambda}^2 = \norm[\big]{(a+ i \cdot n \cdot e_\lambda)(a- i \cdot n \cdot e_\lambda)} &= \norm*{a^2 +n^2 \cdot e_\lambda^2 + i \cdot n \cdot \enbrace*{e_\lambda a - a e_\lambda}} \\
		&\le 1+ n^2 + n \cdot \norm*{e_\lambda a - a e_\lambda}
	\end{align}
	Erweitern mit $\norm*{\varphi}^2$ liefert
	\[
		\begin{array}{rlccl}
			&& \mbox{}\hspace{-2em}\mathclap{\abs[\big]{\varphi(a- i \cdot n \cdot e_\lambda)}^2 \le \norm*{\varphi}^2 \enbrace*{1+ n^2 + n \cdot \norm*{e_\lambda a - a e_\lambda}}}\hspace{2cm}\mbox{} && \\
			&\hspace{4em}\rotatebox{270}{\mbox{ }\hspace{.4ex}\ensuremath{\grenzw{\rotatebox{90}{\ensuremath{\lambda}}}}\hspace{1.4ex}\mbox{ }} &&  \rotatebox{270}{\mbox{ }\hspace{.4ex}\ensuremath{\grenzw{\rotatebox{90}{\ensuremath{\lambda}}}}\hspace{1.4ex}\mbox{ }}& \\
			&\abs[\big]{\varphi(a)- i \cdot n \cdot \norm*{\varphi}}^2 && \norm*{\varphi}^2 (1+ n^2) & \\
			=&\abs[\big]{\alpha+ i \cdot \beta - i \cdot n \cdot \norm*{\varphi}}^2 &&& \\
			=&\alpha^2 + \beta^2 - 2 \cdot n \cdot \norm*{\varphi} \cdot \beta + n^2 \norm*{\varphi}^2 &&&
		\end{array}
	\]
	% \begin{alignat}{4}
	% 	&& \abs[\big]{\varphi(a- i \cdot n \cdot e_\lambda)}^2 \le &\norm*{\varphi}^2 \enbrace*{1+ n^2 + n \cdot \norm*{e_\lambda a - a e_\lambda}} && \\
	% 	 \ArrowBetweenLines[{\rotatebox{270}{\ensuremath{\grenzw{\rotatebox{90}{\ensuremath{\lambda}}}}}}] &&& &&\ArrowBetweenLines*[{\rotatebox{270}{\ensuremath{\grenzw{\rotatebox{90}{\ensuremath{\lambda}}}}}}]\\
	% 	\abs[\big]{\varphi(a)- i \cdot n \cdot \norm*{\varphi}}^2 &&& \norm*{\varphi}^2 (1+ n^2)
	% \end{alignat}
	% \begin{align}
	% 	&\abs[\big]{\varphi(a- i \cdot n \cdot e_\lambda)}^2 \le \norm*{\varphi}^2 \enbrace*{1+ n^2 + n \cdot \norm*{e_\lambda a - a e_\lambda}} \grenzw{\lambda}  \norm*{\varphi}^2 (1+ n^2) \\
	% 	\abs[\big]{\varphi(a)- i \cdot n \cdot \norm*{\varphi}}^2 & \grenzwl{\lambda} \\
	% 	=\abs[\big]{\alpha+ i \cdot \beta - i \cdot n \cdot \norm*{\varphi}}^2 & =\alpha^2 + \beta^2 - 2 n \norm*{\varphi} \beta + n^2 \norm*{\varphi}^2
	% \end{align}
	% \begin{align}
	% 	\alpha^2 + \beta^2 - 2 n \norm*{\varphi} \beta + n^2 \norm*{\varphi}^2=\ldots =\abs*{\varphi(a)- i \cdot n \cdot \norm*{\varphi}} &\grenzwl{\lambda}\abs*{\varphi(a- i \cdot n \cdot e_\lambda)}^2 \\
	% 	&\le \norm*{\varphi}^2 \enbrace*{1+ n^2 + n \cdot e_\lambda a - a e_\lambda} \\
	% 	&\grenzw{\lambda}  \norm*{\varphi}^2 (1+ n^2)
	% \end{align}
	Wir erhalten 
	\[
		- 2 \cdot n \cdot \norm*{\varphi} \cdot \beta \le \norm*{\varphi}^2 - \alpha^2 -\beta^2
	\]
	Damit folgt $\norm*{\varphi}= 0$ oder $\beta\ge 0$. Insgesamt also in jedem Fall $\varphi(A_\sa) \subset \mathbb{R}$.
	Sei $a \in A_+^1$. Dann gilt $e_\lambda-a \in A_\sa$ und $\norm*{e_\lambda-a} \le 1$ (denn $- \ind_{A^\simm}\le e_\lambda-a \le \ind_{A^\simm}$).
	Also ist $\varphi(e_\lambda-a) \le \abs*{\varphi(e_\lambda-a)} \le \norm*{\varphi} \cdot \norm*{e_\lambda-a} \le \norm*{\varphi}$.
	Dann gilt aber
	\begin{align}
		\norm*{\varphi}- \varphi(a) = \lim_{\lambda} \varphi(e_\lambda) - \varphi(a) = \lim_\lambda \varphi \enbrace*{e_\lambda-a} \le \norm*{\varphi}
	\end{align}
	Also ist $\varphi(a) \ge 0$ und wir sind fertig.
\end{beweis}

\begin{satz}[{name=[Existenz von Zuständen zu gegebenen normalen Elementen]},label=satz:55]
	Sei $A$ eine $C^*$-Algebra und $0\neq a \in A$ normal.
	Dann existiert ein Zustand $\varphi$ auf $A$ mit $\norm*{a}=\abs*{\varphi(a)}$.
\end{satz}
\begin{beweis}
	Die $C^*$-Algebra $B := C^*(a, \ind)\subset A^\simm$ ist abelsch, da $a$ normal ist. 
	Insbesondere existiert ein Charakter $\varphi_2$ auf $B$ mit $\abs*{\varphi_2(a)}=\norm*{a}$.
	Nach Hahn-Banach können wir $\varphi_2$ stetig und linear zu $\varphi_1$ auf $A^\simm$ fortsetzen mit $\norm*{\varphi_1} = \norm*{\varphi_2}=1$.
	Weiter gilt $\varphi_1(\ind)= \varphi_2(\ind)=1$ und damit ist $\varphi_1$ ein stetiges lineares Funktional auf $A^\simm$ mit 
	\[
		\norm*{\varphi_1} = \norm*{\varphi_2}\StackText{\ref{satz:54}}{=} \varphi_2(\ind) = \varphi_1(\ind)
	\]
	Mit \autoref{satz:54} folgt dann auch, dass $\varphi_1$ positiv ist. $\varphi := \varphi_1|_A$ ist also positiv mit $\norm*{a}= \abs*{\varphi(a)}$ wie gewünscht.
	Weiter ist $\norm*{\varphi} \le \norm*{\varphi_1} \le 1$ und außerdem 
	\[
		\norm*{a} = \abs*{\varphi(a)} \le \norm*{\varphi} \cdot \norm*{a} \implies \norm*{\varphi} \ge 1
	\]
	Also ist $\varphi$ auch ein Zustand.
\end{beweis}

\begin{proposition}[label=prop:56,{name=[{positive Funktionale und Elemente der Form $a^*a$}]}]
	Sei $\varphi$ ein positives stetiges Funktional auf $A$. 
	\begin{enumerate}[(i),itemsep=0pt]
		\item Für alle $a \in A$ gilt $\varphi(a^*a)=0 \iff \varphi(ba)=0$ für alle $b \in A$.
		\item Für alle $a,b \in A$ gilt: $\varphi(b^*a^*ab) \le \norm*{a^*a}\varphi(b^*b)$.
	\end{enumerate}
\end{proposition}
\begin{beweis}
	\leavevmode
	\begin{enumerate}[(i),itemsep=0pt]
		\item Die nicht triviale Implikation folgt mit Cauchy-Schwarz.
		\item Benutze $b^*a^*ab \le b^* \enbrace*{\norm*{a^*a} \cdot \ind_{A^\simm}}b = \norm*{a^*a} \cdot b^*b$. \qedhere
	\end{enumerate}
\end{beweis}

\begin{satz}[label=satz:57,{name=[positive Fortsetzung positiver Funktionale]}]
	Sei $B \subset A$ eine $C^*$-Unteralgebra und $\varphi$ ein positives Funktional auf $B$. 
	Dann besitzt $\varphi$ eine positive Fortsetzung $\overline{\varphi}$ auf $A$ mit $\norm*{\overline{\varphi}} = \norm*{\varphi}$.
\end{satz}
\begin{beweis}
	Definiere $\varphi^+ \colon B^+ \to \mathbb{C}$ durch $\varphi^+ \enbrace*{(b,\mu)} := \varphi(b) + \mu \cdot \norm*{\varphi}$.\marginnote{vgl. \ref{bem:27} für $B^+$}
	Für $(b,\mu) \in B^+$ erhalten wir 
	\begin{align}
		\abs*{\varphi^+ \enbrace*{(b,\mu)}} = \abs[\big]{\varphi(b) + \mu \cdot \norm*{\varphi}} = \abs[\Big]{\lim_\lambda \varphi \enbrace*{b e_\lambda} + \mu \cdot \lim_\lambda \varphi(e_\lambda) } &= \abs[\Big]{\lim_\lambda \varphi \enbrace[\big]{(b+ \mu \cdot \ind_{B^+})e_\lambda}} \\
		&\le \norm*{\varphi} \cdot  \norm*{b+ \mu \cdot \ind_{B^+}}_{B^+} \\
		&= \norm*{\varphi} \cdot \norm*{(b,\mu)}
	\end{align}
	wobei $(e_\lambda)_\Lambda$ eine approximative Eins für $B$ ist. 
	Damit ist $\norm*{\varphi^+} \le \norm*{\varphi} \enbrace*{\le \norm*{\varphi^+}}$ und es folgt $\norm*{\varphi^+}=\norm*{\varphi}=\varphi^+(\ind_{B^+})$ noch Definition von $\varphi^+$.
	Mit \autoref{satz:54} folgt, dass $\varphi^+$ positiv ist.
	
	Die Einbettung $\iota \colon B \hookrightarrow A \hookrightarrow A^\simm$ setzt sich zu einem unitalen $^*$-Homomorphismus $\iota^+ \colon B^+ \to A^\simm$ fort mittels
	\[
		\enbrace*{\iota^+ \enbrace*{b,\mu}} := \iota(b) + \mu \cdot \ind_{A^\simm}
	\]
	Definiere nun $\hat{\varphi} \colon \iota^+(B^+) \to \mathbb{C}$ durch $\hat{\varphi} \enbrace*{\iota^+(b,\mu)} := \varphi^+(b,\mu)$. 
	Dies ist wohldefiniert, denn nur der Fall $B\ni \ind_B= \ind_A \in A$ ist zu überprüfen.
	Dann ist $\hat{\varphi}$ positiv, denn $\varphi^+$ ist positiv und $\enbrace*{\iota^+(B^+)}_+ = \iota^+(B^+_+) $ (warum?).
	Nach Hahn-Banach existiert eine stetige lineare Fortsetzung $\hat{\hat{\varphi}} \colon A^\simm  \to \mathbb{C}$ von $\hat{\varphi}$ mit $\norm*{\hat{\hat{\varphi}}}=\norm*{\hat{\varphi}}=\hat{\varphi}(\ind_{A^\simm})= \hat{\hat{\varphi}}(\ind_{A^\simm})$ und es folgt wieder mit \autoref{satz:54}, dass $\hat{\hat{\varphi}}$ positiv ist.
	Dann ist $\overline{\varphi} := \hat{\hat{\varphi}}|_A$ die gesuchte positive Fortsetzung. 
\end{beweis}

\begin{satz}[{name=[Eindeutigkeit der Fortsetzung im hereditären Fall]},label=satz:58]
	Sei $B \subher A$ eine hereditäre $C^*$-Unteralgebra und $\varphi$ ein positives Funktional auf $B$. 
	Dann besitzt $\varphi$ eine \emph{eindeutige} positive lineare Fortsetzung $\overline{\varphi}$ auf $A$ mit $\norm*{\overline{\varphi}}=\norm*{\varphi}$.
	Falls $(e_\lambda)_\Lambda \subset B$ eine approximative Eins für $B$ ist, so gilt $\overline{\varphi}(a) = \lim_\lambda \varphi(e_\lambda a e_\lambda)$ für $a \in A$.
\end{satz}
\begin{beweis}
	\emph{Eine} positive Fortsetzung $\overline{\varphi}$ existiert nach \autoref{satz:57}. 
	Sei nun $\psi \colon A \to \mathbb{C}$ irgendeine positive lineare Fortsetzung von $\varphi$ mit $\norm*{\psi}=\norm*{\varphi}$.
	Sei $\psi^\simm \colon A^\simm \to \mathbb{C}$ eine (die) positive lineare Fortsetzung von $\psi$ mit $\norm*{\psi}=\norm*{\psi^\simm}$ nach \autoref{satz:57}.
	Wir haben $\lim_\lambda \varphi(e_\lambda)=\norm*{\varphi} = \norm*{\psi^\simm} = \psi^\simm(\ind_{A^\simm})$, also ist 
	\[
		\lim_\lambda \psi^\simm(\ind_{A^\simm} -e_\lambda)=0
	\]
	Für  $a \in A$ ergibt sich nun
	\begin{align}
		\abs[\big]{\psi(a)- \varphi(e_\lambda a e_\lambda)} &\le \abs[\big]{\psi^\simm(a)- \psi^\simm(e_\lambda a)} + \abs[\big]{ \psi^\simm \enbrace*{e_\lambda a - e_\lambda a e_\lambda}}\\
		&= \abs[\big]{\psi^\simm \enbrace[\big]{(\ind -e_\lambda)a}} + \abs[\big]{\psi^\simm \enbrace[\big]{e_\lambda a(\ind -e_\lambda)}} \\ 
		&\StackTextClap{C.S.}{\le} \psi^\simm \enbrace*{(\ind-e_\lambda)^2}^{\sfrac{1}{2}} \cdot \psi^\simm \enbrace*{a^*a}^{\sfrac{1}{2}} + \psi^\simm(a^* e_\lambda^2 a)^{\sfrac{1}{2}} \cdot \psi^\simm \enbrace*{(\ind-e_\lambda)^2}^{\sfrac{1}{2}} \\
		&\le \Underbracket{\psi^\simm (\ind-e_\lambda)^{\sfrac{1}{2}}}{\grenzw{} 0}\cdot  \psi^\simm (a^*a)^{\sfrac{1}{2}} + \psi^\simm (a^*e_\lambda^2 a) \cdot  \Underbracket{\psi^\simm(\ind-e_\lambda)^{\sfrac{1}{2}}}{\grenzw{} 0}
	\end{align}
	Damit folgt $\psi(a) = \lim_\lambda \varphi(e_\lambda a e_\lambda)$ und die Eindeutigkeit.
\end{beweis}

\begin{beispiel}[{name=[{Zustände,Spurzustände}]}]
	\leavevmode
	\begin{enumerate}[(i)]
		\item Folgende Abbildung ist ein Zustand aber kein Charakter, da das Integral nicht multiplikativ ist Positivität mit \autoref{bem:32}:
		\mapdef{C([0,1])}{\mathbb{C}}{f}{\int f \,\mathd \lambda}{}
		\item Die normalisierte Spurabbildung
		\mapdef{M_n}{\mathbb{C}}{X}{\tr X}{}
		bezeichnet man als \Index{Spurzustand}, was allgemein ein Zustand $\tau \colon A \to \mathbb{C}$ ist mit $\tau(a^*a)=\tau(aa^*)$ für alle $a \in A$.
		\item Sei 
		$\begin{tikzcd}[cramped,sep=small]
			A_0 \rar & A_1 \rar & \ldots \rar & A =\Dlim\, A_i
		\end{tikzcd}$
		ein unitales induktives System von $C^*$-Algebren.
		Falls jedes $A_i$ einen Spurzustand $\tau_i$ besitzt, so auch $A$: Wir haben
		\[
			A \xhookrightarrow{\minwidthbox{}{1.4em}} \Faktor{\prod_\mathbb{N} A_i}{\bigoplus_\mathbb{N}(A_i)}
		\]
		Für $a=\benbrace*{(a_i)_\mathbb{N}}\in A \subset \sfrac{\prod_\mathbb{N} A_i}{\bigoplus_\mathbb{N}(A_i)}$ definieren wir $\tau(a) := \lim_\omega \tau_i(a_i)$, wobei $\omega \in \beta \mathbb{N} \setminus \mathbb{N}$ ein freier \hrefsymX{https://de.wikipedia.org/wiki/Ultrafilter}{Ultrafilter} ist.
		Man verifiziert nun leicht, dass $\tau$ wohldefiniert, unital, positiv und linear ist und außerdem die Spurbedingung erfüllt.  
		
		Alternativ kann man auch wie folgt vorgehen: Betrachte folgendes Diagramm
		\[
			\begin{tikzcd}
				& (a_i)_\mathbb{N} \rar[mapsto] & \enbrace*{\tau_i(a_i)}_\mathbb{N} \\
				&\prod_\mathbb{N} A_i \rar \dar & \ell^\infty(\mathbb{N}) \dar \\
				A \rar[hook] &\Faktor{\prod_\mathbb{N} A_i}{\bigoplus_\mathbb{N}(A_i)} \rar & \Faktor{\ell^\infty(\mathbb{N})}{c_0(\mathbb{N})} \rar["\chi"] & \mathbb{C}
			\end{tikzcd}
		\]
		Die Menge rechts unten ist eine abelsche, unitale $C^*$-Algebra. Wähle einen Charakter $\chi$. Man sieht auch, dass die freien Ultrafilter in 1-zu-1-Zusammenhang zu $\beta \mathbb{N} \setminus \mathbb{N}$ stehen.
	\end{enumerate}
\end{beispiel}

\begin{definition}[{name=[Darstellung]},label=def:510]
	Sei $A$ eine $C^*$-Algebra. 
	Eine \Index{Darstellung} von $A$ ist ein Paar $(\mathcal{H},\pi)$, wobei $\mathcal{H}$ ein Hilbertraum ist und $\pi$ ein $^*$-Homomorphismus $A \to \mathcal{B}(\mathcal{H})$.
	\begin{itemize}[itemsep=0pt]
		\item $(\mathcal{H},\pi)$ heißt \bet{treu}\index{Darstellung!treu}, falls $\pi$ injektiv ist.
		\item $(\mathcal{H},\pi,\xi)$ heißt \bet{zyklisch}\index{Darstellung!zyklisch}, falls $(\mathcal{H},\pi)$ eine Darstellung ist und $\xi \in \mathcal{H}$ ein Vektor mit $\overline{\pi(A) \xi}=\mathcal{H}$.
		\item $(\mathcal{H},\pi)$ heißt \bet{algebraisch irreduzibel}\index{Darstellung!algebraisch irreduzibel}, falls gilt: 
		Für einen beliebigen Untervektorraum $\mathcal{H}_0 \subset \mathcal{H}$ mit ${\pi(A)\mathcal{H}_0 \subset \mathcal{H}_0}$ ist $\mathcal{H}_0=\set{0}$ oder $\mathcal{H}_0=\mathcal{H}$.
		\item $(\mathcal{H},\pi)$ heißt \bet{(topologisch) irreduzibel}\index{Darstellung!irreduzibel (topologisch)}, falls gilt: 
		Für einen abgeschlossenen Untervektorraum $\mathcal{H}_0 \subset \mathcal{H}$ mit $\pi(A)\mathcal{H}_0 \subset \mathcal{H}_0$ ist $\mathcal{H}_0=\set{0}$ oder $\mathcal{H}_0=\mathcal{H}$.
		\item $(\mathcal{H},\pi)$ heißt \bet{nicht degeneriert}\index{Darstellung!nicht degeneriert}, falls $\overline{\pi(A)\mathcal{H}}=\mathcal{H}$. Äquivalent: $\pi(A)\xi = \set{0} \implies \xi=0$.
	\end{itemize}
\end{definition}

\begin{satz}[{name={GNS-Konstruktion}},label=gns]
	Sei $A$ eine $C^*$-Algebra und $\varphi \colon A \to \mathbb{C}$ ein positives Funktional. 
	Wir setzen 
	\[
		N_\varphi := \set*{a \in A \given \varphi(a^*a)=0}\StackText{\ref{prop:56}(i)}{=}\set*{a \in A \given \varphi(ba)=0 \text{ für alle }b \in A}
	\]
	Dann ist $N_\varphi \subset A$ ein abgeschlossenes Linksideal. 
	Definiere eine Abbildung
	\mapdef{\Faktor{A}{N_\varphi} \times \Faktor{A}{N_\varphi}}{\mathbb{C}}{(a+N_\varphi,b+N_\varphi)}{\skal*{a+N_\varphi}{b+N_\varphi}:= \varphi(a^*b)}{}
	$\skal*{\cdot}{\cdot}$ ist dann ein wohldefiniertes Skalarprodukt auf $\sfrac{A}{N_\varphi}$:
	Für die Wohldefiniertheit seien $a+N_\varphi=c+N_\varphi$ und $b+N_\varphi=d+N_\varphi$.
	Dann gilt $a-c \in N_\varphi$ und $b-d \in N_\varphi$. Damit folgt dann
	\begin{align}
		\abs[\big]{\varphi(a^*b)-\varphi(c^*d)} \le \abs[\big]{\varphi(a^*b)-\varphi(c^*b)} + \abs[\big]{\varphi(c^*b)- \varphi(c^*d)} &= \abs*{\varphi \enbrace[\big]{(a^*-c^*)b}} + \abs*{\varphi \enbrace[\big]{c^*(b-d)}} \\[-1em]
		&\StackTextClap{\ref{satz:53}}{=} \abs[\Big]{\overline{\varphi \enbrace[\big]{b^*(a-c)}}} + \abs*{\varphi \enbrace[\big]{c^*(b-d)}}\\ &=0
	\end{align}
	Dass $\skal*{\cdot}{\cdot}$ sesquilinear ist, ist klar.
	Die Antisymmetrie folgt mit $\varphi(a^*b)=\overline{\varphi(b^*a)}$ (nach \autoref{satz:53}). 
	Die positive Definitheit rechnet man einfach wie folgt nach:
	\[
		\skal*{a+N_\varphi}{a+N_\varphi}= 0 \iff \varphi(a^*a)=0 \iff a \in N_\varphi \iff a+ N_\varphi = N_\varphi =0 \in \Faktor{A}{N_\varphi}
	\]
	Insgesamt ist $\sfrac{A}{N_\varphi}$ also ein Prä-Hilbertraum mit $\skal*{\cdot}{\cdot}$.
	Sei $\mathcal{H}_\varphi := \overline{\sfrac{A}{N_\varphi}}^{\skal*{\cdot}{\cdot}}$ die Hilbertraumvervollständigung und definiere $\dot{\pi} \colon A \to \mathcal{B}(\sfrac{A}{N_\varphi})$ durch $\dot{\pi}(a)(b+N_\varphi) := ab + N_\varphi$ für $a,b \in A$.
	Dann ist $\dot{\pi}$ ein wohldefinierter kontraktiver Homomorphismus von Algebren:
	Da $N_\varphi$ ein Linksideal ist, ist $\dot{\pi}(a)$ linear auf $\sfrac{A}{N_\varphi}$. Weiter ist $\dot{\pi}$ linear und multiplikativ als Abbildung $A \to \mathcal{L}(\sfrac{A}{N_\varphi})$. Kontraktivität:
	\begin{align}
		\norm[\big]{\dot{\pi}(a)(b+N_\varphi)}^2_{\sfrac{A}{N_\varphi}} = \skal*{ab+N_\varphi}{ab+N_\varphi} = \varphi \enbrace*{b^*a^*ab} &\StackTextClap{\ref{prop:56}(ii)}{\le} \norm*{a}_A^2 \cdot \varphi(b^*b) \\
		&= \norm*{a}_A^2 \cdot \norm*{b+N_\varphi}^2_{\sfrac{A}{N_\varphi}}
	\end{align}
	Damit gilt $\norm*{\dot{\pi}(a)(b+N_\varphi)}_{\sfrac{A}{N_\varphi}} \le \norm*{a}_A \cdot \norm*{b+N_\varphi}_{\sfrac{A}{N_\varphi}}$, also $\norm*{\dot{\pi}(a)}_{\mathcal{B}(\sfrac{A}{N_\varphi})} \le \norm*{a}_A$.
	Insgesamt ist $\dot{\pi}\colon A \to \mathcal{B}(\sfrac{A}{N_\varphi})$ also ein kontraktiver Homomorphismus.
	Für jedes $a \in A$ ist $\norm*{\dot{\pi}(a)}$ beschränkt durch $\norm*{a}$, daher besitzt $\dot{\pi}(a)$ eine eindeutig bestimmte Fortsetzung $\pi_\varphi(a) \in \mathcal{B}\enbrace{\mathcal{H}_\varphi}$ mit $\norm*{\pi_\varphi(a)} = \norm*{\dot{\pi}(a)}\le \norm*{a}$.
	Wir zeigen nun, dass $\pi_\varphi \colon A \to \mathcal{B}(\mathcal{H}_\varphi)$ ein $^*$-Homomorphismus und damit nach \autoref{kor:211} kontraktiv ist.
	Dazu zeigen zunächst $\dot{\pi}(a^*)=\dot{\pi}(a)^*$. Für $b +N_\varphi,c+N_\varphi \in \sfrac{A}{N_\varphi}$ gilt
	\[
		\skal[\big]{c+N_\varphi}{\dot{\pi}(a) b+N_\varphi} = \varphi(c^*ab) = \varphi \enbrace[\big]{(a^*c)^*b} = \skal*{a^*c +N_\varphi}{b+N_\varphi}= \skal[\big]{\dot{\pi}(a^*)c+N_\varphi}{b+N_\varphi}
	\]
	Da $\dot{\pi}$ stetig fortgesetzt wird, gilt dann auch $\pi_\varphi(a^*)=\pi_\varphi(a)^*$.
	
	$(\mathcal{H}_\varphi,\pi_\varphi)$ ist die \Index{GNS-Darstellung}\footnote{nach \href{https://de.wikipedia.org/wiki/Israel_Moissejewitsch_Gelfand}{\textsc{\textbf{\emph{G}}elfand}}, \href{https://de.wikipedia.org/wiki/Mark_Neumark}{\textsc{\textbf{\emph{N}}aimark}}, \href{https://de.wikipedia.org/wiki/Irving_Segal}{\textsc{\textbf{\emph{S}}egal}}} von $\varphi$.
	$(\mathcal{H}_\varphi,\pi_\varphi,\xi_\varphi)$ ist zyklisch mit $\xi_\varphi = \lim_\lambda e_\lambda+N_\varphi$, wobei $(e_\lambda)_\lambda$ eine approximative Eins von $A$ ist. \emph{Übungsaufgabe 2 von Blatt 8, siehe \cref{sub:gns_zyklisch}}
\end{satz}

\begin{satz}[{name=[Existenz treuer Darstellungen]},label=satz:512]
	Jede $C^*$-Algebra $A$ besitzt eine treue Darstellung, die \Index{universelle Darstellung}. 
	Falls $A$ separabel ist, so existiert eine treue Darstellung auf einem separablen Hilbertraum.
\end{satz}
\begin{beweis}
	Wir definieren die universelle Darstellung $(\mathcal{H}_u,\pi_u)$ durch 
	\[
		\mathcal{H}_u := \bigoplus_{\varphi \in S(A)} \!\!\!\mathcal{H}_\varphi = \set*{(\eta_\varphi)_{\varphi \in S(A)} \given \eta_\varphi \in \mathcal{H}_\varphi, \sum\nolimits_{\varphi} \skal*{\eta_\varphi}{\eta_\varphi} <\infty}
	\]
	und $\pi_u \colon A \to \mathcal{B}(\mathcal{H}_u)$ mit $\pi_u (a) \enbrace[\big]{(\eta_\varphi)_\varphi} := \enbrace[\big]{\pi_\varphi(a) \eta_\varphi}_\varphi$ für $a \in A$.
	Sei $0 \neq a \in A^1$. Dann existiert nach \autoref{satz:55} ein $\overline{\varphi} \in S(A)$ mit $\overline{\varphi}(a^*a)=\norm*{a^*a}=\norm*{a}^2$.
	Es gilt nun
	\begin{align}
		\norm*{a}^2 = \overline{\varphi}(a^*a) \StackTextClap{\ref{sub:gns_zyklisch}}{=}  \skal[\big]{\pi_{\overline{\varphi}} (a^*a) \xi_{\overline{\varphi}}}{\xi_{\overline{\varphi}}}
		\StackTextClap{C.S.}{\le} \norm[\big]{\pi_{\overline{\varphi}}(a^*a)} \cdot \norm*{\xi_{\overline{\varphi}}}^2_{\mathcal{H}_{\overline{\varphi}}} 
		&= \norm[\big]{\pi_{\overline{\varphi}}(a^*a)} \cdot \skal*{\xi_{\overline{\varphi}}}{\xi_{\overline{\varphi}}} \\
		&=\norm[\big]{\pi_{\overline{\varphi}}(a^*a)} \cdot \overline{\varphi} \enbrace*{\lim\nolimits_\lambda e_\lambda^2}\\
		&\StackTextClap{\ref{satz:54}}{=} \norm*{\pi_{\overline{\varphi}}(a)}^2 \cdot \Underbracket{\norm*{\overline{\varphi}}}{=1}
	\end{align}
	Da $\pi_\varphi$ nach \autoref{kor:211} kontraktiv ist, folgt also $\norm*{\pi_\varphi(a)}= \norm*{a}$.
	Damit erhalten wir weiter $\norm*{a} = \norm*{\pi_\varphi(a)} \le \norm*{\pi_u(a)} \le \norm*{a}$ und somit ist $\pi_u$ als Isometrie treu.
	% \begin{align}
	% 	\norm*{a}^2 = \overline{\varphi}(a^*a) &= \overline{\varphi} \enbrace*{ (a^*a)^{\sfrac{1}{2}} (a^*a)^{\sfrac{1}{2}}}\\
	% 	&= \skal*{(a^*a)^{\sfrac{1}{2}}+N_{\overline{\varphi}}}{(a^*a)^{\sfrac{1}{2}}+N_{\overline{\varphi}}}_{\mathcal{H}_{\overline{\varphi}}} \\
	% 	&= \skal*{\pi_{\overline{\varphi}}\enbrace[\big]{(a^*a)^{\sfrac{1}{4}}} \enbrace[\big]{(a^*a)^{\sfrac{1}{4}} + N_{\overline{\varphi}}}}{\pi_{\overline{\varphi}} \enbrace[\big]{(a^*a)^{\sfrac{1}{4}}} \enbrace[\big]{(a^*a)^{\sfrac{1}{4}}+N_{\overline{\varphi}}}}_{\mathcal{H}_{\overline{\varphi}}} \\
	% 	&= \norm*{\pi_{\overline{\varphi}} \enbrace[\big]{(a^*a)^{\sfrac{1}{4}}} \enbrace[\big]{(a^*a)^{\sfrac{1}{4}}+N_{\overline{\varphi}}}}^2_{\mathcal{H}_{\overline{\varphi}}} \\
	% 	&\le \norm*{\pi_{\overline{\varphi}} \enbrace[\big]{(a^*a)^{\sfrac{1}{4}}}}^2_{\mathcal{B}(\mathcal{H}_{\overline{\varphi}})}  \cdot \norm*{(a^*a)^{\sfrac{1}{4}}+ N_{\overline{\varphi}}}_{\mathcal{H}_{\overline{\varphi}}}^2\\
	% 	% &\le \norm*{\pi_{\overline{\varphi}} \enbrace[\big]{(a^*a)^{\sfrac{2}{4}}}}_{\mathcal{B}(\mathcal{H}_{\overline{\varphi}})} \\
	% 	% &\le \norm*{\pi_u \enbrace*{(a^*a)^{\sfrac{1}{2}}}}_{\mathcal{B}(\mathcal{H}_u)} = \norm*{\pi_u(a)}
	% 	&= \norm*{\pi_{\overline{\varphi}} \enbrace[\big]{(a^*a)^{\sfrac{2}{4}}}}_{\mathcal{B}(\mathcal{H}_{\overline{\varphi}})} \cdot \abs*{\overline{\varphi}\enbrace[\big]{(a^*a)^{\sfrac{1}{2}}}} \\
	% 	&\StackTextClap{\ref{satz:53}}{\le} \norm*{\pi_{\overline{\varphi}} \enbrace*{a^*a}^{\sfrac{1}{2}}}_{\mathcal{B}(\mathcal{H}_{\overline{\varphi}})} \cdot \Underbracket{\norm*{\overline{\varphi}}^{\sfrac{1}{2}}}{=1} \cdot \overline{\varphi}(a^*a)^{\sfrac{1}{2}}\\
	% 	&= \norm*{\pi_{\overline{\varphi}}(a)}_{\mathcal{B}(\mathcal{H}_{\overline{\varphi}})} \cdot \norm*{a}\le \norm*{\pi_u(a)}_{\mathcal{B}(\mathcal{H}_{u})}\cdot \norm*{a}
	% \end{align}
	% Damit ist $\pi_u$ injektiv also treu.
	Falls $A$ separabel ist, so wählen wir eine abzählbare dichte Teilmenge $\set*{a_0,a_1,\ldots } \subset A^1$ und Zustände $\varphi_0, \varphi_1,\ldots \in S(A)$ mit $\varphi_i(a_i^*a_i) = \norm*{a_i^*a_i}= \norm*{a_i}^2$. 
	Für $a \in A^1$ existiert dann ein $i \in \mathbb{N}$ mit $\varphi_i(a_i^*a_i) \ge \frac{1}{2} \norm*{a}^2$.
	Dann ist 
	\[
		\norm*{\pi_{\varphi_i}(a)} \ge \norm[\big]{\pi_{\varphi_i}(a) \xi_{\varphi_i}}_{\mathcal{H}_{\varphi_i}}\hspace{-1ex}
		= \varphi_i \enbrace*{a^*a}^{\sfrac{1}{2}} \ge \sfrac{1}{\sqrt{2}}\cdot \norm*{a}
	\]
	Wir definieren $\mathcal{H}_{\sep} := \bigoplus_{\mathbb{N}} \mathcal{H}_{\varphi_i}$ und $\pi_{\sep} \colon A \to \mathcal{B}(\mathcal{H}_{\sep})$ durch 
	\(
		\pi_{\sep}(a) \enbrace[\big]{(\eta_i)_{i \in \mathbb{N}}} := \enbrace[\big]{\pi_{\varphi_i}(a)(\eta_i)}_{i \in \mathbb{N}}
	\).
	Mit obiger Ungleichung folgt $\norm*{\pi_{\sep}(a)} \ge \sfrac{1}{\sqrt{2}} \norm*{a}\ge 0$ und damit ist $\pi_{\sep}$ treu, denn nur für $a=0$ kann dann $\pi_{\sep}(a)=0$ gelten.
\end{beweis}

\begin{korollar}[{name=[Matrizen über einer C^*-Algebra]}]
	Sei $A$ eine $C^*$-Algebra und $n\ge 1$ eine natürliche Zahl.
	Dann ist $M_n(A)$ eine $C^*$-Algebra mit Operationen 
	\[
		(a_{ij})_{ij} \cdot (b_{ij})_{ij} = \enbrace*{\sum_{k=1}^{n} a_{ik} b_{kj}}_{ij} \quad \text{ und } \quad (a_{ij})^*_{ij} = (a_{ji}^*)_{ij}
	\]
	Wegen \autoref{korr:212} ist die $C^*$-Norm auf $M_n(A)$ eindeutig bestimmt.
\end{korollar}
\begin{beweis}
	Wir müssen eine geeignete Norm auf $M_n(A)$ definieren.
	Sei dazu $(\mathcal{H}_u,\pi_u)$ die universelle Darstellung von $A$.
	Definiere
	\[
		\begin{tikzcd}[row sep=0pt]
			\pi_u^\ssbrace{n} \colon M_n(A) \rar & M_n \enbrace[\big]{ \mathcal{B} \enbrace*{\mathcal{H}_u}} \rar & \mathcal{B} \enbrace*{\mathcal{H}_u^{\oplus n}} \\
			\mbox{ } \hspace{2em}(a_{ij})_{ij} \rar[mapsto] & \enbrace[\big]{\pi_u(a_{ij})}_{ij} \rar[mapsto] & \pi_u^\ssbrace{n} \enbrace[\big]{(a_{ij})_{ij}}
		\end{tikzcd}
	\]
	wobei analog zur Matrix-Vektor-Multiplikation gilt
	\[
		\pi_u^\ssbrace{n} \enbrace[\big]{(a_{ij})_{ij}} %\enbrace[\big]{(\xi_j)_{j=1,\ldots ,n}}
		\begin{psmallmatrix}
			\xi_1 \\ \vdots \\ \xi_n
		\end{psmallmatrix}
		= \enbrace*{\sum\nolimits_{j} \pi_u(a_{ij}) \xi_j}_{i=1,\ldots ,n}
	\]
	Damit ist $\pi_u^\ssbrace{n}$ ein $^*$-Homomorphismus; $\pi_u^\ssbrace{n}$ ist injektiv, da $\pi_u$ injektiv ist.\marginnote{Hat jemand mal Lust das alles nachzurechnen und mir zu schicken?}
	Nun definiert 
	\[
		\norm[\big]{(a_{ij})_{ij}}_{M_n(A)} := \norm*{\pi_u^\ssbrace{n} \enbrace[\big]{(a_{ij})_{ij}}}_{\mathcal{B} (\mathcal{H}_u^{\oplus n})} 
	\]
	eine $^*$-Algebra Norm, welche die $C^*$-Bedingung erfüllt.
	Weiter gilt für $i,j =1,\ldots n$
	\begin{align}
		\norm*{a_{ij}}_A =  \norm*{\pi_u(a_{ij})}_{\mathcal{B}(\mathcal{H}_u)} \le \norm*{\pi_u^\ssbrace{n} \enbrace[\big]{(a_{kl})_{kl}}}_{\mathcal{B} (\mathcal{H}_u^{\oplus n})}
		= \norm[\big]{(a_{kl})_{kl}}_{M_n(A)} &\le \sum_{k,l} \norm[\big]{\pi_u(a_{kl})}_{\mathcal{B}(\mathcal{H}_u)} \\
		&= \sum_{k,l} \norm*{a_{kl}}_A
	\end{align}
	Mit der Vollständigkeit von $A$ folgt, dass $M_n(A)$ vollständig ist bezüglich $\norm*{\cdot}_{M_n(A)}$.
\end{beweis}

\begin{korollar}[{name=[selbstadjungiert und positiv mittels Zuständen]},label=korr:514]
	Sie $A$ eine $C^*$-Algebra und $a \in A$.
	Dann gilt
	\begin{enumerate}[(i),itemsep=0pt]
		\item $a \in A_{\sa}$ genau dann, wenn $\varphi(a) \in \mathbb{R}$ für jeden Zustand $\varphi \in S(A)$.
		\item $a \in A_+$ genau dann, wenn $\varphi(a) \ge 0$ ist für jeden Zustand $\varphi \in S(A)$.
	\end{enumerate}
\end{korollar}
\begin{beweis}
	\leavevmode
	\begin{enumerate}[(i)]
		\item Die erste Implikation folgt mit \autoref{satz:53} (ii).
		Für die umgekehrte Implikation betrachte das Element $i(a-a^*) \in A_{\sa}$. Dann gilt für $\varphi \in S(A)$
		\[
			\varphi \enbrace[\big]{i \cdot (a- a^*)} = i \enbrace[\big]{\varphi(a) - \varphi(a^*)} = i \cdot  \enbrace[\big]{\varphi(a) - \overline{\varphi(a)}} =0
		\]
		Da $\varphi$ beliebig war folgt mit \autoref{satz:55} $i \cdot (a-a^*)=0$ und damit $a=a^*$.
		\item Hier ist die erste Implikation trivial. 
		Für die umgekehrte Implikation sei $(\mathcal{H}_u,\pi_u)$ die universelle Darstellung von $A$ und $a \in A$, sodass $a$ unter jedem Zustand auf eine positive Zahl abgebildet wird. 
		Für einen Vektor $\xi \in \mathcal{H}_u$ ist $\varphi_\xi \colon A \to \mathbb{C}$, $b \mapsto \skal*{\xi}{\pi_u(b)\xi}$ ein positives Funktional, wie man sich leicht anhand \autoref{satz:54} überlegt.
		Also ist $\skal*{\xi}{\pi_u(a)\xi}=  \varphi_\xi(a)\ge 0$ für $\xi \in \mathcal{H}_u$.
		Sei $\lambda<0$ und $\xi\in \mathcal{H}_u$.
		Dann gilt, da $a$ nach (i) selbstadjungiert ist
		\begin{align}
			\norm[\big]{\enbrace*{\pi_u(a)-\lambda} \xi}^2_{\mathcal{H}_u} &= \skal[\big]{\enbrace*{\pi_u(a)-\lambda}\xi}{ \enbrace*{\pi_u(a)-\lambda}\xi} \\
			&= \norm*{\pi_u(a)\xi}^2 + \abs*{\lambda}^2 \cdot \norm*{\xi}^2 \Underbracket{- 2 \cdot \lambda \cdot \skal*{\xi}{\pi_u(a)\xi}}{\ge 0} \ge \abs*{\lambda}^2 \cdot \norm*{\xi}^2 
		\end{align}
		Damit folgt $\ker \enbrace[\big]{\pi_u(a)-\lambda}= \set*{0}$ und $\enbrace*{\pi_u(a)-\lambda}(\mathcal{H}_u) \subset \mathcal{H}_u$ ist abgeschlossen.
		Also ist 
		\[
			\im \enbrace[\big]{\pi_u(a)-\lambda} = \enbrace[\big]{\ker \enbrace*{\pi_u(a)-\lambda}^*}^\bot \stackrel{a^*=a}{=} \enbrace[\big]{\ker \enbrace*{\pi_u(a)-\lambda}}^\bot = \mathcal{H}_u
		\]
		und folglich ist $\pi_u(a)-\lambda$ surjektiv und somit auch invertierbar.  
		Damit ist $\lambda$ also kein Element des Spektrums von $\pi_u(a)$. 
		Da $\pi_u$ isometrisch ist, ist $\lambda$ auch kein Element von $\sigma(a)$ und es folgt $a \ge 0$ wie gewünscht.\qedhere
	\end{enumerate}
\end{beweis}

\begin{beispiel}[{name=[Toeplitz-Algebra]}]
	Sei $\mathcal{H}= \ell^2(\mathbb{N})$ mit der Standard-Orthonormalbasis $(\eta_i)_{i \in \mathbb{N}}$.
	Definiere den unilateralen Shift\index{unilateraler Shift} $v \in \mathcal{B}(\mathcal{H})$ durch $v(\eta_i) := \eta_{i+1}$.
	Dann ist $v^* \in \mathcal{B}(\mathcal{H})$ gegeben durch 
	\[
		v^*(\eta_i) = \begin{cases}
			\eta_{i-1} &\text{ falls }i\ge 1\\
			0 &\text{ sonst }
		\end{cases}
	\]
	Es gilt $v^*v=\ind_{\mathcal{B}(\mathcal{H})}(= \id)$ und 
	\[
		\ind -vv^* = p_{\eta_0}=e_{00} \in \mathcal{K} = \mathcal{K} \enbrace*{\ell^2(\mathbb{N})}
	\]
	Wir definieren nun die \Index{Toeplitz-Algebra} $\mathcal{T}$ durch
	\[
		\mathcal{T} := C^*(v) \subset \mathcal{B} \enbrace*{\ell^2(\mathbb{N})}
	\]
	Weiter sei $J := \langle \ind - vv^* \rangle \lhd \mathcal{T}$ das von $\ind-vv^*$ in $\mathcal{T}$ erzeugte (abgeschlossene zweiseitige) Ideal.
	Wir zeigen, dass $J=\mathcal{K}$ ist: Offensichtlich ist $J \subset \mathcal{K}$, denn $\ind -vv^* \in \mathcal{K}$.\marginnote{$e_{ij}(\eta_k) =\delta_{jk} \eta_i$}
	Für die Matrixeinheiten $e_{ij} \in \mathcal{K}$ gilt $e_{ij} = v^i(\ind-vv^*)(v^*)^j \in J$ wie man leicht auf Basiselementen nachrechnet.
	Damit ist
	\[
		\mathcal{K}\stackrel[\text{Blatt 14}]{\text{F.A.}}{=\joinrel=\joinrel=\joinrel=}\overline{\Span \set*{e_{ij}}} \subset J
	\]
	Für die Quotientenabbildung $\pi \colon \mathcal{T} \to \sfrac{\mathcal{T}}{J}$ gilt $\sfrac{\mathcal{T}}{J} = C^* \enbrace[\big]{\pi(v)}$ und $\pi(\ind - vv^*)=0$, also ist 
	\[
		\pi(v)^* \pi(v)= \ind_{\sfrac{\mathcal{T}}{J}} =  \pi(v)\pi(v^*)
	\]
	und damit ist $\pi(v) \in \sfrac{\mathcal{T}}{J}$ unitär.
	Man zeigt mit dem Fredholmindex $\sigma \enbrace[\big]{\pi(v)} = \mathbb{T}$ und erhält einen Isomorphismus $\sfrac{\mathcal{T}}{J} \cong C(\mathbb{T})$.
	Wir haben nun eine kurze exakte Sequenz von $C^*$-Algebren
	\[
		\begin{tikzcd}
			0 \rar & \mathcal{K} \rar[hook] & \mathcal{T} \rar[two heads] & C(\mathbb{T}) \rar & 0
		\end{tikzcd}
	\]
	Für jede unitale $C^*$-Algebra $A$ mit einer Isometrie $w$, das heißt $w^*w=\ind$, existiert ein $^*$"=Homomorphismus $\varphi\colon\mathcal{T}\to A$ mit $\varphi(v)=w$, dies besagt der \emph{Satz von Coburn}.
\end{beispiel}

\begin{beispiel}[{name=[Cuntz-Algebren]}]
	Sei $\mathcal{H}= L^2([0,1])$.
	Definiere $v_1,v_2 \in \mathcal{B}(\mathcal{H})$ durch
	\[
		v_1(f)(t) := \begin{cases}
			\sqrt{2} \cdot f(2t) &\text{ falls }t \in [0,\sfrac{1}{2}]\\
			0 &\text{ falls } t \in (\sfrac{1}{2},1]
		\end{cases}
		\quad \text{ und } \quad 
		v_2(f)(t) := \begin{cases}
			\sqrt{2} \cdot f(2t-1) &\text{ falls }t \in [\sfrac{1}{2},1]\\
			0 &\text{ falls } t \in [0,\sfrac{1}{2})
		\end{cases} 
	\]
	Wie man sich leicht überlegt, ist $v_i^* v_i=\ind$ und $v_1v_1^* + v_2v_2^*=\ind$. 
	Die von $v_1$ und $v_2$ erzeugte $C^*$-Algebra  
	\[
		\mathcal{O}_2 := C^*(v_1,v_2) \subset \mathcal{B}(\mathcal{H}) 
	\]
	ist eine \href{https://de.wikipedia.org/wiki/Cuntz-Algebra}{\Index{Cuntz-Algebra}}. Ähnlich lassen sich $\mathcal{O}_n$ und $\mathcal{O}_\infty$ definieren.
	$\mathcal{O}_2$ besitzt keinen Spurzustand und ist einfach (d.h. es gibt keine nichttrivialen abgeschlossenen zweiseitigen Ideale).
	Ist nun $A$ eine beliebige unitale $C^*$-Algebra mit $s_1,s_2 \in A$, sodass $s_1^*s_1 = s_2^*s_2=\ind = s_1s_1^* + s_2s_2^*$.
	Dann existiert ein $^*$-Homomorphismus $\pi \colon \mathcal{O}_2 \to A$ mit $\pi(v_i)=s_i$ für $i=1,2$.
	\begin{figure}[bht]
		\centering{
		\begin{tikzpicture}
			\begin{axis}[
				xtick distance=0.5,
				% ytick distance=0.5,
				minor x tick num=4,
				% minor y tick num=4,
				xticklabel style={/pgf/number format/frac},
				declare function={f(\x)= sin(deg(2*pi*\x)) + 2 * \x + 3;}, %2*(\x)-0.5*(\x)^2
				width=8cm,
			]
				\addplot[
					Firebrick2,
					domain=0:1,
					samples=200,
					very thick,
				]
				{f(x)} node[pos=0.65,below]{$f$};
				\addplot[
					DodgerBlue3,
					domain=0:0.5,
					samples=200,
					very thick,
				]
				{sqrt(2)*f(2*x)} node[pos=0.5,above right]{$v_1(f)$};
				\addplot[
					SeaGreen4,
					domain=0.5:1,
					samples=200,
					very thick,
				]
				{(sqrt(2))*f(2*x-1)} node[pos=0.5,above right]{$v_2(f)$};
			\end{axis}
		\end{tikzpicture}
		\caption{Wirkung von $v_1$ und $v_2$ in der Cuntz-Algebra $\mathcal{O}_2$ über $\mathcal{H}=L^2([0,1])$}
		}
	\end{figure}
\end{beispiel}
% section 5 (end)
\newpage

\section{Von-Neumann-Algebren und der Bikommutantensatz} % (fold)
\label{sec:6}

\begin{erinnerungA}[{name=[starke Operatortopologie]},label=erinnerung:61]
	Die \Index{starke Operatortopologie} $\Tso$ auf $\mathcal{B}(\mathcal{H})$ ist die lokalkonvexe Topologie, welche von den Halbnormen $p_\xi \colon \mathcal{B}(\mathcal{H}) \to \mathbb{R}_+$, $x \mapsto \norm*{x \xi}_\mathcal{H}$ für $\xi \in \mathcal{H}$ erzeugt wird.
	Das heißt $\Tso$ ist die gröbste Topologie, sodass Addition, skalare Multiplikation und alle Halbnormen $p_\xi$ stetig sind.
	\enquote{Typische} offene Mengen sind der Form
	\[
		U_{x,\xi,\varepsilon} := \set[\big]{y \in \mathcal{B}(\mathcal{H}) \given \norm*{(x-y)\xi}<\varepsilon}
	\]
	wobei $x \in \mathcal{B}(\mathcal{H})$, $\xi \in \mathcal{H}$ und $\varepsilon>0$. 
	Die $U_{x,\xi,\varepsilon}$ bilden eine Subbasis für $\mathcal{T}_\so$.
	Für ein Netz $(x_\lambda)_\Lambda \subset \mathcal{B}(\mathcal{H})$ und $x \in \mathcal{B}(\mathcal{H})$ gilt\marginnote{Damit ist $\Tso$ die Topologie der punktweisen Konvergenz}
	\begin{align}
		x_\lambda \grenzw{\so} x &\iff \norm*{(x_\lambda-x)\xi}_{\mathcal{H}} \grenzw{ }0 \quad \text{ für jedes } \xi \in \mathcal{H} \\
		&\iff x_\lambda \xi \grenzw{} x \xi \qquad \text{ für jedes } \xi \in \mathcal{H} 
	\end{align}
	$\Tso$ ist schwächer als $\Top{\norm*{\cdot}}$, das heißt $\enbrace*{\mathcal{B}(\mathcal{H}), \Top{\norm*{\cdot}}} \grenzw{\raisemath{-2pt}{\id}} \enbrace*{\mathcal{B}(\mathcal{H}), \Tso}$ ist stetig.
	Äquivalent dazu: Wenn $A \subset \mathcal{B}(\mathcal{H})$ $\so$-abgeschlossen ist, so ist $A$ auch $\norm*{\cdot}$-abgeschlossen. Alternativ: Wenn $x_\lambda\to x$ in der Normtopologie, dann auch in der starken Operatortopologie.
	
	Die Norm $\norm*{\cdot}$, die Multiplikation und Involution sind nicht stetig bezüglich der starken Operatortopologie:
	Sei dazu $\mathcal{H}=\ell^2(\mathbb{N})$ mit Standardorthonormalbasis $(\eta_k)_\mathbb{N}$ und Matrixeinheiten\index{Matrixeinheiten}\marginnote{$e_{ij} \eta_k= \delta_{jk} \cdot \eta_i$} $e_{ij}$.
	Dann gilt $e_{0,n}\grenzw{\raisemath{-2pt}{\so}} 0$ aber weder $1=\norm*{e_{0,n}}$ noch $e_{0,n}^* = e_{n,0}$ konvergieren gegen $0$ für $n \to \infty$.
	Dass die Multiplikation nicht stetig ist, ist etwas subtiler (Übungsaufgabe 5 von Blatt 9).
\end{erinnerungA}

\begin{satz}[label=satz:62,{name=[monotone Netze selbstadjungierter Operatoren]}]
	Sei $(x_\lambda)_\Lambda \subset \mathcal{B}(\mathcal{H})_\sa$ ein monotones und beschränktes Netz selbstadjungierter Operatoren.
	Dann existiert ein $x \in \mathcal{B}(\mathcal{H})$ mit 
	\[
		x_\lambda \grenzw{\so} x
	\]
	Falls $(x_\lambda)_\Lambda$ ein aufsteigendes Netz von orthogonalen Projektionen ist, so ist $x$ die (orthogonale) Projektion auf $\overline{\bigcup_{\Lambda} x_\lambda \mathcal{H}}$.
	Falls $(x_\lambda)_\Lambda$ ein absteigendes Netz aus Projektionen ist, so ist $x$ die Projektion auf $\bigcap_\Lambda\overline{x_\lambda \mathcal{H}}$.
\end{satz}
\begin{beweis}
	Der erste Teil wurde schon in Funktionalanalysis Satz 7.32 bewiesen. 
	Hier noch einmal kurz die Idee: 
	Mittels Polarisierung und der Monotonie sieht man, dass $\skal*{x_\lambda \xi}{\cdot} \in \mathcal{H}^*\cong \mathcal{H}$ schwach gegen ein $\eta:= x \xi$ konvergiert.
	Also
	\[
		x_\lambda \grenzw{\mathrm{wo}} x
	\]
	da $\skal*{\eta}{x_\lambda \xi} \to \skal*{\eta}{x \xi}$ für $\xi,\eta \in \mathcal{H}$. 
	Mit Beschränktheit folgt dann $x_\lambda \grenzw{\so} x$.
	In dem Fall, dass die $x_\lambda$ Projektionen sind, folgt $x^2=x$ mittels 
	\[
		\skal*{\eta}{x \xi} = \lim_\lambda \skal*{\eta}{x_\lambda \xi} = \lim_\lambda \skal*{x_\lambda \eta}{x_\lambda \xi} = \skal*{x \eta}{x \xi} = \skal*{\eta}{x^* x \xi}
	\]
	Also ist $x=x^*x=x^2$. \emph{Der Rest ist Übungsaufgabe 4 von Blatt 9, siehe \cref{sub:mon_proj}}
\end{beweis}

\begin{bemerkung}[{name=[verschiedene Definitionen von Positivität]},label=bem:63]
	In Funktionalanalysis hatten wir $0\le x \in \mathcal{B}(\mathcal{H})$ definiert als $\skal*{\xi}{x \xi} \ge 0$ für alle $\xi \in \mathcal{H}$ (und $x \in \mathcal{B}(\mathcal{H})_\sa$).
	Dies ist äquivalent zu $0\le x$ in unserem Sinne: 
	Ist $x$ positiv in unserem Sinne, so haben wir $x=y^*y$ und es folgt
	\[
		\skal*{\xi}{x \xi} = \skal*{\xi}{y^*y \xi} = \skal*{y \xi}{y \xi} \ge 0
	\]
	für $\xi \in \mathcal{H}$. 
	Die Gegenrichtung zeigt man wie im Beweis von \autoref{korr:514}.
\end{bemerkung}

\begin{definition}[{name=[von-Neumann-Algebra]}]
	Eine \Index{von-Neumann-Algebra} auf dem Hilbertraum $\mathcal{H}$ ist eine in der starken Operatortopologie abgeschlossene $^*$-Unteralgebra von $\mathcal{B}(\mathcal{H})$.
\end{definition}

\begin{definitionP}[{name=[Kommutante und Bikommutante]},label=def:65]
	Sei $B$ eine Algebra und $S \subset B$ eine Teilmenge.
	\[
		S' := \set[\big]{b \in B \given bs=sb, s \in S}
	\]
	heißt \Index{Kommutante} von $S$ in $B$. 
	$S'' := (S')'$ heißt \Index{Bikommutante}.
	Sowohl $S'$ als auch $S''$ sind Unteralgebren.
	Falls $B$ eine $^*$-Algebra und $S=S^*$ ist, so sind $S'$ und $S''$ auch $^*$-Unteralgebren.
	Falls $A \subset \mathcal{B}(\mathcal{H})$ eine $^*$-Unteralgebra ist, so ist $A' \subset \mathcal{B}(\mathcal{H})$ eine von-Neumann-Algebra.
\end{definitionP}
\begin{beweis}
	Wir zeigen nun, dass $A'$ abgeschlossen ist in der starken Operatortopologie.
	Sei dazu $(x_\lambda) \subset A'$ ein Netz mit $x_\lambda \grenzw{\raisemath{-2pt}{\so}} x$, $a \in A$ beliebig sowie $\xi,\eta \in \mathcal{H}$.
	Dann gilt
	\begin{align}
		\skal*{\eta}{a x \xi} = \skal*{\eta}{ a \lim\nolimits_\lambda x_\lambda \xi} =\lim\nolimits_\lambda \skal*{\eta}{a x_\lambda \xi} = \lim\nolimits_\lambda \skal*{\eta}{x_\lambda a \xi} = \skal*{\eta}{x a \xi}
	\end{align}
	Also ist $a x= x a$.
	Alle weiteren Behauptungen sind trivial.
\end{beweis}

\begin{lemma}[label=lem:66,{name=[{$A$ liegt stark dicht in der Bikommutante}]}]
	Sei $\mathcal{H}$ ein Hilbertraum und $\ind_{\mathcal{B}(\mathcal{H})} \in  A \subset \mathcal{B}(\mathcal{H})$ eine unitale $^*$-Unteralgebra.
	Dann ist $A$ $\so$-dicht in der Bikommutante $A''$.
\end{lemma}
\begin{beweis}
	Sei $x \in A''$. Zu zeigen: Jede $\so$-offene Umgebung $U$ von $x$ mit $x \in U= \bigcap_{i=1}^n U_{x,\xi_i,\varepsilon}$\marginnote{die $U_{x,\xi‚‘,\varepsilon}$ bilden eine Subbasis von $\Tso$} für $\xi_1, \ldots ,\xi_n \in \mathcal{H}$ besitzt nichtleeren Schnitt mit $A$.
	Die Abbildung 
	\mapdef{\pi^\ssbrace{n} \colon \mathcal{B}(\mathcal{H})}{M_n \enbrace*{\mathcal{B}(\mathcal{H})}\cong \mathcal{B} \enbrace*{\mathcal{H}^{\oplus n}}}{y}{\begin{psmallmatrix}
		y & & 0 \\  & \ddots& \\ 0 & & y
	\end{psmallmatrix}}{}
	ist ein $^*$-Homomorphismus und $\pi^\ssbrace{n}(A) \subset \mathcal{B} \enbrace*{\mathcal{H}^{\oplus n}}$ ist $^*$-Unteralgebra von $\mathcal{B}(\mathcal{H}^{\oplus n})$ mit $\ind_{\mathcal{B}(\mathcal{H}^{\oplus n})}\in \pi^\ssbrace{n}(A)$.
	Wir zeigen, dass $\pi^\ssbrace{n}(x) \in \pi^\ssbrace{n}(A)''$ ist: 
	Sei dafür $y \in \pi^\ssbrace{n}(A)'$ und $a \in A$. 
	Dann gilt
	\[
		\pi^\ssbrace{n}(a) \,e_{ii}\, y\, e_{jj} = e_{ii}\, \pi^\ssbrace{n}(a)\, y \,e_{jj} = e_{ii}\, y \,\pi^\ssbrace{n}(a) \,e_{jj} = \Underbracket{e_{ii} \,y \,e_{jj}}{=y_{ij}\in \mathcal{B}(\mathcal{H})} \,\pi^\ssbrace{n}(a)\marginnote{$e_{ij}(\eta_k) =\delta_{jk} \eta_i$}
	\]
	Dann ist $y_{ij}\,a = a \,y_{ij}$ für $i,j \in \set*{1,\ldots ,n}$ und es folgt $y_{ij} \in A'$, also $y_{ij} \,x = x \,y_{ij}$ für $i,j \in \set*{1,\ldots ,n}$.
	Damit folgt dann aber auch $\pi^\ssbrace{n}(x) \, y = y \, \pi^\ssbrace{n}(x)$ und somit die Zwischenbehauptung.
	Sei nun
	\[
		\underline{\xi} := \begin{pmatrix}
			\xi_1 \\ \vdots \\ \xi_n
		\end{pmatrix} \in \mathcal{H}^{\oplus n} \qquad \text{ und } \qquad K:= \overline{\pi^\ssbrace{n}(A) \underline{\xi}} \subset \mathcal{H}^{\oplus n}
	\]
	Dann ist $K$ ein abgeschlossener Unterraum und $\pi^\ssbrace{n}(A)$-invariant.
	Da $\pi^\ssbrace{n}(A)= \pi^\ssbrace{n}(A)^*$ ist, ist $K$ \emph{reduzierend} für jedes $\pi^\ssbrace{n}(a) \in \pi^\ssbrace{n}(A)$, das heißt 
	\begin{itemize}[itemsep=3pt]
		\item $\pi^\ssbrace{n}(a) K \subset K$ und $\pi^\ssbrace{n}(a)^*K \subset K$ bzw.
		\item $\pi^\ssbrace{n}(a) K \subset K$ und $\pi^\ssbrace{n}(a)(K^\bot) \subset K^\bot$ bzw.
		\item $p_K\, \pi^\ssbrace{n}(a)= \pi^\ssbrace{n}(a)\, p_K$ wobei $p_K \colon H^{\oplus n} \to K$ die orthogonale Projektion ist.
	\end{itemize}
	% $\pi^\ssbrace{n}(a) K \subset K$, $\pi^\ssbrace{n}(a)^*K \subset K$ bzw. $\pi^\ssbrace{n}(a) K \subset K$, $\pi^\ssbrace{n}(a)(K^\bot) \subset K^\bot$, bzw. $p_K \pi^\ssbrace{n}(a)= \pi^\ssbrace{n}(a) p_K$ für $a  \in A$, wobei $p_K \colon H^{\oplus n} \to K$ die orthogonale Projektion ist.
	Damit ist $p_K $ in der Kommutante $ \pi^\ssbrace{n}(A)'$ und somit auch $\pi^\ssbrace{n}(x)\, p_K = p_K \,\pi^\ssbrace{n}(x)$. 
	Insgesamt gilt
	\[
		\pi^\ssbrace{n}(x) \underline{\xi}=\pi^\ssbrace{n}(x)\, \Underbracket{p_K\underline{\xi}}{=\underline{\xi}} = p_K \,\pi^\ssbrace{n}(x) \underline{\xi} \in K
	\]
	Also existiert eine Folge $(a_m)_{m \in \mathbb{N}} \subset A$, so dass $\pi^\ssbrace{n}(x) \underline{\xi} = \lim_{m\to \infty} \pi^\ssbrace{n}(a_m) \underline{\xi}$.
	Dieser Ausdruck hat folgende Gestalt
	\[
		\begin{pmatrix}
			x \xi_1 \\ \vdots \\ x \xi_n 
		\end{pmatrix}
		= \lim_{m \to \infty}
		\begin{pmatrix}
			a_m \xi_1 \\ \vdots \\ a_m \xi_n
		\end{pmatrix}
		= \lim_{m \to \infty}
		\begin{pmatrix}
			a_m & & 0 \\
			& \ddots & \\
			0 & & a_m
		\end{pmatrix}
		\begin{pmatrix}
			\xi_1 \\ \vdots \\ \xi_n
		\end{pmatrix}
	\]
	Also gilt $a_m \xi_j \grenzw{m \to \infty} x \xi_j$ für $j=1,\ldots ,n$ und damit $\norm*{a_m \xi_j - x \xi_j} < \varepsilon$ für $m$ groß genug. 
	Damit ist $a_m \in U_{x,\xi_j,\varepsilon}$ für $m$ groß genug und $j=1,\ldots ,n$.
	Dies war zu zeigen.
\end{beweis}

\begin{korollar}[{name={von Neumanns Bikommutantensatz}},label=korr:67]
	Sei $A \subset \mathcal{B}(\mathcal{H})$ eine unitale $^*$-Unteralgebra mit $\ind_{\mathcal{B}(\mathcal{H})} \in A$.
	Dann ist $A$ eine von-Neumann-Algebra genau dann, wenn $A=A''$ ist.
\end{korollar}
\begin{beweis}
	Für die erste Implikation ist $A= \overline{A}^{\so}$, also $A=A''$ nach \autoref{lem:66}
	Für die Rückrichtung beachte, dass $A''$ eine $\so$-abgeschlossene $^*$-Unteralgebra ist nach \autoref{def:65}.
\end{beweis}

\begin{proposition}[{name=[von-Neumann-Algebren haben Einselemente]},label=prop:68]
	Jede von-Neumann-Algebra $\set*{0}\neq A \subset \mathcal{B}(\mathcal{H})$ besitzt ein Einselement.
\end{proposition}
\begin{beweis}
	Sei $(e_\lambda)_\Lambda \subset A$ eine approximative Eins für $A$.
	Dann gilt $e_\lambda \grenzw{\raisemath{-3pt}{\so}}e$ für ein $e \in A_\sa$ nach \autoref{satz:62}.
	Für $\xi \in \mathcal{H}$ und $a \in A$ gilt
	\[
		(ea)(\xi) = e \enbrace*{a \xi} = \enbrace[\big]{\so\text{-}\lim\nolimits_\lambda e_\lambda} (a \xi) = \enbrace[\big]{\so\text{-}\lim\nolimits_\lambda (e_\lambda a)}(\xi) = \enbrace[\big]{\norm{\cdot}\text{-}\lim\nolimits_\lambda (e_\lambda a)}(\xi) = a(\xi)
	\]
	Damit ist $e\,a =a$ für alle $a$ und damit auch $e\,a^*= a^*$. Also ist $a\,e =a$ und $e$ somit eine Eins.
\end{beweis}

\begin{proposition}[{name=[von-Neumann-Algebra als unitale von-Neumann-Algebra auffassen]}]
	Sei $A \subset \mathcal{B}(\mathcal{H})$ eine von-Neumann-Algebra mit Einselement $e$.
	Dann ist $A$ aufgefasst als Unteralgebra $A_e$ von $\mathcal{B}(e \mathcal{H})$ eine unitale von-Neumann-Algebra ($\ind_{A_e}= \ind_{\mathcal{B}(e \mathcal{H})}$) und es gilt $A_e = (A_e)''$.
\end{proposition}
\begin{beweis}
	Es gilt $e=e^2 \in \mathcal{B}(\mathcal{H})_\sa$, das heißt $e$ ist eine orthogonale Projektion.
	Weiter ist $e \mathcal{H} \subset \mathcal{H}$ abgeschlossen, denn $e \xi_\lambda \to \xi \in \mathcal{H}$ impliziert
	\[
		\enbrace*{\ind_{\mathcal{B}(\mathcal{H})} - e} \xi = \enbrace*{\ind_{\mathcal{B}(\mathcal{H})}- e} \enbrace[\big]{\lim\nolimits_\lambda e \xi_\lambda} = \lim\nolimits_\lambda \enbrace*{(\ind_{\mathcal{B}(\mathcal{H})}-e) e \xi_\lambda} = 0
	\]
	Also ist $e\xi=\xi \in e \mathcal{H}$.
	Damit folgt, dass $e \mathcal{H} \subset \mathcal{H}$ ein Unterhilbertraum ist. Wir haben\marginnote{$\kappa$ ist $^*$-linear und kontraktiv, aber nicht multiplikativ}
	\[
		\begin{tikzcd}[row sep=0pt]
			\mathcal{B}(e \mathcal{H}) \rar[hook,"\iota"] & \mathcal{B}(\mathcal{H}) \rar["\kappa"] & \mathcal{B}(e \mathcal{H}) \\
			& x \rar[mapsto] & e \, x|_{e \mathcal{H}}
		\end{tikzcd}
	\]
	Dabei gilt $\kappa \circ \iota = \id_{\mathcal{B}(e\mathcal{H})}$ und $\kappa|_A \colon A \to \kappa(A) =: A_e$ ist ein $^*$-Isomorphismus.
	Wir müssen noch zeigen, dass $A_e \subset \mathcal{B}(e\mathcal{H})$ $\so$-abgeschlossen ist:
	Sei dazu $(x_\lambda)_\Lambda \subset A_e$ ein Netz mit $x_\lambda \grenzw{\raisemath{-3pt}{\so}} x \in \mathcal{B}(e \mathcal{H})$.
	Für $\xi \in \mathcal{H}$ gilt dann
	\[
		\iota(x_\lambda) (\xi) = \iota(x_\lambda) \enbrace[\big]{e \xi + (\ind -e) \xi} = \iota(x_\lambda)(e \xi) \grenzw{} \iota(x)(e \xi) = \iota(x) \enbrace[\big]{e \xi + (\ind -e) \xi} = \iota(x) (\xi)
	\]
	Also gilt $\iota(x_\lambda) \grenzw{\raisemath{-3pt}{\so}} \iota(x)$ und damit ist $\iota(x) \in A$, da $A$ $\so$-abgeschlossen ist.
	Damit folgt weiter $x = \kappa \circ \iota(x) \in  A_e$ und $\ind_{\mathcal{B}(e \mathcal{H})} = \kappa(e) \in A_e$. Nach \autoref{korr:67} gilt außerdem $A_e= (A_e)''$.
\end{beweis}

\begin{definitionP}[label=def:610,{name=[Bildprojektion]}]
	Sei $A \subset \mathcal{B}(\mathcal{H})$ eine von-Neumann-Algebra und $x \in A$.
	Wir definieren die \Index{Bildprojektion} $[x]$ als orthogonale Projektion auf $\overline{\im x} = \overline{x \mathcal{H}} \subset \mathcal{H}$. 
	Es gilt
	\[
		[x] = [xx^*] \in A
	\]
\end{definitionP}
\begin{beweis}
	Mit den Implikationen $x^* \xi=0 \implies xx^* \xi=0 \implies \skal*{\xi}{xx^*\xi}=0 \implies \skal*{x^* \xi}{x^*\xi}=0 \implies x^* \xi=0$ folgt
	\[
		\overline{\im x} = \enbrace*{\ker x^*}^\bot = \enbrace[\big]{\ker (xx^*)}^\bot = \overline{\im xx^*}
	\]
	Damit ist $[x]=[xx^*]$. 
	Die Folge $\enbrace[\big]{\enbrace*{\sfrac{1}{n} \cdot \ind + xx^*}^{-1} xx^*}_{n \in \mathbb{N}} \subset A_+^1$\marginnote{$xx^*$ ist positiv} ist monoton aufsteigend und beschränkt nach Funktionalkalkül, konvergiert also nach \autoref{satz:62} stark gegen ein $e \in A_\sa$.
	Für $\xi \in \mathcal{H}$ gilt dann
	\begin{align}
		e xx^* \xi = \lim_n \enbrace*{\sfrac{1}{n} \cdot \ind + xx^*}^{-1} xx^* xx^* \xi = \lim_n \Underbracket{\enbrace*{\sfrac{1}{n} \cdot \ind + xx^*}^{-1} xx^*}{\text{beschränkt} } \enbrace*{\sfrac{1}{n} \cdot \ind + xx^*} \xi 
		&= \lim_n xx^* \xi \\[-.8em] &= xx^* \xi
	\end{align}
	Also ist $exx^* = xx^*$ und damit auch $xx^*e=(e xx^*)^* = (xx^*)^* = xx^*$, es folgt also $e[xx^*] = [xx^*] = [xx^*]e$.
	Andererseits gilt für $\xi \in \mathcal{H}$
	\begin{align}
		e \xi = \lim_n \enspace\Underbracket{\enbrace*{\sfrac{1}{n} \cdot \ind + xx^*}^{-1} xx^*}{\mathclap{\text{kommutiert da in $C^*(xx^*,\ind)$}}}\enspace \xi \in \overline{\im xx^*}= [xx^*] \mathcal{H}
	\end{align}
	Damit ist $[xx^*]e \xi = e \xi$, also $[xx^*]e=e$, woraus insgesamt $[xx^*]=e \in A$ folgt.
\end{beweis}

\begin{satz}[{name={Polarzerlegung}},label=satz:611]
	Sei $A \subset \mathcal{B}(\mathcal{H})$ eine von-Neumann-Algebra und $x \in A$.
	Dann existiert genau eine partielle Isometrie\index{partielle Isometrie} $u \in A$ mit $u^*u=[\abs*{x}]$ und $x=u \cdot \abs*{x}$, wobei $\abs*{x}=\enbrace*{x^*x}^{\sfrac{1}{2}}$.
\end{satz}
\begin{beweis}
	Setze $u_n := x\enbrace*{\sfrac{1}{n} \cdot \ind + \abs*{x}}^{-1} \subset A^1$.\footnote{Für die Normabschätzung siehe \cref{sub:normabschatzung_im_beweis_der_polarzerlegung}}
	Nach \autoref{def:610} ist $[\abs*{x}]=[x^*x]= [x^*]$ und $[x^*]x^*=x^*$, also folgt $x=x[x^*]=x [\abs*{x}]$.
	Wegen $\abs*{x} = [\abs*{x}]\, \abs*{x}= \abs*{x}\, [\abs*{x}]$ kommutiert $\benbrace*{\abs*{x}}$ mit $\enbrace*{\sfrac{1}{n} \cdot \ind + \abs*{x}}$ und damit auch mit dem Inversen.
	Zusammen mit $x=x \benbrace*{\abs*{x}}$ folgt dann $u_n= u_n[\abs*{x}]$ für $n \in \mathbb{N}$.
	Weiter gilt $u_n \abs*{x} \grenzw{n \to \infty} x$ in Norm, denn 
	\begin{align}
		\norm*{x-u_n \abs*{x}} \le \Underbrace{\norm[\big]{x- u_n \enbrace[\big]{\sfrac{1}{n}\cdot \ind + \abs*{x} }}}{=0} + \Underbrace{\norm[\big]{ u_n \enbrace[\big]{\sfrac{1}{n} \cdot \ind + \abs*{x}} - u_n \abs*{x}}}{= \norm*{u_n \cdot \sfrac{1}{n}}}\le \frac{1}{n}  \grenzw{n \to \infty} 0
	\end{align}
	Wegen $x^*x= \abs*{x}^2$ haben wir für $n,m \in \mathbb{N}$
	\begin{align}
		\enbrace*{u_n^*- u_m^*} (u_n - u_m) = \enbrace*{ \enbrace[\big]{\sfrac{1}{n} \cdot \ind + \abs*{x}}^{-1} - \enbrace[\big]{\sfrac{1}{m} \cdot \ind + \abs*{x}}^{-1}}^2 \,\abs*{x}^2 
	\end{align}
	Für $\xi \in \abs*{x}^2 \mathcal{H}$ sehen wir genau wie in \autoref{def:610}, dass
	\begin{align}
		\enbrace*{u_n^*- u_m^*} (u_n - u_m) \xi \grenzw{n,m \to \infty} 0
	\end{align}
	Wegen der Beschränktheit der $\enbrace*{u_n^*- u_m^*} (u_n - u_m)$ gilt dies auch für $\xi \in \overline{\abs*{x}^2 \mathcal{H}} \stackrel{[\abs*{x}]=[x^*x]}{=\joinrel=\joinrel=\joinrel=} \overline{\abs*{x} \mathcal{H}}$.
	Für $\xi \in \overline{\abs*{x}\mathcal{H}}^{\raisemath{-3pt}{\bot}}$ gilt $\enbrace*{u_n^*- u_m^*} (u_n - u_m)\xi=0$, also insgesamt
	\(
		\enbrace*{u_n^*- u_m^*} (u_n - u_m) \xi \grenzw{\raisemath{-2pt}{n,m \to \infty}} 0
	\)
	für beliebiges $\xi \in \mathcal{H}$.
	Es folgt $\norm*{(u_n-u_m)\xi} \grenzw{\raisemath{-2pt}{n,m \to \infty}} 0$ für $\xi \in \mathcal{H}$ und damit ist $(u_n \xi)_{n \in \mathbb{N}}$ eine Cauchy-Folge in $\mathcal{H}$, also konvergent.
	$(u_n)_\mathbb{N}$ ist also $\so$-konvergent, das heißt es existiert ein $u \in A$ mit $u = \solim_n u_n$.
	Wegen dem oben gezeigten ist
	\[
		u \xi = \lim_n u_n \xi = \lim_n u_n [\abs{x}] \xi =  \enbrace*{\solim u_n} [\abs*{x}] \xi = u [\abs*{x}] \xi
	\]
	für beliebiges $\xi \in \mathcal{H}$.
	Es gilt also $u= u [\abs*{x}]$.
	Außerdem haben wir für $\xi \in \mathcal{H}$
	\[
		\enbrace[\big]{u \abs*{x}- x} \xi = \lim_n \enbrace[\big]{u_n \abs*{x} - x} \xi = 0 
	\]
	Also ist $u \abs*{x}=x$.
	Für $\xi,\eta \in \mathcal{H}$ gilt 
	\[
		\skal[\big]{\abs*{x}\, \eta}{\abs*{x}\, \xi} = \skal[\big]{\eta}{x^*x \xi} = \skal[\big]{\eta}{\abs*{x}\, u^*u \,\abs*{x}\, \xi} = \skal[\big]{\eta}{ \abs*{x}\, [\abs*{x}]\, u^*u \,\abs*{x}\, [\abs*{x}]\, \xi} 
		= \skal[\big]{\abs*{x}\,\eta}{[\abs*{x}]\, u^*u \,[\abs*{x}]\, \abs*{x}\, \xi}
	\]
	Mit Stetigkeit folgt $\skal[\big]{\eta}{[\abs*{x}]\,\xi} = \skal[\big]{\eta}{[\abs*{x}]\, u^*u \,[\abs*{x}]\, \xi}$ für alle $\eta,\xi \in \overline{\abs*{x}\mathcal{H}}$.
	Damit erhalten wir $[\abs*{x}] = [\abs*{x}]\, u^*u \,[\abs*{x}]$.
	Wegen $u=u[\abs*{x}]$ gilt weiter $[\abs*{x}]=u^*u$, insbesondere ist $u$ eine partielle Isometrie.
	Falls $v \in \mathcal{B}(\mathcal{H})$ die Gleichungen $x=v \,\abs*{x}$ und $v^*v= [\abs*{x}]$ erfüllt, so gilt wegen $v \,\abs*{x}=x=u \,\abs*{x}$ auch 
	\[
		v= v\, [\abs*{x}] = u \,[\abs*{x}] = u 
	\]
	wobei man die Gleichung wieder zunächst auf $\overline{\abs*{x} \mathcal{H}}$ nachrechnet.
	$u$ ist also eindeutig bestimmt.
\end{beweis}

\begin{proposition}[label=prop:612,{name=[approximative Einsen in offener Teilmenge des Spektrums]}]
	Sei $A \subset \mathcal{B}(\mathcal{H})$ eine von-Neumann-Algebra und $x \in A$ normal.
	Für jede offene Menge $U \subset \sigma(x)$ existiert eine eindeutig bestimmte Projektion $e \in A \cap \set*{x}'$ mit 
	\[
		e=\solim_\lambda e_\lambda(x)
	\]
	wobei $(e_\lambda)_\Lambda$ (irgend)eine aufsteigende approximative Eins von $C_0(U)$ ist.
	Wir schreiben auch $\chi_U(x)$ für $e$.
	Es gilt $e f(x) = f(x) e = f(x)$ für alle $f \in C_0(U)$.
	Für eine weitere offene Menge $V \subset \sigma(x)$ gilt 
	\[
		\chi_U(x) \cdot \chi_V(x)=\chi_V(x) \cdot \chi_U(x) = \chi_{U \cap V}(x) \qquad\! \text{ sowie } \qquad\!\chi_U(x) + \chi_V(x) - \chi_{U \cap V}(x)= \chi_{U \cup V}(x) 
	\]
	Für $W \subset \sigma(x)$ abgeschlossen schreiben wir auch $\chi_W(x) := \ind_{\mathcal{B}(\mathcal{H})}-\chi_{\sigma(x)\setminus W}(x)$.
\end{proposition}
\begin{beweis}
	Für eine beliebige aufsteigende approximative Eins $(e_\lambda)_\Lambda \subset C_0(U)$ existiert der starke Limes $e := \solim e_\lambda(x)$ nach \autoref{satz:62}.
	Wir im Beweis von \autoref{prop:68} zeigt man, dass für $f \in C_0(U)$ und alle $\xi \in \mathcal{H}$ gilt
	\[
		e f(x) \xi = f(x) \xi \enspace\implies\enspace e f(x) = f(x), f(x) e = f(x)
	\]
	Sei nun $(h_\gamma)_{\Gamma}$ eine weitere aufsteigende approximative Eins für $C_0(U)$ und $h := \solim h_\gamma(x)$.
	Dann gilt
	\[
		e h \xi = \lim_\lambda \enbrace[\Big]{e_\lambda(x) \enbrace[\big]{\lim\nolimits_\gamma h_\gamma(x) \xi}} = \lim_{\lambda} \enbrace[\big]{\lim\nolimits_\gamma e_\lambda h_\gamma(x)(\xi)}
		= \lim_\lambda \enbrace[\big]{e_\lambda(x)(\xi)} = e \xi
	\]
	Es folgt $ e h = e = h e$ und genauso $h e = h =eh$, also $h=e$ und damit die Eindeutigkeit. 
	Die restlichen Aussagen folgen auf ähnliche Art und Weise.
\end{beweis}

\begin{korollar}[label=korr:613,{name=[selbstadjungierte und unitäre Element in von-Neumann-Algebren]}]
	Sei $A \subset \mathcal{B}(\mathcal{H})$ eine von-Neumann-Algebra.
	\begin{enumerate}[(i),itemsep=2pt]
		\item $A$ hat reellen Rang\index{reeller Rang} $0$, das heißt $A_\sa=\overline{\set*{a \in A_\sa \given \sigma(a) \text{ ist endlich}}}$.
		\item Für die unitäre Gruppe\index{unitäre Gruppe} gilt $\mathcal{U}(A)= \overline{\set[\big]{e^{i \pi h} \given h \in A_\sa^1}}$.
	\end{enumerate}
\end{korollar}
\begin{beweis}
	\leavevmode
	\begin{enumerate}[(i)]
		\item Sei $a \in A_\sa$. Ohne Einschränkungen sei $\sigma(a) \subset [0,1]$ (Warum?).
		Für $N \in \mathbb{N}$ definieren wir eine Zerlegung der Eins von $[0,1]$ wie in \cref{fig:613:zerlEins}.
		\begin{figure}[hbt]
			\centering{
			\begin{tikzpicture}[x=7cm,y=2cm]
				\draw[very thick] (-.1,0) -- (0,0)  -- (1,0) -- (1.1,0);
				\draw[very thick] (0,-.075) node[below]{$0$} -- ++(0,.15);
				\draw[very thick] (1,-.075) node[below]{$1$} -- ++(0,.15);
				\foreach \x in {10,20,...,90}{
					\draw[thick] (${\x/100}*(1,0) +(0,1)$) -- (${\x/100}*(1,0) + (0.1,0)$); 
				}
				\foreach \x in {10,20,...,90,100}{
					\draw[thick] (${\x/100}*(1,0) -(0.1,0)$) -- (${\x/100}*(1,0) + (0,1)$) node[above]{$f_{\pgfmathparse{int(\x/10)}\pgfmathresult}$}; 
				}
				\draw[very thick,Firebrick3]  (0,1) node[above]{$f_0$} -- (0.1,0) -- (1,0);
				\draw[ultra thick,DodgerBlue3,dashed]  (0,0)  -- (0.3,0) -- (0.4,1) node[above,fill=white]{$f_4$} -- (0.5,0) -- (1,0);
			\end{tikzpicture}
			\caption{Zerlegung der Eins auf $[0,1]$ für gegebenes $N \in N$, hier $N=10$}\label{fig:613:zerlEins}
			}
		\end{figure}
		Nach Aufgabe 2 von Blatt 4 gilt für eine Projektion $p$ stets $\sigma(p) \subset \set*{0,1}$.
		Insbesondere gilt dies für die eben definieren $\chi_U$ für $U \subset \sigma(a)$ offen.
		Es gilt dann
		\begin{align}
			\SwapAboveDisplaySkip
			a=\id_{[0,1]}(a)= \sum_{i=1}^N \frac{i}{N}\cdot  f_i(a) = \sum_{i=1}^{N} \frac{1}{N} \sum_{j=i}^{N} f_j(a) \ge \Overbracket{\sum_{i=1}^{N} \frac{1}{N} \cdot \chi_{(\sfrac{i}{N},1]}(a)}{\text{hat endliches Spektrum}} &\ge \sum_{i=1}^{N} \frac{1}{N} \sum_{j=i+1}^{N} f_j(a) \\
			&\ge \id_{[0,1]}(a) - \sfrac{1}{N} \cdot \ind 
		\end{align}
		Also lässt sich $a$ durch Elemente mit endlichem Spektrum beliebig approximieren.
		\item Die Abbildung $(-1,1] \to \mathbb{T}$, $t \mapsto e^{i \pi t}$ ist bijektiv und stetig; die Umkehrfunktion $f$ ist lediglich stetig auf $\mathbb{T} \setminus \set*{-1}$.
		Sei nun $u \in A$ unitär. 
		Dann ist $\sigma(u) \subset \mathbb{T}$ nach \autoref{prop:210} (ii).
		Falls $-1 \notin \sigma(u) \subset \mathbb{T}$, so ist $h=f(u) \in A_\sa^1$ und $u=e^{i \pi h}$.
		
		Falls $-1$ im Spektrum von $u$ enthalten ist, so betrachten wir eine offene Umgebung $V \subset \mathbb{T}$ von $-1$.
		Dann ist $p:= \chi_{V\cap \sigma(u)}(u)$ eine Projektion in $A \cap \set*{u}'$ nach \autoref{prop:612}.
		Für diese ist dann $u(\ind-p) \in B := (\ind-p)A(\ind-p)$ unitär.
		Es gilt $-1 \notin \sigma_B \enbrace*{u(1-p)}$, denn für $f \in C_0(V)$ gilt $f \enbrace[\big]{u(\ind-p)}= f(u)(\ind-p)=0$\marginnote{$f$ mit Polynomen approximieren} und damit folgt $\sigma_B \enbrace[\big]{u(\ind-p)} \subset \mathbb{T}\setminus V$.
		Wir im ersten Teil gilt dann
		\[
			u(\ind-p)= e^{i \pi h_0} \in B
		\]
		für ein $h_0 \in B_\sa^1$. 
		Ähnlich wie in (i) sieht man, dass 
		\(
			\norm*{\enbrace[\big]{u(\ind-p) -p} -u} \le \diam V
		\).
		Weiter gilt, da $h_0$ und $p$ orthogonal zueinander sind
		\[
			u(\ind-p) - p = e^{i \pi(h_0+p)} 
		\]
		\todo{bei dem $e$-Term kommt wieder $u(\ind-p)$ raus meine ich \ldots}
		Wir können nun $V$ beliebig klein wählen und erhalten somit die Behauptung.\qedhere
	\end{enumerate}
\end{beweis}

\begin{bemerkung}[{name=[Folgerungen aus \autoref{korr:613}]},label=bem:614]
	\leavevmode
	\begin{enumerate}[a),itemsep=0pt]
		\item Aus dem Beweis von (i) folgt insbesondere, dass $A= \overline{\Span \set*{p \in A \given p\text{ ist Projektion}}}$.
		\item Die Gleichung in (i) gilt entsprechend auch für normale Elemente.
		\item Mit dem Borel-Funktionalkalkül, \autoref{satz:615}, folgt, dass die Gleichung in (ii) auch ohne den Normabschluss gilt.
	\end{enumerate}
\end{bemerkung}

\begin{satz}[label=satz:615,name={Borel-Funktionalkalkül}]
	Sei $A \subset \mathcal{B}(\mathcal{H})$ eine von-Neumann-Algebra und $a \in A$ normal.
	Dann existiert ein $^*$-Homomorphismus\marginnote{beschränkte Borel-Funktionen bilden $C^*$-Algebra mit $\norm*{\cdot}_\infty$; $W^*(x)$ erzeugte von-Neumann-Algebra}
	\mapdef{\mathbb{B}_b \enbrace[\big]{\sigma(x)}}{W^*(x)\subset A}{f}{f(x)}{} 
	welcher den Gelfandhomomorphismus $C(\sigma(x)) \to C^*(x)$ fortsetzt.
	Für eine aufsteigende beschränkte Folge $(f_n)_\mathbb{N} \subset \mathbb{B}_b(\sigma(x))_\sa$ mit punktweisem Limes $f=\lim_n f_n$ gilt
	\[
		f(x) =\solim f_n(x)
	\]
\end{satz}
\begin{beweis}[name={Skizze}]
	Für $\xi,\eta \in \mathcal{H}$ können wir mit Hilfe von
	\(
		\mu_{\xi,\eta}(U) := \skal[\big]{\xi}{\chi_U(x)(\eta)}
	\)
	für $U \subset \sigma(a)$ offen bzw. abgeschlossen ein Integral auf Treppenfunktionen definieren.
	Hieraus erhält man ein Integral $\Phi_{\xi,\eta}$ auf $\mathbb{B}_b(\sigma(x))$:
	\begin{itemize}[itemsep=0pt]
		\item entweder mit $L^1$-Approximation wie in Analysis \RM{3}.
		\item oder als Regelintegral zunächst für stetige Funktionen, also $C(\sigma(x))$, dann für $\mathbb{B}_b(\sigma(x))$ mit dem Rieszschen Darstellungssatz.
	\end{itemize}
	Damit hat man nun $\Phi_{\xi,\eta} \colon \mathbb{B}_b(\sigma(x)) \to \mathbb{C}$ linear und stetig.
	Für $f \in \mathbb{B}_b(\sigma(x))$ und $\eta \in \mathcal{H}$ definiert $\xi \mapsto \Phi_{\xi,\eta}(f) \in \mathbb{C}$ ein konjugiert lineares, stetiges Funktional $\varphi_{f,\eta}$ auf $\mathcal{H}$.
	Mittels $\skal*{\xi}{\cdot }$ lässt sich dann schließen, dass man $\varphi_{f,\eta}$ als Element von $\mathcal{H}$ auffassen kann, d.h. $\varphi_{f,\eta}(\xi) = \skal*{\xi}{\varphi_{f,\eta}}$.
	Die Abbildung $\eta \mapsto \varphi_{f,\eta}$ ist linear und stetig, definiert also einen Operator $f(x) \in \mathcal{B}(\mathcal{H})$.
	Auf Vektoren in $\mathcal{H}$ rechnet man nach, dass $f \mapsto f(x)$ ein isometrischer $^*$-Homomorphismus $\mathbb{B}_b(\sigma(x)) \to \mathcal{B}{\mathcal{H}}$ ist. 
	Weiter gilt $f(x) \in A''= \overline{A}^{\so}$ nach dem Bikommutantensatz.
	Außerdem gilt $\id_{\sigma(x)} \mapsto x$ und daher setzt das Borel-Funktionalkalkül den Gelfandhomomorphismus fort.
\end{beweis}

\begin{erinnerungA}[{name=[schwache Operatortopologie]},label=erinnerung:616]
	Die \Index{schwache Operatortopologie} $\TWeakOp$ auf $\mathcal{B}(\mathcal{H})$ ist die lokalkonvexe Topologie, welche von den Halbnormen
	\[
		p_{\eta,\xi}(x) := \abs[\big]{\skal*{\eta}{x \xi}} 
	\]
	für $\eta,\xi \in \mathcal{H}$ erzeugt wird.
	Die Mengen $U_{x,\eta,\xi,\varepsilon} := \set[\big]{y \in \mathcal{B}(\mathcal{H}) \given \abs*{\skal*{\eta}{(x-y)\xi}} < \varepsilon}$ bilden eine Subbasis von $\TWeakOp$.
	Die Abbildungen
	\[
		\begin{tikzcd}
			\enbrace[\big]{\mathcal{B}(\mathcal{H}),\Top{\norm*{\cdot}}} \rar["\id"] & \enbrace[\big]{\mathcal{B}(\mathcal{H}),\Top{\so}} \rar["\id"] & \enbrace[\big]{\mathcal{B}(\mathcal{H}), \TWeakOp}
		\end{tikzcd}
	\]
	sind stetig, denn $\norm*{(x-y)\xi} \le \norm*{x-y} \cdot \norm*{\xi}$ und $\abs*{ \skal*{\eta}{(x-y)\xi}} \le \norm*{\eta} \cdot \norm*{(x-y) \xi}$.
	Im Gegensatz zur starken Operatortopologie (vergleiche \autoref{erinnerung:61}) ist die Involution $\wo$-stetig:
	\begin{align}
		x_\lambda \grenzw{\wo} x &\iff \abs[\big]{\skal*{\eta}{(x_\lambda-x)\xi}} \grenzw{} 0 \text{ für } \eta,\xi \in \mathcal{H} \\
		&\iff \abs[\big]{\skal*{(x_\lambda-x)^* \eta}{\xi}}   \grenzw{} 0 \text{ für } \eta,\xi \in \mathcal{H} \\
		&\iff \abs[\big]{\skal*{\xi}{(x_\lambda^* -x^*) \eta}} \grenzw{} 0 \text{ für } \eta,\xi \in \mathcal{H} \\
		&\iff x_\lambda^* \grenzw{\wo} x^* 
	\end{align}
	Die Multiplikation ist allerdings auch nicht $\wo$-stetig: Sei dazu $S$ der Shift. Dann gilt einerseits $S^n,{S^*}^n \grenzw{\wo} 0$ aber andererseits ${S^*}^n S^n=\ind$ für alle $n \in \mathbb{N}$.
\end{erinnerungA}

\begin{satz}[{name=[Einheitskugel der beschränkten Operatoren ist schwach kompakt]}]
	Die Einheitskugel $\mathcal{B}(\mathcal{H})^1$ ist $\wo$-kompakt.
\end{satz}
\begin{beweis}[name={Skizze}]
	% $\id \colon \enbrace*{\mathcal{B}(\mathcal{H}),\TWeakStar} \to \enbrace*{\mathcal{B}(\mathcal{H}),\TWeakOp}$ ist stetig.
	% Nehme $\mathcal{H}=\ell^2(\mathbb{N})$ an und setze $X= \overline{\Span\set*{e_{ij} \given i,i=1,2,\ldots }}^{\norm*{\cdot}_1}$.
	% $x \mapsto \Tr(\abs*{x})$ ist auf $\Span (e_{ij})$ definiert und ist Norm.
	% $X$ ist ein abgeschlossenes Ideal in $\mathcal{B}(\mathcal{H})$.
	% $x \mapsto  \enbrace*{y \mapsto \Tr(xy)}$ liefert eine Abbildung $\mathcal{B}(\mathcal{H}) \to X^*$ (tatsächlich ein Isomorphismus).
	% Sei $\TWeakStar$ die zurückgezogene Topologie auf $\mathcal{B}(\mathcal{H})$.
	% Man überprüft: $x_\lambda \grenzw{\w^*}x \implies x_\lambda\grenzw{\wo}x$ und damit ist
	% \[
	% 	\id \colon \enbrace*{\mathcal{B}(\mathcal{H}), \TWeakStar} \implies \enbrace*{\mathcal{B}(\mathcal{H}),\TWeakOp}
	% \]
	% stetig.
	% Aber $\mathcal{B}(\mathcal{H})^1$ ist $\w^*$ kompakt nach Banach-Alaoglu.
	% Damit ist $\mathcal{B}(\mathcal{H})^1$ $\wo$-kompakt und die Abbildung ist ein Homöomorphismus.
	%
	% \emph{Variante, wie sie in den Notizen steht:}
	Fixiere eine Orthonormalbasis $(\xi_\gamma)_\Gamma$ und definiere die \href{https://de.wikipedia.org/wiki/Spurklasseoperator}{\Index{Spurklasse-Operatoren}}
	\[
		\mathcal{L}^1(\mathcal{H}) := \set[\big]{x \in \mathcal{B}(\mathcal{H}) \given \norm*{x}_1 := {\textstyle\sum\nolimits_\Gamma} \skal*{\xi_\gamma}{\abs*{x} \xi_\gamma } < \infty} 
	\]
	Man zeigt, dass $\norm*{\cdot}_1$ eine Norm ist und dass folgendes gilt:
	\[
		\mathcal{L}^1(\mathcal{H}) = \overline{\set[\big]{x \in \mathcal{B}(\mathcal{H}) \given x \text{ hat endlichen Rang} }}^{\norm{\cdot }_1}
	\]
	Für $y \in \mathcal{B}(\mathcal{H})$ definiert man nun $\varphi_y \in \mathcal{L}^1(\mathcal{H})^*$ durch $x \mapsto \varphi_y(x) := \sum_\Gamma \skal*{\xi_\gamma}{yx \xi_\gamma}$.
	Man zeigt nun, dass $\varphi \colon \mathcal{B}(\mathcal{H}) \to \mathcal{L}^1(\mathcal{H})^*$, $y \mapsto \varphi_y$ einen isometrischen Isomorphismus definiert. 
	Sei $\TWeakStar$ die zurückgezogene Topologie auf $\mathcal{B}(\mathcal{H})$.
	Mit Banach-Alaoglu folgt, dass $\mathcal{B}(\mathcal{H})^1$ $\w^*$-kompakt und hausdorffsch ist.
	Die Abbildung $\id \colon \enbrace*{\mathcal{B}(\mathcal{H})^1, \TWeakStar} \to \enbrace*{\mathcal{B}(\mathcal{H})^1,\TWeakOp}$ ist stetig, was man wie folgt verifiziert: 
	\begin{align}
		y_\lambda \grenzw{\w^*} y \text{ in } \mathcal{B}(\mathcal{H})^1 &\implies \sum_{\gamma} \skal*{\xi_\gamma}{(y_\lambda-y)x \xi_\gamma} \to 0 \text{ für jedes } x \in \mathcal{L}^1(\mathcal{H}) \\[-.9em]
		&\StackTextClap{Übung}{\implies} \skal*{\xi}{(y_\lambda-y) \eta} \to 0 \text{ für alle } \xi,\eta \in \mathcal{H} \\
		&\implies y_\lambda \grenzw{\wo} y 
	\end{align}
	Damit ist $\mathcal{B}(\mathcal{H})^1$ auch $\wo$-kompakt.
\end{beweis}

\begin{satz}[{name=[schwacher Abschluss entspricht der Bikommutante]}]
	Sei $\ind_{\mathcal{B}(\mathcal{H})} \in A \subset \mathcal{B}(\mathcal{H})$ eine $^*$-Algebra. Dann gilt
	\begin{enumerate}[(i),itemsep=0pt]
		\item $\overline{A}^\wo = A''$
		\item $A$ ist eine von-Neumann-Algebra genau dann, wenn $A= \overline{A}^\wo$.
	\end{enumerate}
\end{satz}
\begin{beweis}
	Zur ersten Aussage: Es gilt $A'' \StackText{\ref{lem:66}}{=}\overline{A}^\so \subset \overline{A}^\wo$. 
	Sei nun $(x_\lambda)_\Lambda \subset A$ ein Netz mit $x_\lambda \grenzw{\wo} x$.
	Falls $y \in A'$, so gilt für $\eta,\xi \in \mathcal{H}$
	\[
		\skal[\big]{\eta}{(x y)\xi} = \lim_\lambda \skal[\big]{\eta}{x_\lambda y \xi} =\lim_\lambda \skal[\big]{\eta}{y x_\lambda \xi} = \lim_\lambda \skal[\big]{y^* \eta}{ x_\lambda \xi} = \skal[\big]{y^* \eta}{x \xi} = \skal[\big]{\eta}{(yx) \xi}
	\]
	Es folgt $xy=yx$ und somit ist $x \in A''$, also $\overline{A}^\wo \subseteq A''$.
	Die zweite Aussage folgt offensichtlich direkt aus (i) und dem Bikommutantensatz.
\end{beweis}

Wir wollen uns nun noch ein bisschen mehr mit dem Zusammenspiel der verschiedenen Topologien beschäftigen. 
Für ein lineares Funktional $\varphi$ auf $\mathcal{B}(\mathcal{H})$ gilt: 
\[
	\begin{tikzcd}[row sep=1.6em]
		\varphi \text{ $\wo$-stetig} \rar[Rightarrow] & \varphi \text{ $\so$-stetig}  \rar[Rightarrow] & \varphi \text{ $\norm*{\cdot}$-stetig} \\[-1.5em]
		x_\lambda \grenzw{\wo}x \dar[Rightarrow]& x_\lambda\grenzw{\so}x \lar[Rightarrow]\dar[Rightarrow] & x_\lambda\grenzw{\norm*{\cdot}} x \lar[Rightarrow]\dar[Rightarrow] \\
		\varphi(x_\lambda) \to \varphi(x) \rar[Leftrightarrow] &\varphi(x_\lambda) \to \varphi(x) \rar[Leftrightarrow] & \varphi(x_\lambda)\to \varphi(x) 
	\end{tikzcd}
\]
Die Umkehrung der zweiten Implikation in der obersten Zeile gilt im Allgemeinen nicht. 
Wie wir jetzt sehen werden ist die erste aber tatsächlich eine Äquivalenz:

\begin{satz}[label=satz:619,{name=[für Funktionale ist schwach stetig äquivalent zu stark stetig]}]
	Für ein lineares Funktional $\varphi$ auf $\mathcal{B}(\mathcal{H})$ sind äquivalent:
	\begin{enumerate}[(i),itemsep=0pt]
		\item $\varphi$ ist $\wo$-stetig.
		\item $\varphi$ ist $\so$-stetig.
		\item $\varphi(x)=\sum_{i=1}^{n} \skal*{\eta_i}{x \xi_i} $ für $\eta_1,\xi_1, \ldots , \eta_n, \xi_n \in \mathcal{H}$.
	\end{enumerate}
\end{satz}
\begin{beweis}
	Die Implikation (i) $\Rightarrow$ (ii) ist klar, da $\TWeakOp \subset \Tso$ ist.
	Die Implikation (iii) $\Rightarrow$ (i) ist auch klar nach Definition der schwachen Operatortopologie und wir müssen nur (ii) $\Rightarrow$ (iii) zeigen:
	Sei also $\varphi$ $\so$-stetig. Die Menge $S_\varphi := \set*{x \in \mathcal{B}(\mathcal{H}) \given \abs*{\varphi(x)}<1} \subset \mathcal{B}(\mathcal{H})$ ist dann eine $\so$-offene Umgebung der Null und es existieren $\xi_1, \ldots ,\xi_n \in \mathcal{H}$ und ein $\alpha>0$, sodass für die entsprechenden Elemente $U_{0,\xi_i,\alpha}$ der Subbasis\footnote{Es gilt $U_{x,\xi_i,\alpha}=\set[\big]{y \in \mathcal{B}(\mathcal{H}) \given \norm*{(x-y)\xi_i}< \alpha}$, vergleiche \autoref{erinnerung:61}} von $\Tso$ gilt
	\[
		\bigcap_{i=1}^n U_{0,\xi_i,\alpha}  \subset S_\varphi
	\]
	Wir bezeichnen die dazugehörigen Halbnormen\marginnote{Details hinzugefügt in Anlehnung an \cite[S.~267~f]{Murphy}} mit $p_{i}(x) := \norm*{x \xi_i}$ und setzen $p(x) := \max_{i} p_i(x)$.
	Für $x \in \mathcal{B}(\mathcal{H})$ beliebig gilt dann: $p(x)<\alpha \implies \abs*{\varphi(x)}<1$.
	Wir setzen $M := \sfrac{2}{\alpha}$ und erhalten, falls $p(x)>0$ ist,
	\(
		p \enbrace[\big]{\frac{x}{M \cdot p(x)}} = \frac{\alpha}{2} < \alpha 
	\).
	Mit obiger Implikation folgt $\abs[\big]{\varphi \enbrace[\big]{\frac{x}{M \cdot p(x)}}} < 1$ und wir erhalten insgesamt für alle $x \in \mathcal{B}(\mathcal{H})$
	\begin{equation}
		\abs*{\varphi(x)} \le M \cdot p(x) = M \cdot \max_i \norm*{x \xi_i}
	\end{equation}
	Nach entsprechender Skalierung der $\xi_i$ können wir $M=1$ annehmen und erhalten dann
	\begin{equation}
		\abs*{\varphi(x)} \le  \max_i \norm*{x \xi_i} \le \enbrace*{\sum_{i=1}^{n} \norm*{x \xi_i}^2}^{\mkern-5mu\raisemath{-2.2pt}{\sfrac{1}{2}}} \label{eq:619}\tag{\dag}
	\end{equation}
	Setze nun 
	\[
		\mathcal{H}_0 := \set*{\begin{psmallmatrix}
			x \xi_1 \\ \raisemath{2pt}{\vdots} \\ x \xi_n
		\end{psmallmatrix} \given x \in \mathcal{B}(\mathcal{H})} \subset \mathcal{H}^{\oplus n}
	\]
	und definiere $\psi \colon \mathcal{H}_0 \to \mathbb{C}$ durch $\enbrace[\big]{x \xi_1, \ldots , x \xi_n} \mapsto \varphi(x)$. 
	Man sieht nun, dass $\psi$ wohldefiniert, linear und stetig ist  mit $\norm*{\psi}\le 1$ (unter Benutzung von \eqref{eq:619}).
	$\psi$ setzt sich also fort zu $\overline{\psi} \colon \overline{\mathcal{H}_0} \to \mathbb{C}$.
	Mit Riesz folgt die Existenz eines eindeutigen $\eta= \enbrace*{\eta_1, \ldots ,\eta_n} \in \overline{\mathcal{H}_0}$ mit $\overline{\psi}(x) = \skal*{\eta}{x}$, also gilt
	\[
		\varphi(x) = \overline{\psi} \begin{pmatrix}
			x \xi_1 \\ \vdots \\ x \xi_n
		\end{pmatrix} = \sum_{i=1}^{n} \skal*{\eta_i}{x \xi_i}\qedhere
	\]
\end{beweis}

\begin{proposition}[{name=[konvexe Menge stark g.d.w schwach abgeschlossen]},label=prop:620]
	Sei $C \subset \mathcal{B}(\mathcal{H})$ konvex.
	Dann ist $C$ $\so$-abgeschlossen genau dann, wenn $C$ $\wo$-abgeschlossen ist.
\end{proposition}
\begin{beweis}[name={Idee}]
	Wenn $C$ $\wo$-abgeschlossen ist, dann ist $C$ auch $\so$-abgeschlossen, da die schwache Topologie in der starken enthalten ist.
	Für die andere Implikation sei $x \in \overline{C}^{\raisemath{-1.5pt}{\wo}}$. Zu zeigen ist nun, dass $x$ auch in $\overline{C}^{\raisemath{-1.5pt}{\so}}$ enthalten ist. 
	Es existiert ein Netz $(x_\lambda)_\Lambda \subset C$, dass schwach gegen $x$ konvergiert.
	Für jedes $\wo$-stetige Funktional $\varphi$ auf $\mathcal{B}(\mathcal{H})$ gilt $\varphi(x_\lambda) \to \varphi(x)$.
	Mit \autoref{satz:619} folgt, dass für jedes $\so$-stetige Funktional $\varphi$ auf $\mathcal{B}(\mathcal{H})$ auch $\varphi(x_\lambda) \to \varphi(x)$ gilt.
	Falls $x \notin \overline{C}^{\raisemath{-1.5pt}{\so}}$, so folgt mit dem Lemma von Zorn (bzw. mit Hahn-Banach für lokalkonvexe Räume), dass $\overline{\varphi} \in \mathcal{B}(\mathcal{H})^*$ $\so$-stetig existiert mit $\dist(\overline{\varphi}(x),\overline{\varphi}(C))>0$.
	Dies ist ein Widerspruch.
\end{beweis}

\begin{proposition}[label=prop:621,{name=[Involution stark stetig auf normalen Elementen]}]
	Die Involution ist $\so$-stetig auf $\set*{x \in \mathcal{B}(\mathcal{H}) \given x \text{ normal} }$.
\end{proposition}
\begin{beweis}
	Für $x,y \in \mathcal{B}(\mathcal{H})$ normal und $\xi \in \mathcal{H}$ gilt 
	\begin{align}
		\norm*{(x^*-y^*) \xi}^2 &= \norm*{x \xi}^2 - \norm*{y \xi}^2 + 2 \cdot \skal*{\xi}{yy^* \xi} - \skal*{\xi}{xy^* \xi} - \skal*{\xi}{yx^* \xi} \\
		&= \norm*{x \xi}^2 - \norm*{y \xi}^2 + \skal[\big]{\xi}{(y-x)y^* \xi} + \Underbracket{\skal[\big]{\xi}{y(y^*-x^*)\xi}}{=\skal*{(y-x)y^* \xi}{\xi}} \\[-1em]
		&\StackTextClap{C.S.}{\le} \norm*{x \xi}^2 - \norm*{y \xi}^2 + 2 \cdot \norm*{(y-x)y^* \xi} \cdot \norm*{\xi}
	\end{align}
	Falls nun $x_\lambda \grenzw{\so} y$ in $\mathcal{B}(\mathcal{H})$, so gilt 
	\[
		\norm*{(x_\lambda^* -y^*)\xi} \le \Underbracket{\norm*{x_\lambda \xi}^2 - \norm*{y \xi}^2}{\to 0} + \Underbracket{2 \cdot \norm*{(y-x_\lambda)y^* \xi} \cdot \norm*{\xi}}{\to 0} \grenzw{} 0 \qedhere
	\]
\end{beweis}

\begin{definition}[label=def:622,{name=[{stark stetige Funktion}]}]
	Eine stetige Funktion $f \colon \mathbb{R} \to \mathbb{C}$ heißt \Index{stark stetig}, falls für jeden Hilbertraum $\mathcal{H}$ und jedes Netz $(x_\lambda)_\Lambda \subset \mathcal{B}(\mathcal{H})_\sa$ mit $x_\lambda \grenzw{\so} x \in \mathcal{B}(\mathcal{H})_\sa$ gilt $f(x_\lambda) \grenzw{\so}f(x)$.
\end{definition}

\begin{proposition}[label=prop:623,{name=[jede beschränkte stetige Funktion ist stark stetig]}]
	Jedes $h \in C_b(\mathbb{R})$ ist stark stetig.
\end{proposition}
\begin{beweis}
	Sei $S := \set*{f \colon \mathbb{R}\to \mathbb{C} \given f \text{ stark stetig}}$ und $S_b := S \cap C_b(\mathbb{R})$.
	Dann sind $S$ und $S_b$ Vektorräume. 
	Wegen $\norm*{ab \xi - a_0 b_0 \xi} \le \norm*{a} \cdot \norm*{(b-b_0)\xi} + \norm*{(a-a_0)b_0\xi}$ haben wir für $f \in S_b$, $g \in S$ und $x_\lambda \grenzw{\so}x$ in $\mathcal{B}(\mathcal{H})_\sa$
	\[
		\norm*{f(x_\lambda) g(x_\lambda) \xi - f(x) g(x)\xi} \le \norm*{f}_\infty \cdot \norm*{\enbrace[\big]{g(x_\lambda)- g(x)}\xi} + \norm*{ \enbrace[\big]{f(x_\lambda)- f(x)} g(x) \xi}
		\grenzw{} 0
	\]
	Damit ist $f \cdot g \in S$, das heißt $S_b \cdot S \subset S$ und $S_b \cdot S_b \subset S_b$.
	Weiter ist $S_b$ selbstadjungiert, also $S_b^*=S_b$, denn  $f(x)$ für $x$ selbstadjungiert  ist normal und nach \autoref{prop:621} ist $f(x) \mapsto f(x)^*$ $\so$-stetig. 
	Außerdem ist $S_b$ abgeschlossen unter $\norm*{\cdot}_\infty$.
	Insbesondere ist $S_b \subset C_b(\mathbb{R})$ eine $C^*$-Unteralgebra.
	Ebenso $S_b \cap C_0(\mathbb{R}) =: S_0 \subset C_0(\mathbb{R})$.
	
	Die Funktionen $t \mapsto f(t):= \frac{1}{1+t^2}$ und $t \mapsto g(t) := \frac{t}{1+t^2}$ sind in $C_0(\mathbb{R})$ mit $\norm*{f}_\infty, \norm*{g}_\infty\le 1$.
	Behauptung: $f$ und $g$ sind in $S_0$ enthalten.
	Für $x,y \in \mathcal{B}(\mathcal{H})_\sa$ gilt 
	\[
		g(x)-g(y)= \ldots = (\ind+x^2)^{-1} \enbrace[\big]{x-y+x(y-x)y} (\ind+y^2)^{-1} 
	\]
	Für $\xi \in \mathcal{H}$ gilt dann weiter, da $\norm*{(\ind+x^2)^{-1}} \le 1$ ist
	\[
		\norm*{\enbrace[\big]{g(x)-g(y)}\xi} \le \norm*{(x-y)(\ind+y^2)^{-1} \xi} + \norm*{x(y-x) y(\ind+y^2)^{-1} \xi}
	\]
	Damit ist $g$ stark stetig und somit in $S_0$.
	Weiter ist $(t \mapsto t) \in S$, also $\enbrace[\big]{t \mapsto t \cdot g(t)} \in S$ und $\enbrace[\big]{t \mapsto 1-t \cdot g(t)}=f \in S_0$. 
	Damit ist die Behauptung gezeigt.
	
	Nach Stone-Weierstraß ist $^*\text{-}\Alg(f,g) \subset C_0(\mathbb{R})$ dicht, denn diese Algebra trennt Punkte von $\mathbb{R}$ und enthält mit $f$ eine strikt positive Funktion.
	Damit folgt $S_0=C_0(\mathbb{R})$.
	Für $h \in C_b(\mathbb{R})$ gilt $hf, hg \in S_0 \subset S_b$, also ist auch $\enbrace*{t \mapsto t \cdot h(t) \cdot g(t)} \in S$ und
	\[
		\enbrace[\big]{t \longmapsto h(t) \cdot f(t) + t \cdot h(t) \cdot g(t)=h(t)} \in S_b
	\]
	Also gilt $S_b \subset C_b(\mathbb{R})$ und damit ist die Behauptung gezeigt.
\end{beweis}

Zum krönenden Abschluss dieses Kapitels wollen wir noch den Dichtesatz von Kaplansky beweisen, über den \citeauthor{PedC*} in seinem Buch \citetitle{PedC*} schreibt:
\begin{quote}
	\foreignblockquote{english}[{\cite[S. 25]{PedC*}}]{\itshape The density theorem is Kaplansky's great gift to mankind. It can be used every day, and twice on Sundays.}
\end{quote}

\begin{satz}[label=satz:624,name={Kaplanskys Dichtesatz}]
	Sei $\mathcal{H}$ ein Hilbertraum, $A \subset \mathcal{B}(\mathcal{H})$ eine $C^*$-Algebra\marginnote{$B$ ist vN-Algebra nach Aufgabe 1\newline von Blatt 10} und $B= \overline{A}^\so=\overline{A}^\wo$ die von $A$ in $\mathcal{B}(\mathcal{H})$ erzeugte von-Neumann-Algebra.
	Dann gilt
	\begin{enumerate}[(i),itemsep=2pt]
		\item $\overline{A_\sa}^\so = \overline{A_\sa}^\wo= B_\sa$
		\item $\overline{A_\sa^1}^\so= \overline{A^1_\sa}^\wo = B^1_\sa$
		\item $\overline{A_+^1}^\so = \overline{A_+^1}^\wo= B_+^1$
		\item $\overline{A^1}^\so = \overline{A^1}^\wo = B^1$
		\item $\overline{\mathcal{U}(A)}^\so = \mathcal{U}(B)$ falls $\ind_{\mathcal{B}(\mathcal{H})} \in A$
	\end{enumerate}
\end{satz}
\begin{beweis}
	Die starken und schwachen Abschlüsse stimmen jeweils überein nach \autoref{prop:620}.
	\begin{enumerate}[(i)]
		\item Falls $a \in B_\sa$, so existiert $(a_\lambda)_\Lambda \subset A$ mit $a_\lambda \grenzw{\wo}a$.
		 Dann gilt aber auch $a_\lambda^* \grenzw{\wo}a^*=a$ nach \autoref{erinnerung:616} und damit 
		 \[
		 	A_\sa \ni\sfrac{1}{2} \enbrace*{a_\lambda + a_\lambda^*} \grenzw{\wo} a
		 \]
		 Also ist $B_\sa \subset \overline{A_\sa}^\wo \subset B_\sa$, wobei die zweite Inklusion offensichtlich ist.
		 \item Sei $a \in B^1_\sa$. Nach (i) existiert ein Netz $(a_\lambda)_\Lambda \subset A_\sa$ mit $a_\lambda \grenzw{\so}a$.
		 Sei $f \in C_b(\mathbb{R})$ gegeben durch \cref{fig:624}.
		 \begin{figure}[bht]
		 	\centering{
			\begin{tikzpicture}
				\draw (-2,0) -- (2,0);
				\draw (0,-1.2)-- (0,1.2);
				\draw (-.1,-1) node[left]{$-1$} -- ++(.2,0);
				\draw (-.1,1) node[left]{$1$} -- ++(.2,0);
				\draw (-1,-.1) node[below]{$-1$} -- ++ (0,.2);
				\draw (1,-.1) node[below]{$1$} -- ++ (0,.2);
				\draw[Firebrick2,very thick] (-2,-1) -- (-1,-1) -- (1,1) -- (2,1) node[below left]{$f$};
			\end{tikzpicture}
			\hspace{2cm}
			\begin{tikzpicture}
				\draw[] (-2,0) -- (2,0);
				\draw (0,-1.2)-- (0,1.2);
				\draw (-.1,-1) node[left]{$-1$} -- ++(.2,0);
				\draw (-.1,1) node[left]{$1$} -- ++(.2,0);
				\draw (-1,-.1) node[below]{$-1$} -- ++ (0,.2);
				\draw (1,-.1) node[below]{$1$} -- ++ (0,.2);
				\draw[Firebrick2,very thick] (-2,0) -- (0,0) -- (1,1) -- (2,1) node[below left]{$g$};
			\end{tikzpicture}
			\caption{Funktionen im Beweis von \autoref{satz:624}\label{fig:624}}}
		 \end{figure}
		 $f$ ist stark stetig nach \autoref{prop:623}.
		 Also ist $A_\sa^1 \ni  f(a_\lambda) \grenzw{\so} f(a)=a$ und damit $a \in \overline{A^1_\sa}^\so$.
		 Wir haben also $B^1_\sa \subseteq \overline{A^1_\sa}^\so$ gezeigt. Die andere Inklusion ist wieder klar.
		 \item Wie (ii) mit $g$ anstatt $f$, siehe \cref{fig:624}.
		 \item Sei $a \in B^1$.
		 Dann ist $\begin{psmallmatrix*}[l] 0 & a \\ a^* & 0 \end{psmallmatrix*} \in M_2(B)_\sa^1 \subset \mathcal{B} \enbrace*{\mathcal{H}^{\oplus 2}}$. 
		 Es gilt $\overline{M_2(A)}^\so=M_2(B)$ und nach (ii) existiert $(b_\lambda)_\Lambda \subset M_2(A)^1_\sa$ mit 
		 \(
		 	\begin{psmallmatrix}
		 		b_{\lambda,1,1} & b_{\lambda,1,2} \\
				b_{\lambda,2,1} & b_{\lambda,2,2}
		 	\end{psmallmatrix}
			=b_\lambda \grenzw{\so} \begin{psmallmatrix*}[l]
				0 & a \\
				a^* & 0
			\end{psmallmatrix*}
		 \).
		 Dann gilt
		 \[
		 	\begin{pmatrix}
		 		1 &0 \\
				0 & 0
		 	\end{pmatrix} 
			\begin{pmatrix}
		 		b_{\lambda,1,1} & b_{\lambda,1,2} \\
				b_{\lambda,2,1} & b_{\lambda,2,2}
			\end{pmatrix}
			\begin{pmatrix}
				0 & 0 \\
				1 & 0
			\end{pmatrix} =
			\begin{pmatrix}
				0 & b_{\lambda,1,2} \\
				0 & 0
			\end{pmatrix}
		 \]
		 Also ist $b_{\lambda,1,2} \in A^1$ und wir erhalten $b_{\lambda,1,2} \grenzw{\so} a$, also $a \in \overline{A^1}^\so$.
		 Wir haben somit die Inklusion $B^1 \subseteq \overline{A^1}^\so$ gezeigt; die andere ist wieder klar.
		 \item Nach \autoref{korr:613} (ii) gilt $\mathcal{U}(B) = \overline{\set*{e^{i \pi h} \given h \in B^1_\sa}}^{\norm*{\cdot}}$.
		 Zu $u \in \mathcal{U}(B)$ und $\varepsilon>0$ existiert ein $h \in B^1_\sa=\overline{A^1_\sa}^\so$ mit $\norm*{e^{i \pi h}-u}<\varepsilon$.
		 Aber $\enbrace*{t \mapsto e^{i \pi t}} \in C_b(\mathbb{R})$ ist stark stetig nach \autoref{prop:623} und somit ist $e^{i \pi h} \in \overline{\mathcal{U}(A)}^\so$.
		 Damit haben wir
		 \[
		 	\mathcal{U}(B) \subset \overline{\overline{\mathcal{U}(A)}^\so}^{\norm*{\cdot}} = \overline{\mathcal{U}(A)}^\so
		 \]
		 Die umgekehrte Inklusion ist wieder klar. \qedhere
	\end{enumerate}
\end{beweis}
\newpage
% section 6 (end)

\section{Reine Zustände und irreduzible Darstellungen} % (fold)
\label{sec:7}

Bevor man in dieses Kapitel abtaucht, ist es empfehlenswert, sich noch einmal an \autoref{def:510} zu erinnern, denn dort wurden die Begriffe \emph{irreduzibel} und \emph{zyklisch} für Darstellungen eingeführt.
Wenig überraschend wird auch die GNS-Konstruktion, siehe \autoref{gns} und \cref{sub:gns_zyklisch}, wieder eine wichtige Rolle spielen.

\begin{proposition}[label=prop:71,{name=[über positive Funktionale, die kleiner als Zustände sind]}]
	Sei $A$ eine $C^*$-Algebra, $\varphi \in S(A)$ ein Zustand und $\rho \in A^*$ ein weiteres positives Funktional.
	Falls $\varphi-\rho$ positiv ist, so existiert genau ein $0 \le h \le \ind \in \pi_\varphi(A)' \subset \mathcal{B}(\mathcal{H}_\varphi)$ mit\marginnote{Es gilt auch die Umkehrung, wie man sich leicht mit \ref{satz:37} überlegt} 
	\[
		\rho(a) = \skal[\big]{\xi_\varphi}{\pi_\varphi(a)h \xi_\varphi}
	\]
	wobei $\enbrace*{\mathcal{H}_\varphi,\pi_\varphi,\xi_\varphi}$ die GNS-Darstellung zu $\varphi$ ist.
\end{proposition}
\begin{beweis}
	Im Zuge der GNS-Konstruktion, siehe \autoref{gns} haben wir ein Ideal in $A$ definiert durch $N_\varphi = \set*{a \in A \given \varphi(a^*a)=0}$ und dann den Hilbertraum $\mathcal{H}_\varphi = \overline{\sfrac{A}{N_\varphi}}^{\raisemath{-2pt}{\skal*{\cdot}{\cdot}}}$ betrachtet.
	Definiere auf $\sfrac{A}{N_\varphi}$ eine hermitesche Sesquilinearform $\sigma$ durch 
	\[
		\sigma \enbrace*{a + N_\varphi, b+ N_\varphi} := \rho(a^* b)
	\]
	$\sigma$ ist wohldefiniert, da $0 \le \rho \le \varphi$ ist, woraus $N_\varphi \subset N_\rho$ folgt.
	Es gilt für $a,b \in A$
	\[
		\sigma \enbrace*{a+ N_\varphi, b+N_\varphi} = \rho(a^*b) \StackTextClap{C.S.}{\le} \rho(a^*a)^{\sfrac{1}{2}} \rho(b^*b)^{\sfrac{1}{2}} \le \varphi(a^*a)^{\sfrac{1}{2}} \varphi(b^*b)^{\sfrac{1}{2}} 
		= \norm*{a + N_\varphi} \cdot \norm*{b+ N_\varphi}
	\]
	Damit ist $\norm*{\sigma}\le 1$ und $\sigma$ lässt sich stetig auf $\mathcal{H}_\varphi$ fortsetzen.
	Für $\xi \in \mathcal{H}_\varphi$ ist $\eta \mapsto \sigma(\xi,\eta)$ eine stetige Linearform auf $\mathcal{H}_\varphi$.
	Nach Riesz existiert also genau ein Vektor $\xi_\sigma \in \mathcal{H}_\varphi$ mit $\sigma(\xi,\eta) = \skal*{\xi_\sigma}{\eta}$.
	$\xi \mapsto \xi_\sigma$ ist linear und beschränkt, also existiert genau ein $h \in \mathcal{B}(\mathcal{H}_\varphi)$ mit $h \xi= \xi_\sigma$.
	Es gilt
	\[
		\skal*{h \xi}{\eta} = \sigma(\xi,\eta) = \overline{\sigma(\eta,\xi)} = \overline{\skal*{h \eta}{\xi}} = \skal*{\xi}{h \eta} = \skal*{h^* \xi}{\eta}
	\]
	Also ist $h$ selbstadjungiert.
	% Es gilt
	% \[
	% 	\skal*{h^*\xi}{\xi}=\overline{\sigma(\xi,\xi)}=\sigma(\xi,\xi)= \skal*{h \xi}{\xi} \in \mathbb{R}
	% \]
	% Damit ist $h=h^*$ nach Polarisierung.
	Weiter gilt $\norm*{h} \le \norm*{\sigma} \le 1$ und für $a \in A$ haben wir
	\[
		\skal*{a + N_\varphi}{h \enbrace*{a+ N_\varphi}}= \sigma(a^*a)=\rho(a^*a)\ge 0
	\]
	Also ist $\skal*{\xi}{h \xi} \ge 0$ für $\xi \in \mathcal{H}_\varphi$, womit $h \ge 0$ ist nach \autoref{bem:63}. 
	Damit haben wir $0 \le h \le \ind$.
	Für $a,b,c \in A$ gilt 
	\begin{align}
		\skal[\big]{a + N_\varphi}{\pi_\varphi(c)h (b+N_\varphi)} = \skal[\big]{c^*a + N_\varphi}{h(b+N_\varphi)} = \rho \enbrace*{a^*cb} &= \skal[\big]{a + N_\varphi}{h(cb+N_\varphi)} \\
		&= \skal[\big]{a+ N_\varphi}{h \pi_\varphi(c)(b+ N_\varphi)}
	\end{align}
	Damit ist $\pi_\varphi(c) h = h \pi_\varphi(c)$, also ist $h$ in der Kommutante von $\pi_\varphi(A)$ wie behauptet.
	Um die behauptete Formel für $\rho$ zu beweisen, betrachten wir eine approximative Eins $(e_\lambda)_\Lambda \subset A^1_+$. 
	Dann ist $\xi_\varphi = \lim_\lambda e_\lambda+ N_\varphi \in \mathcal{H}_\varphi$ und es gilt
	\begin{align}
		\rho(a)= \lim_{\lambda,\lambda'} \rho \enbrace*{e_\lambda a e_{\lambda'}}= \lim_{\lambda,\lambda'} \skal[\big]{e_\lambda + N_\varphi}{h(a e_{\lambda'}+ N_\varphi)}
		&= \lim_{\lambda,\lambda'} \skal[\big]{e_\lambda+ N_\varphi}{h\,\pi_\varphi(a) (e_{\lambda'} + N_\varphi)} \\
		&= \skal[\big]{\xi_\varphi}{h\, \pi_\varphi(a)\,\xi_\varphi}
	\end{align}
	Für die Eindeutig sei $d \in \enbrace*{\pi_\varphi(A)'}^1_+$ mit $\rho(a)= \skal*{\xi_\varphi}{\pi_\varphi(a)d \xi_\varphi}$ für alle $a \in A$.
	Dann gilt für $a,b \in A$, da $\xi_\varphi = \lim_\lambda e_\lambda+ N_\varphi$ zyklisch ist
	\begin{align}
		\skal[\big]{a+ N_\varphi}{d(b+N_\varphi)} = \skal[\big]{\pi_\varphi(a) \xi_\varphi}{d \,\pi_\varphi(b) \xi_\varphi} = \skal[\big]{\xi_\varphi}{\pi_\varphi(a^*b)d \xi_\varphi} 
		&= \rho(a^*b) \\ &= \skal[\big]{a+ N_\varphi}{h(b+N_\varphi)}
	\end{align}
	Also ist $d=h$, womit die Eindeutigkeit gezeigt ist.
\end{beweis}

\begin{satz}[{name=[Darstellung als direkte Summe]}]
	Sei $(\mathcal{H},\pi)$ eine \hyperref[def:510]{nicht degenerierte} Darstellung einer $C^*$-Algebra $A$.
	Dann ist $(\mathcal{H},\pi)$ eine direkte Summe zyklischer Darstellungen.
\end{satz}
\begin{beweis}
	Sei $\mathcal{S} := \set[\big]{S \subset \mathcal{H}\setminus \set{0} \given {\overline{\pi(A) \xi}} \perp {\overline{\pi(A)\eta}} \text{ für alle }\xi \neq \eta \in S }$.
	$\mathcal{S}$ ist partiell geordnet bezüglich Inklusion und abgeschlossen bezüglich Vereinigungen von aufsteigenden Teilmengen.
	Nach Zorns Lemma besitzt $\mathcal{S}$ ein maximales Element $\tilde{S}$.
	Sei $\eta \in \enbrace[\big]{\bigcup_{\xi \in \tilde{S}} \overline{\pi(a)\xi}}^{\raisemath{-1pt}{\bot}}$, dann gilt für $a,b \in A$, $\xi \in \tilde{S}$.
	\[
		\skal[\big]{\pi(a) \eta}{\pi(b)\xi} = \skal[\big]{\eta}{ \pi(a^*b)\xi} =0
	\]
	Dann ist $\overline{\pi(A)\eta} \perp \overline{\pi(A)\xi}$ für alle $\xi \in \tilde{S}$.
	Da $\tilde{S}$ maximal ist, muss $\eta \in \tilde{S}$ oder $\eta=0$ sein.
	Falls $\eta \in \tilde{S}$, so gilt $\overline{\pi(A)\eta}\perp \overline{\pi(A)\eta}$, was ein Widerspruch ist, denn $(\mathcal{H},\pi)$ ist nicht degeneriert ($\Rightarrow  \pi(e_\lambda)\grenzw{\so}\ind_{\mathcal{B}(\mathcal{H})}$).
	Also ist $\eta=0$ und damit
	\[
		\mathcal{H} = \bigoplus_{\xi \in \tilde{S}} \overline{\pi(A)\xi}
	\]
	Die eingeschränkten Darstellungen $\pi(\cdot )\big|_{\overline{\pi(A)\xi}}$ haben jeweils $\xi \in \overline{\pi(A)\xi}$ als zyklischen Vektor.
\end{beweis}

\begin{satz}[{name=[Eigenschaften irreduzibler Darstellungen]},label=satz:73]
	Sei $(\mathcal{H},\pi)$ eine Darstellung einer $C^*$-Algebra $A$ mit $\pi \neq 0$. Dann gilt
	\begin{enumerate}[(i),itemsep=1pt]
		\item $(\mathcal{H},\pi)$ ist irreduzibel $\iff \pi(A)'= \mathbb{C} \cdot \ind_{\mathcal{B}(\mathcal{H})}$.\marginnote{$\iff \pi(A)$ $\so$-dicht in $\mathcal{B}(\mathcal{H})$, siehe \cref{sub:das_bild_irreduzibler_darstellungen}}
		\item $(\mathcal{H},\pi)$ ist irreduzibel $\implies$ jedes $0 \neq \xi \in \mathcal{H}$ ist zyklisch, also $\overline{\pi(A)\xi}=\mathcal{H}$.
	\end{enumerate}
\end{satz}
\begin{beweis}
	\leavevmode
	\begin{enumerate}[(i)]
		\item Für die Rückrichtung nehmen wir an, $(\mathcal{H},\pi)$ sei nicht irreduzibel. 
		Dann existiert $\tilde{\mathcal{H}} \subset \mathcal{H}$ abgeschlossen mit $\set*{0} \neq \tilde{\mathcal{H}}\neq \mathcal{H}$ und $\pi(A)\tilde{\mathcal{H}} \subset \tilde{\mathcal{H}}$.
		Setze $p := p_{\tilde{\mathcal{H}}}$, dann gilt $0 \neq p \neq \ind_{\mathcal{B}(\mathcal{H})}$ und 
		\[
			p \pi(a)=\enbrace[\big]{\pi(a)p}^* =\enbrace[\big]{p \pi(a)p}^*=p \pi(a) p=\pi(a)p 
		\]
		für $a \in A_\sa$. 
		Also ist $p	\in \pi(A_\sa)'$ und somit $p \in \pi(A)'$.
		Damit ist $p= \lambda \cdot \ind_{\mathcal{B}(\mathcal{H})}$ für ein $\lambda \in \mathbb{C}$.
		Da $p$ ein Projektion ist, folgt $\lambda \in \set*{0,1}$ und somit ist $p=0$ oder $p=\ind_{\mathcal{B}(\mathcal{H})}$, was einen Widerspruch darstellt.
		
		Für ein Projektion $p \in \mathcal{B}(\mathcal{H})$ gilt $p \in \pi(A)'$ genau dann, wenn $\overline{p \mathcal{H}} \subset \mathcal{H}$ invariant ist unter $\pi(A)$.
		Da $\pi(A)'$ eine von-Neumann-Algebra ist, ist sie nach \autoref{bem:614} der Abschluss des Spans ihrer Projektionen.
		Wenn die Darstellung irreduzibel ist, enthält $\pi(A)'$ also nur die trivialen Projektionen und es folgt $\pi(A)' = \mathbb{C} \cdot \id_{\mathcal{B}(\mathcal{H})}$.
		% Für die andere Implikation stellen wir fest, dass $\pi(A)' \subset \mathcal{B}(\mathcal{H})$ eine von-Neumann-Algebra ist.
		% Falls $\pi(A)' \neq \mathbb{C} \cdot \ind_{\mathcal{B}(\mathcal{H})}$, so existiert $h \in \pi(A)'$ selbstadjungiert, sodass $\sigma(h)$ mehr als einen Punkt enthält.
		% Ohne Einschränkungen sei $h \in \enbrace*{\pi(A)'}^1_+$ und $\set*{0,1} \in \sigma(h)$.\marginnote{Funktionalkalkül}
		% $p :=\chi_{(\sfrac{1}{2},1]}(h) \in \pi(A)'$ ist eine Projektion mit $0\neq p \neq \ind_{\mathcal{B}(\mathcal{H})}$.
		% Dann ist $\overline{p \mathcal{H}}$ ein $\pi(A)$-invarianter Unterraum und damit ist $(\mathcal{H},\pi)$ nicht irreduzibel.
		\item Sei $(\mathcal{H},\pi)$ irreduzibel.
		Dann ist $\overline{\pi(A)\xi}=\set*{0}$ oder $\overline{\pi(A)\xi}=\mathcal{H}$ für $\xi \in \mathcal{H}$, da die Darstellung irreduzibel ist.
		Da $\pi\neq 0$ ist, existiert ein $\eta \in \mathcal{H}$ mit $\pi(A)\eta \neq \set*{0}$.
		Also ist $\overline{\pi(A)\eta}=\mathcal{H}$.
		Damit ist $\pi$ nicht degeneriert und es folgt $\overline{\pi_n(A)\xi} = \mathcal{H}$ für $\xi\neq 0$, also ist $\xi$ zyklisch wie behauptet.\qedhere
	\end{enumerate}
\end{beweis}

\begin{definition}[{name=[reiner Zustand]}]
	Sei $A$ eine $C^*$-Algebra.
	Ein Zustand $\varphi \in S(A)$ heißt \bet{rein}\index{Zustand!reiner}, falls gilt: Für jedes $\rho \in A^*$ positiv mit $\rho \le \varphi$ existiert $\lambda \in [0,1]$ mit $\rho= \lambda \cdot \varphi$.
	Wir schreiben $PS(A) := \set*{\text{reine Zustände auf }A}$.\marginnote{engl. \emph{\textbf{p}ure \textbf{s}tates}}
\end{definition}

\begin{satz}[{name=[Charakterisierung reiner Zustände]},label=satz:75]
	Sei $A$ eine $C^*$-Algebra, $\varphi \in S(A)$. Dann gilt
	\[
		\varphi \text{ ist rein } \iff (\mathcal{H}_\varphi, \pi_\varphi) \text{ ist irreduzibel}  
	\]
\end{satz}
\begin{beweis}
	Sei $\varphi$ rein und $0 \le h \le \ind_{\mathcal{B}(\mathcal{H})} \in \pi_\varphi(A)'$.
	Dann definiert 
	\mapdef{\rho \colon A}{\mathbb{C}}{a}{\skal[\big]{\xi_\varphi}{h\,\pi_\varphi(a)\xi_\varphi}}{}
	ein positives Funktional auf $A$ mit $\rho \le \varphi$, siehe \autoref{prop:71}.
	$\varphi$ ist rein, also existiert $\lambda \in [0,1]$ mit $\rho =\lambda \cdot \varphi$.
	Für $a,b \in A$ erhalten wir:
	\begin{align}
		\skal*{a +N_\varphi}{h (b+ N_\varphi)} = \skal*{\pi(a) \xi_\varphi}{h\,\pi(b)\xi_\varphi} &= \skal*{\xi_\varphi}{h\,\pi(a^*b)\xi_\varphi} \\
		&= \rho(a^*b) =\lambda \cdot \varphi(a^*b) \\
		&= \lambda \cdot \skal*{a+ N_\varphi}{b + N_\varphi} \\
		&= \skal*{a + N_\varphi}{\lambda \cdot \ind_{\mathcal{B}(\mathcal{H})}(b+ N_\varphi)}
	\end{align}
	Also ist $\skal*{\xi}{h \eta}= \skal*{\xi}{\lambda \cdot \ind_{\mathcal{B}(\mathcal{H}_\varphi)}(\eta)}$ für $\xi,\eta \in \mathcal{H}_\varphi$ und damit $h= \lambda \cdot \ind_{\mathcal{B}(\mathcal{H}_\varphi)}$.
	Also ist $\pi_\varphi(A)'= \mathbb{C} \cdot \ind_{\mathcal{B}(\mathcal{H})}$ und mit \autoref{satz:73} (i) folgt, dass $(\mathcal{H}_\varphi,\pi_\varphi)$ irreduzibel ist.
	
	Für die andere Implikation sei $(\mathcal{H}_\varphi,\pi_\varphi)$ irreduzibel und $\rho \in A^*$ positiv mit $\rho \le \varphi$.
	Nach \autoref{prop:71} existiert $0 \le h \le \ind \in \pi_\varphi(A)'$ mit $\rho(a) = \skal*{\xi_\varphi}{h\,\pi_\varphi(a)\xi_\varphi}$ für $a \in A$.
	Nach \autoref{satz:73}~(i) gilt $\pi_\varphi(A)'=\mathbb{C} \cdot \ind_{\mathcal{B}(\mathcal{H}_\varphi)}$, also $h= \lambda \cdot \ind_{\mathcal{B}(\mathcal{H})}$ für ein $\lambda \in [0,1]$ und $\rho(a)= \skal*{\xi_\varphi}{h\,\pi_\varphi(a)\xi_\varphi}= \lambda \cdot \skal*{\xi_\varphi}{\pi_\varphi(a)\xi_\varphi}= \lambda \cdot \varphi(a)$.
\end{beweis}

\begin{satz}[{name=[{Reine Zustände von $C_0(X)$}]}]
	Sei $X$ ein lokalkompakter Hausdorffraum.
	Dann gilt $PS\enbrace[\big]{C_0(X)}=\Omega \enbrace[\big]{C_0(X)}$.
\end{satz}
\begin{beweis}
	Sei $\varphi \in PS \enbrace*{C_0(X)}$ ein reiner Zustand.
	Dann gilt
	\[
		\pi_\varphi \enbrace[\big]{C_0(X)} \subset \pi_\varphi \enbrace[\big]{C_0(X)}' \StackText{\ref{satz:73}, \ref{satz:75}}{=\joinrel=\joinrel=} \mathbb{C} \cdot \ind_{\mathcal{B}(\mathcal{H}_\varphi)} \subset \mathcal{B}(\mathcal{H}_\varphi)
	\]
	und $\mathcal{B}(\mathcal{H}_\varphi) \subset \enbrace*{\mathbb{C}\cdot \ind_{\mathcal{B}(\mathcal{H}_\varphi)}}' \subset \pi_\varphi \enbrace[\big]{C_0(X)}'=\mathbb{C}\cdot \ind_{\mathcal{B}(\mathcal{H}_\varphi)}$.
	Also ist $\mathcal{B}(\mathcal{H}_\varphi)= \mathbb{C} \cdot \ind_{\mathcal{B}(\mathcal{H}_\varphi)} \cong \mathbb{C}$.
	Wir zeigen jetzt, dass $\varphi$ multiplikativ und damit ein \hyperref[def:charakter]{Charakter} ist:
	Für $a,b \in C_0(X)$ gilt 
	\begin{align}
		\varphi(ab) = \skal*{\xi_\varphi}{\pi_\varphi(ab)\xi_\varphi} = \Underbracket{\pi_\varphi(a)}{\in \mathbb{C}}\cdot  \skal*{\xi_\varphi}{\pi_\varphi(b) \xi_\varphi} 
		&= \pi_\varphi(a) \Underbracket{\skal*{\xi_\varphi}{\xi_\varphi}}{=1} \cdot \skal*{\xi_\varphi}{\pi_\varphi(b)\xi_\varphi} \\
		&= \skal*{\xi_\varphi}{\pi_\varphi(a)\xi_\varphi} \cdot \skal*{\xi_\varphi}{\pi_\varphi(b)\xi_\varphi} \\
		&= \varphi(a) \cdot \varphi(b)
	\end{align}
	Also ist $\varphi \in \Omega \enbrace*{C_0(X)}$ ein Charakter. 
	Sei nun umgekehrt $\chi \in \Omega(C_0(X))$ und $\rho \in C_0(X)^*$ positiv mit $\rho \le \chi$.
	Für $a \in C_0(X)$ gilt dann
	\[
		\abs*{\rho(a)} \StackText{\ref{satz:53}}{\le} \norm*{\rho}^{\sfrac{1}{2}} \rho(a^*a)^{\sfrac{1}{2}} \le \chi(a^* a)^{\sfrac{1}{2}}= \enbrace*{\overline{\chi(a)}\cdot  \chi(a)}^{\sfrac{1}{2}} \le \abs*{\chi(a)} 
	\]
	Also ist $\ker \chi \subset \ker \rho$.
	$\rho$ induziert nun ein Funktional $\tilde{\rho}$ auf $\sfrac{C_0(X)}{\ker \chi} \cong \mathbb{C}$ mit $\rho = \tilde{\rho} \circ \chi$.
	\[
		\begin{tikzcd}[column sep=.5em]
			C_0(X) \ar[rr,"\rho"] \drar["\chi",two heads]& & \mathbb{C} \\
			& \sfrac{C_0(X)}{\ker \chi} \cong \mathbb{C} \urar["\tilde{\rho}"]
		\end{tikzcd}
	\]
	Damit ist $\rho = \lambda \cdot \chi$ für ein $\lambda \in \mathbb{C}$.
	Für $a \in C_0(X)$ mit $\chi(a)=1$ gilt $0\le\rho(a^*a)= \lambda \cdot \chi(a^*a)=\lambda$ und außerdem $\rho(a^*a) \le \chi(a^*a)=1$, womit $0 \le \lambda\le 1$ ist.
	Daraus folgt, dass $\chi$ rein ist.
\end{beweis}

\begin{definition}[{name=[Seiten und Extrempunkte]}]
	Sei $X$ ein hausdorffscher lokalkonvexer Vektorraum, das heißt $X$ ist ein topologischer Vektorraum und die Topologie ist durch eine punkte trennende Familie $(p_\lambda)_\Lambda$ von Halbnormen induziert.
	Sei $C \subset X$ eine nichtleere, (kompakte,) konvexe Teilmenge und sei $F \subset C$ eine nichtleere konvexe Teilmenge.
	$F$ heißt \Index{Seite} von $C$, falls gilt:  
	Für $x \in F$, $y,z \in C$ und $t \in (0,1)$ gilt 
	\[
		x= t \cdot y +(1-t) \cdot z \implies y,z \in F
	\]
	$x \in C$ heißt \Index{Extrempunkt} von $C$, falls $\set*{x}$ eine Seite von $C$ ist.
	Wir schreiben 
	\[
		\partial_e C := \set[\big]{\text{Extrempunkte von } C}
	\]
\end{definition}

\begin{proposition}[{name=[grundlegende Eigenschaften von Seiten]},label=prop:78]
	Seien $X,C$ wie oben.
	\begin{enumerate}[(i),itemsep=0pt]
		\item $C$ ist eine Seite von $C$.
		\item Eine nichtleerer Durchschnitt von Seiten ist wieder eine Seite.
		\item Für eine Seite $F$ von $C$ gilt $\partial_e F \subset \partial_e C$.
	\end{enumerate}
\end{proposition}
\begin{beweis}
	Alle drei Aussagen ergeben sich direkt aus der Definition.
\end{beweis}

\begin{lemma}[label=lem:79,{name=[durch Funktional definierte Seite]}]
	Seien $X$ und $C$ wie oben und $\varphi \colon X \to \mathbb{C}$ ein stetiges lineares Funktional.
	Setze $M :=\sup_{x \in C} \re \varphi(x)$.
	Dann ist $F:= \set*{x \in C \given \varphi(x)=M}$ eine kompakte Seite von $C$.
\end{lemma}
\begin{beweis}
	Da $C$ kompakt ist, existiert ein $x_0 \in C$, sodass $\re \varphi(x_0)=M$, also ist $F \neq \emptyset$.
	Offensichtlich ist $F$ konvex.
	Da $\varphi$ stetig ist, ist $F \subset C$ abgeschlossen.
	Seien nun $x \in F$, $y,z \in C$ und $t \in (0,1)$ mit $x= t \cdot y +(1-t) \cdot z$.
	Falls $y$ oder $z$ nicht in $F$ liegt, so gilt $\re \varphi(y) <M$ oder $\re \varphi(z)<M$.
	Damit ist 
	\[
		M = \re \varphi(x) = t \cdot \re \varphi(y) + (1-t) \cdot \re \varphi(z) <M
	\]
	Dies ist ein Widerspruch, also müssen $y$ und $z$ in $F$ liegen und damit ist $F$ eine Seite.
\end{beweis}

\begin{satz}[{name={{Krein-Milman}}},label=satz:krein_milman]
	Sei $X$ ein hausdorffscher lokalkonvexer Vektorraum\marginnote{nach \href{https://de.wikipedia.org/wiki/Mark_Grigorjewitsch_Krein}{Mark Grigorjewitsch \textsc{Krein}} und \href{https://de.wikipedia.org/wiki/David_Milman}{David \textsc{Milman}}} und $C \subset X$ eine nichtleere, kompakte, konvexe Teilmenge.
	Dann gilt 
	\[
		C=\overline{\conv(\partial_e C)}
	\]
\end{satz}
\begin{beweis}
	Sei $\Lambda := \set*{\text{kompakte Seiten von }C}$.
	Dann ist $C \in \Lambda$, das heißt $\Lambda \neq \emptyset$.
	$\Lambda$ ist partiell geordnet bezüglich $\supset$.
	Sei $\Gamma \subset \Lambda$ eine totalgeordnete Teilmenge, dann gilt unter der Annahme $\bigcap_{\gamma \in \Gamma} \gamma =\emptyset$
	\[
		C= C \setminus \bigcap_{\gamma \in \Gamma} \gamma = \bigcup_{\gamma \in \Gamma} \Underbracket{C \setminus \gamma}{\text{offen in }C}\StackText{C \text{ kpt.} }{=\joinrel=\joinrel=} 
		\bigcup_{i=1}^n C \setminus \gamma_i \StackText{totalgeord.}{=\joinrel=\joinrel=} C \setminus \gamma_0 \neq C
	\]
	Damit ist $\bigcap_{\gamma \in \Gamma} \gamma \neq \emptyset$ eine kompakte Seite von $C$ und $\Gamma$ besitzt eine untere Schranke in $\Lambda$.
	Also ist $\Lambda$ induktiv geordnet und mit dem Lemma von Zorn folgt die Existenz eines minimalen Elements $F_0$.
	Behauptung: $F_0= \set*{x_0}$ für ein $x_0 \in C$.
	Angenommen es gäbe $x_0 \neq y_0 \in F_0$. 
	Dann existiert $\varphi \colon X \to \mathbb{C}$ stetig und linear mit $\re \varphi(x_0) \neq \re \varphi(y_0)$ (Hahn-Banach für lokalkonvexe Vektorräume\footnote{siehe zum Beispiel \cite[Cor. A.9, S. 271]{Murphy}}).
	Dann ist 
	\[
		F_1 :=  \set*{x \in F_0 \given \re \varphi(x)=M=\sup\nolimits_{x \in F_0} \re \varphi(x)}
	\]
	eine kompakte Seite von $F_0$ nach \autoref{lem:79} und \autoref{prop:78}.
	Es gilt $F_1 \subset F_0$, also folgt wegen der Minimalität $F_0=F_1$ und damit $M= \re \varphi(x_0) = \re \varphi(y_0)$.
	Dies ist ein Widerspruch und somit kann $F_0$ nur genau ein Element $x_0 \in \partial_e C$ besitzen.
	
	Angenommen es gibt ein $z \in C \setminus \overline{\conv(\partial_e C)}$.
	Dann existiert (wieder nach Hahn-Banach für lokalkonvexe Räume\footnote{siehe \cite[Th. A.7, S. 270]{Murphy}}) ein stetiges lineares $\psi \colon X \to \mathbb{C}$ mit 
	\[
		\sup\set*{\re \psi(y) \given y \in \overline{\conv(\partial_e C)}} < \re \psi(z)
	\]
	Setze $M' := \sup_{x \in C} \re \psi(x)$ und $F':= \set*{x \in C \given \re \psi(x) =M'}$.
	Dann ist $F'$ eine kompakte Seite von $C$ nach \autoref{lem:79}.
	Wie oben sieht man, dass $F'$ einen Extrempunkt $y_0$ besitzt, der dann auch ein Extrempunkt von $C$ ist, also $y_0 \in \partial_e C$.
	\[
		M' = \re \psi(y_0) \le \sup \set*{\re \psi(y) \given y \in \overline{\conv(\partial_e C)}} < \re \psi(z) \le M'
	\]
	Dies ist ein Widerspruch und damit folgt die Behauptung.
\end{beweis}

\begin{satz}[label=satz:711,{name=[reine Zustände sind Extrempunkte]}]
	Sei $A$ eine $C^*$-Algebra und $S:= \set*{\varphi \in A^* \given \varphi \text{ ist positiv}, \norm*{\varphi}\le 1}$.
	Dann ist $S$ konvex und $\w^*$-kompakt. 
	Weiter gilt 
	\[
		\partial_e S = PS(A) \cup \set*{0} \quad \text{ und } \quad  S= \overline{\conv(\partial_e S)}^{\w^*}
	\]
	Falls $A$ unital ist, so gilt $\partial_e \enbrace*{S(A)} =PS(A)$ und $S(A) = \overline{\conv \enbrace*{PS(A)}}^{\raisemath{-.4ex}{\w^*}}$. 
\end{satz}
\begin{beweis}
	Es ist $S \subset B^1(A^*)$ $\w^*$-abgeschlossen, wobei $B^1(A^*)$ nach Banach-Alaoglu $\w^*$-kompakt ist.
	Also ist $S$ auch $\w^*$-kompakt. 
	Offensichtlich ist $S$ außerdem konvex.
	Wir zeigen nun, dass für die Extrempunkte von $S$ gilt: $\partial_e S = PS(A) \cup \set*{0}$:
	\begin{itemize}
		\item Es gilt $0 \in \partial_e S$:
		Falls $0=t \cdot \varphi + (1-t) \cdot \psi$ mit $\varphi,\psi \in S$ und $t \in (0,1)$, so gilt
		\[
			\varphi\big|_{A_+}\!\! = \psi\big|_{A_+}\!\!=0 \implies \varphi=\psi=0
		\]
		\item Es gilt $PS(A) \subset \partial_e S$: Sei $\rho \in PS(A)$ und es gelte $\rho = t \cdot \varphi + (1-t) \cdot \psi$ mit $\varphi,\psi \in S$ und $t \in (0,1)$.
		Dann gilt $t \cdot \varphi, (1-t) \cdot \psi\le \rho$, also folgt 
		\[
			t \cdot \varphi = \lambda_\varphi \cdot \rho \quad \text{ und }\quad  (1-t) \cdot \psi = \lambda_\psi \cdot \rho 
		\]
		für geeignete Skalare $\lambda_\varphi, \lambda_\psi \in [0,1]$.
		Es gilt $1= \norm*{\rho} \le t \cdot \norm*{\varphi} + (1-t) \cdot \norm*{\psi} \le 1$, also ist $\norm*{\varphi}=\norm*{\psi}=1$.
		Damit folgt weiter $t=\lambda_\varphi$ und $(1-t)=\lambda_\psi$ und somit ist $\varphi=\rho=\psi$, also $\rho \in \partial_e S$.
		\item Es gilt $\partial_e S \setminus \set*{0} \subset PS(A)$: Für $0\neq\tau \in \partial_e S$ ist $0<\norm*{\tau}\le1$.
		Es gilt dann
		\[
			\tau= \norm*{\tau} \cdot \Underbracket{\frac{1}{\norm*{\tau}} \cdot \tau}{\in S} + \enbrace[\big]{1- \norm*{\tau}} \cdot 0 
		\]
		Falls $\norm*{\tau}<1$, so ist $\norm*{\tau} \in (0,1)$ und $\tau=0$, da $\tau$ ein Extrempunkt ist.
		Dies ist ein Widerspruch und es folgt also $\tau \in S(A)$.
		Sei nun $0\neq \rho \le \tau$ ein positives Funktional auf $A$.
		Dann gilt $1- \norm*{\rho} = \lim_\lambda \tau(e_\lambda) - \rho(e_\lambda)= \norm*{\tau- \rho}$, wobei $(e_\lambda)_\Lambda \subset A$ eine approximative Eins ist.
		Es folgt
		\[
			\tau = \norm*{\rho} \cdot \Underbracket{\frac{1}{\norm*{\rho}} \cdot \rho}{\in S} + \enbrace[\big]{1- \norm*{\rho}} \cdot \Underbracket{\frac{1}{\norm*{\tau - \rho}} \cdot (\tau - \rho)}{\in S}  
		\]
		Da $\tau \in \partial_e S$ ist, folgt $\tau = \sfrac{1}{\norm*{\rho}} \cdot \rho$ und somit $\norm*{\rho} \cdot \tau =\rho$, also ist $\tau \in PS(A)$.
	\end{itemize}
	Nach dem \hyperref[satz:krein_milman]{Satz von Krein-Milman (\ref*{satz:krein_milman})} ist $S= \overline{\conv(\partial_e S)}^{\raisemath{-.4ex}{\w^*}}\!$.
	Sei nun $A$ unital.
	Dann ist $S(A)$ nichtleer, konvex und $\w^*$-kompakt, denn aus $A$ unital folgt $\norm*{\varphi}=\varphi(\ind)$.
	Mit \hyperref[satz:krein_milman]{Krein-Milman} folgt also
	\[
		S(A) = \overline{\conv \enbrace[\big]{\partial_e \enbrace*{S(A)}}}^{\w^*}
	\]
	Aber $S(A) \subset S$ ist eine Seite, also gilt 
	\[
		PS(A) =S(A) \cap \partial_e S \StackText{\ref{prop:78}(iii)}{=\joinrel=} \partial_e \enbrace*{S(A)} \qedhere
	\]
\end{beweis}

\begin{satz}[label=satz:712,{name=[{reiner Zustand der Norm von $a \in A$ realisiert}]}]
	Sei $A$ eine $C^*$-Algebra und $0\neq a \in A_+$.
	Dann existiert $\rho \in PS(A)$ mit $\rho(a)=\norm*{a}$.
\end{satz}
\begin{beweis}
	Die folgende Menge mit $S$ wie in \autoref{satz:711} ist $\w^*$-abgeschlossen und konvex:
	\[
		F := \set[\big]{\varphi \in S \given \varphi(a)=\norm*{a}} \subset S
	\]
	$F$ ist nichtleer nach \autoref{satz:55} und $\w^*$-kompakt nach Banach-Alaoglu.
	Mit \hyperref[satz:krein_milman]{Krein-Milman} folgt, dass $F$ einen Extrempunkt $\rho$ besitzt.
	Wie in \autoref{lem:79} zeigt man, dass $F$ eine Seite von $S$ ist.
	Nach \autoref{prop:78} ist $\rho$ ein Extrempunkt von $S$.
	Da $\rho\neq0$ ist, folgt also $\rho \in PS(A)$ nach \autoref{satz:711}.
\end{beweis}

\begin{korollar}[label=korr:713,{name=[{irreduzible Darstellung, die die Norm von $a \in A$ erhält}]}]
	Sei A eine $C^*$-Algebra und $a \in A$.
	Dann existiert eine irreduzible Darstellung $(\mathcal{H},\pi)$ von $A$ mit $\norm*{a}= \norm*{\pi(a)}$.
\end{korollar}
\begin{beweis}
	Ohne Einschränkungen sei $a \neq 0$.
	Nach \autoref{satz:712} existiert ein reiner Zustand $\rho \in PS(A)$ mit $\rho(a^*a)=\norm*{a^*a}=\norm*{a}^2$.
	Die GNS-Darstellung $(\mathcal{H}_\rho,\pi_\rho)$ ist irreduzibel nach \autoref{satz:75}.
	Es gilt
	\[
		\norm*{a}^2 = \rho(a^*a) = \skal*{\xi_\rho}{\pi_\rho(a^*a)\xi_\rho} = \norm*{\pi_\rho(a) \xi_\rho}^2 \le \norm*{\pi_\rho(a)}^2 \le \norm*{a}^2 \qedhere
	\]
\end{beweis}


\begin{korollar}[label=korr:714,{name=[Fortsetzung reiner Zustände]}]
	Sei $B \subset A$ eine $C^*$-Unteralgebra und $\rho \in PS(B)$.
	Dann existiert eine Fortsetzung $\rho' \in PS(A)$.
	Falls $B\subher A$, so ist $\rho'$ eindeutig bestimmt. 
\end{korollar}
\begin{beweis}
	Die Menge $F:= \set*{\varphi \in S(A) \given \varphi|_{B}=\rho}$ ist nichtleer nach \autoref{satz:57} und es ist $F \subset S$ $\w^*$-abgeschlossen, also auch $\w^*$-kompakt (vergleiche \autoref{satz:711}).
	Weiter ist $F$ auch wieder konvex.
	$F$ ist nun eine Seite von 
	\[
		S = \set*{\varphi \colon A \to \mathbb{C}\given \text{ positiv, linear, } \norm*{\varphi}\le 1} 
	\]
	wie in Aufgabe 5 von Blatt 12 gezeigt (siehe \cref{sub:aufg5_12}).
	Mit \hyperref[satz:krein_milman]{Krein-Milman} und \autoref{prop:78} folgt, dass ein $\rho' \in \partial_e F \subset\partial_e S \setminus\set*{0}=PS(A)$ existiert.
	Die Eindeutigkeit folgt direkt mit \autoref{satz:58}.
\end{beweis}

\begin{proposition}[label=prop:715,{name=[schwache Variante des Transitivitätssatzes]}]
	Sei $\mathcal{H}$ ein Hilbertraum und $\xi_1,\ldots ,\xi_n,\eta_1 ,\ldots \eta_n \in \mathcal{H}$, wobei $\set*{\xi_1, \ldots ,\xi_n}$ ein Orthonormalsystem bildet.
	Dann existiert ein $x \in \mathcal{B}(\mathcal{H})$ mit $x(\xi_i)=\eta_i$ für $i=1,\ldots,n$ und 
	\[
		\norm*{x} \le \sqrt{2n} \cdot \max_{i=1,\ldots ,n} \norm*{\eta_i}
	\]
	Falls $y \in \mathcal{B}(\mathcal{H})_\sa$ mit $y(\xi_i)=\eta_i$ existiert, so kann man $x$ selbstadjungiert wählen.
\end{proposition}
\begin{beweis}
	Wir benutzen in diesem Beweis die aus Quantenphysik stammende \hrefsymX{https://de.wikipedia.org/wiki/Dirac-Notation}{\emph{Bra-Ket-Notation}}\index{Bra-Ket-Notation}.
	Setze also $x := \sum_{i=1}^{n} \ket*{\eta_i} \bra*{\xi_i}$, wobei der Operator $\ket*{\eta} \bra*{\xi} \in \mathcal{B}(\mathcal{H})$ gegeben ist durch
	\[
		\enbrace[\Big]{\ket[\big]{\eta} \bra[\big]{\xi}} (\zeta) := \skal*{\xi}{\zeta} \cdot \eta
	\]
	Dann gilt $x(\xi_i) = \eta_i$ für $i=1,\ldots ,n$, da die $\xi_i$ orthogonal sind.
	Für $\xi \in \mathcal{H}$ gilt weiter 
	\begin{align}
		\norm*{x \xi} = \norm*{\sum_{i=1}^{n} \bracket{\xi_i}{\xi} \eta_i} \le \sum_{i=1}^{n} \abs[\big]{ \bracket{\xi_i}{\xi}} \cdot \norm*{\eta_i}\enspace &\StackTextClap{\href{https://de.wikipedia.org/wiki/Hölder-Ungleichung}{Hölder}}{\le}\enspace \enbrace*{\sum_{i=1}^{n} \abs[\big]{\bracket{\xi_i}{\xi}}^2 \cdot \sum_{i=1}^{n} \norm*{\eta_i}^2}^{\sfrac{1}{2}}\\
		&\le \norm*{\xi} \cdot \sqrt{n} \cdot \max_i \norm*{\eta_i}
	\end{align}
	Es folgt $\norm*{x} \le \sqrt{n} \cdot \max_i \norm*{\eta_i}$.\todo{ich bin der Meinung, dass aus jeder Summe ein $n$ stammt, sodass am Ende keine Wurzel mehr auftritt}
	Existiert nun $y \in \mathcal{B}(\mathcal{H})_\sa$ mit $y(\xi_i) =\eta_i$, dann gilt
	\[
		x = \sum_{i=1}^{n} \ket{\eta_i} \bra{\xi_i} = \sum_{i=1}^{n} \ket{y \xi_i} \bra{\xi_i} = y \Underbracket{\enbrace*{\sum_{i=1}^{n} \ket{\xi_i} \bra{\xi_i}}}{=:p} =y p
	\]
	wobei $p$ die Projektion auf $\Span \set*{\xi_1, \ldots ,\xi_n}$ ist. 
	Dann ist $x' := yp + py - pyp$ selbstadjungiert und es gilt $x' \xi_i = \eta_i + p \eta_i - p \eta_i=\eta_i$.
	Für die Norm erhalten wir
	\begin{align}
		\norm*{x'}^2 = \norm*{yp + py - pyp}^2 = \norm[\big]{ \enbrace*{yp+py-pyp} \enbrace*{yp + py - pyp}}&= \norm[\big]{ypy + py(\ind -p) yp} \\
		&\le \norm*{ypy} + \norm*{py(\ind-p)^2yp} \\
		&\le \norm*{xx^*} + \norm*{x^*x} = 2 \cdot \norm*{x}^2 \\
		&\le 2 \cdot n \cdot \enbrace[\big]{\max\nolimits_i \norm*{\eta_i}}^2
	\end{align}
	Damit erhalten wir insgesamt die gewünschte Abschätzung.
\end{beweis}

\begin{satz}[{name={Kadisons Transitivitätssatz}},label=satz:716]
	Sei $A$ eine $C^*$-Algebra und $\pi \colon A \to \mathcal{B}(\mathcal{H})$ eine irreduzible Darstellung.
	Seien $\xi_1,\ldots ,\xi_n, \eta_1, \ldots , \eta_n \in \mathcal{H}$, wobei die $\xi_i$ linear unabhängig sind.
	Dann existiert ein $x \in A$ mit $\pi(x)(\xi_i) = \eta_i$ für $i=1,\ldots ,n$.
	\begin{itemize}
		\item Falls $y \in \mathcal{B}(\mathcal{H})_\sa$ existiert mit $y(\xi_i)=\eta_i$, so kann man auch $x \in A_\sa$ wählen.
		\item Falls $A$ und $\pi$ unital sind und es ein $u \in \mathcal{U} \enbrace*{\mathcal{B}(\mathcal{H})}$ gibt mit $u(\xi_i)=\eta_i$, so kann man $x$ unitär wählen (tatsächlich sogar $x=e^{i h}$ mit $h \in A_\sa$).
	\end{itemize}
\end{satz}
\begin{beweis}
	Wir zeigen verschiedene Spezialfälle, die wir schrittweise verallgemeinern:
	\begin{enumerate}[label={Teil \Roman*.},wide,font=\bfseries\scshape,itemsep=2pt]
		\item \label{enum:kad:1} Seien zunächst die $\xi_i$ orthonormal und es existiere ein $y \in \mathcal{B}(\mathcal{H})_\sa$ mit $y(\xi_i)=\eta_i$.
		Wir dürfen wegen Linearität $\max_i \norm*{\eta_i} \le (2n)^{- \sfrac{1}{2}}$ annehmen.
		Zu $\varepsilon>0$ definiere
		\[
			U_\varepsilon := \set[\big]{z \in \mathcal{B}(\mathcal{H}) \given \norm*{z(\xi_i)} < \varepsilon , i=1,\ldots ,n}
		\]
		Dann ist $U_\varepsilon$ eine stark-offene Umgebung von $0 \in \mathcal{B}(\mathcal{H})$.
		Es ist $\pi(A)$ $\so$-dicht in $\mathcal{B}(\mathcal{H})$, wie in einem kleinen Exkurs in \cref{sub:das_bild_irreduzibler_darstellungen} gezeigt.
		Mit \hyperref[satz:624]{Kaplansky} folgt, dass 
		\[
			\overline{\pi(A_\sa^1)}^\so = \mathcal{B}(\mathcal{H})_\sa^1
		\]
		Insbesondere gilt:
		Zu $w \in \mathcal{B}(\mathcal{H})_\sa$ existiert $w' \in A_\sa$ mit $\pi(w')-w \in U_\varepsilon$ und $\norm*{w'} \le \norm*{w}$.
		Nach \autoref{prop:715} existiert $v_0 \in \mathcal{B}(\mathcal{H})_\sa$ mit $v_0(\xi_i)=\eta_i$ und $\norm*{v_0} \le \sqrt{2n} \cdot \max_i \norm*{\eta_i} \le 1$.
		Also existiert ein $x_0 \in A_\sa$ mit $\pi(x_0)-v_0 \in U_{\sfrac{1}{(2 \sqrt{2n})}}$ und $\norm*{x_0} \le 1$.
		Wir konstruieren nun induktiv $(x_k)_k \subset A_\sa$, $(v_k)_k \subset \mathcal{B}(\mathcal{H})_\sa$ mit 
		\[
			\norm*{x_k}, \norm*{v_k} \le 2^{-k} \qquad  \pi(x_k) -v_k \in U_{2^{-(k+1)} (2n)^{- \sfrac{1}{2}}} \qquad v_k(\xi_i) = \enbrace[\big]{v_{k-1} - \pi(x_{k-1})}(\xi_i)
		\]
		$x_0,$ und $v_0$ sind bereits konstruiert.
		Falls auch $x_0, \ldots ,x_r, v_0, \ldots ,v_r$ bereits konstruiert sind, so existiert nach \autoref{prop:715} ein $v_{r+1} \in \mathcal{B}(\mathcal{H})_\sa$ mit 
		$v_{r+1}(\xi_i)= \enbrace*{v_r- \pi(x_r)}(\xi_i)$ für $i=1,\ldots ,n$ und 
		\[
			\norm*{v_{r+1}} \le \sqrt{2n} \cdot \max_{i=1,\ldots ,n} \norm*{\enbrace[\big]{v_r - \pi(x_r)}(\xi_i)} \le \sqrt{2n} \cdot 2^{-(r+1)} (2n)^{- \sfrac{1}{2}} = 2^{-(r+1)}
		\]
		Also existiert $x_{r+1} \in A_\sa$ mit $\pi(x_{r+1}) - v_{r+1} \in U_{2^{-(r+2)} \cdot (2n)^{-\sfrac{1}{2}}}$ und $\norm*{x_{r+1}} \le 2^{-(r+1)}$.
		Dies schließt die Induktion ab.
		Nach Konstruktion konvergiert $x := \sum_{r=0}^{\infty} x_r$ in $A_\sa$ und es gilt
		\begin{align}
			\eta_i - \pi(x)(\xi_i) 
			&= \lim_{k \to \infty} \enbrace[\bigg]{\eta_i - \sum_{r=0}^{k} \pi(x_r) \xi}\\ 
			&= \lim_{k \to \infty} \enbrace[\Bigg]{\Underbracket{\enbrace[\bigg]{ \Underbracket{\enbrace[\big]{v_o(\xi_i) - \pi(x_0)(\xi_i)}}{=v_1(\xi_i)} - \pi(x_1)(\xi_i)}}{=v_2(\xi_i)} - \ldots - \pi(x_k)(\xi_i)} \\[-3pt]
			&= \lim_{k  \to \infty} v_{k+1}(\xi_i) =0
		\end{align}
		Damit ist dieser Spezialfall bewiesen.\todo{Revision 7}
		\item \label{enum:kad:2} Wir nehmen nun nicht mehr an, dass $y$ wie oben existiert.
		Wir wissen aus \autoref{prop:715}, dass $\tilde{y} \in \mathcal{B}(\mathcal{H})$ mit $\tilde{y} \xi_i = \eta_i$ für $i=1,\ldots ,n$ existiert.
		Dann sind $\re \tilde{y}, \im \tilde{y} \in \mathcal{B}(\mathcal{H})_\sa$ und nach \ref{enum:kad:1} existieren $x',x'' \in A_\sa$ mit 
		\[
			\pi(x') \xi_i = \enbrace*{\re \tilde{y}} \xi_i \quad \text{ und } \quad \pi(x'') \xi_i = (\im \tilde{y})\xi_i 
		\]
		Für $x := x' + i \cdot x''$ gilt $\pi(x) \xi_i = \tilde{y}\xi_i = \eta_i$.
		\item \label{enum:kad:3}Falls die $\xi_i$ nicht notwendigerweise orthogonal, aber linear unabhängig sind, wähle eine Orthonormalbasis $\set*{\xi_1', \ldots , \xi_n'}$ für $\Span \set*{\xi_1, \ldots ,\xi_n}$.
		Falls $y \in \mathcal{B}(\mathcal{H})_\sa$ mit $y \xi_i = \eta_i$ existiert, so existiert nach \ref{enum:kad:1} ein $x \in A_\sa$ mit $\pi(x)\xi_i' = y \xi_i'$.
		Aber dann gilt auch $\pi(x)\xi_i = y \xi_i = \eta_i$ für $i=1,\ldots ,n$.
		Falls kein solches $y$ existiert, wende \ref{enum:kad:2} an.
		\item \label{enum:kad:4} Seien nun $A$ und $\pi$ unital und es existiere $u \in \mathcal{U}\enbrace*{\mathcal{B}(\mathcal{H})}$ mit $u(\xi_i)=\eta_i$.
		Wir in \ref{enum:kad:3} reduziert man auf den Fall, dass die $\xi_i$ ein Orthonormalsystem bilden.
		Sei 
		\[
			\mathcal{H}_0 := \Span \set*{\xi_1, \ldots, \xi_n, \eta_i, \ldots ,\eta_n}
		\]
		Dann existiert $u' \in \mathcal{U}\enbrace*{\mathcal{B}(\mathcal{H}_0)} \subset \mathcal{B}(\mathcal{H})$ mit $u' \xi_i = \eta_i$.\footnote{komponieren mit der Projektion auf den Span der $\xi_i$; $\set*{\xi_1, \ldots ,\xi_n}$ und $\set*{\eta_1,\ldots ,\eta_n}$ bilden Orthonormalsysteme in $\mathcal{H}_0$ und lassen sich zu Orthonormalbasen auf $\mathcal{H}_0$ ergänzen}
		Es ist $\mathcal{B}(\mathcal{H}_0) \cong M_m$ und $u' = e^{i h}$ für ein $h \in \mathcal{B}(\mathcal{H}_0)_\sa \subset \mathcal{B}(\mathcal{H})_\sa$.
		Es existiert nach \ref{enum:kad:1} ein $d \in A_\sa$ mit $\pi(d) \zeta_i = h \zeta_i$, wobei $\zeta_1,\ldots , \zeta_m$ eine Orthonormalbasis für $\mathcal{H}_0$ ist.
		Dann folgt $\pi(d)|_{\mathcal{H}_0}=h$ und $\pi(e^{id})|_{\mathcal{H}_0}= e^{ih}$.\qedhere
	\end{enumerate}
\end{beweis}

\begin{korollar}[{name=[irreduzibel impliziert algebraisch irreduzibel]},label=korr:717]
	Eine irreduzible Darstellung $(\mathcal{H},\pi)$ einer $C^*$-Algebra $A$ ist algebraisch irreduzibel.
\end{korollar}
\begin{beweis}
	Nach \autoref{satz:716} gilt $\pi(A)\xi=\mathcal{H}$ für jedes $0 \neq \xi \in \mathcal{H}$.
	Also ist $\pi(A)(\mathcal{H}_0)=\mathcal{H}$ oder $\set*{0}$ für jeden Unterraum $\mathcal{H}_0 \subset \mathcal{H}$.
\end{beweis}

\begin{definition}[label=def:718,{name=[primitives Ideal]}]
	Sei $A$ eine $C^*$-Algebra.
	Ein Ideal $J \lhd A$ heißt \bet{primitiv}\index{primitives Ideal}, falls $J= \ker \pi$ für eine irreduzible Darstellung $\pi$ ist.
	Wir schreiben $\Prim A  := \set*{J \lhd A \text{ primitiv}}$.
	Die \Index{Hülle-Kern-Topologie} bzw. \Index{Jacobson-Topologie} auf $\Prim A$ ist gegeben durch:
	Für $R \subset \Prim A$ sei
	\[
		\overline{R} = \hull\enbrace[\big]{\ker(R)} = \set*{J \in \Prim A \given \enbrace*{\bigcap_{J \in R} J} \subset J}
	\]  
\end{definition}

\begin{satz}[label=satz:719,{name=[Für $A$ kommutativ ist das Charakterspektrum homöomorph zu den primitiven Idealen]}]
	Für eine kommutative $C^*$-Algebra $A$ gilt $\Omega(A) \approx \Prim A$.
\end{satz}
\begin{beweis}
	Jeder Charakter $\chi \in \Omega(A)$ ist eine irreduzible Darstellung $\chi \colon A \to \mathbb{C} = \mathcal{B}(\mathbb{C})$.
	Wir haben also eine Abbildung 
	\mapdef{\gamma \colon\Omega(A)}{\Prim A}{\chi}{\ker \chi}{}
	$\gamma$ ist injektiv: 
	\[
		\begin{tikzcd}[sep=small]
			A \ar[rr,"\chi_1"] \ar[ddd,equal] \drar& & \mathbb{C} \ar[ddd,"\cong"]\\
			& \sfrac{A}{\ker \chi_1} \urar["\cong"'] \dar[equal] \\[-2.3pt]
			& \sfrac{A}{\ker \chi_2} \drar["\cong"] \\
			A \ar[rr,"\chi_2"'] \urar& & \mathbb{C}
		\end{tikzcd}
	\]
	Der Isomorphismus in der rechten Spalte ist Multiplikation mit $C \in \mathbb{C}$. 
	Man überlegt sich, dass $C=1$ gelten muss, sodass $\chi_1=\chi_2$ folgt.
	
	Für die Surjektivität sei $J \in \Prim A$ mit $J= \ker \pi$ für $\pi \colon A \to \mathcal{B}(\mathcal{H})$ irreduzibel.
	Da $A$ abelsch ist, gilt dann $\pi(A) \subset \pi(A)' = \mathbb{C} \cdot \ind_{\mathcal{B}(\mathcal{H})} \cong \mathbb{C}$.
	Also ist $\pi$ ein Charakter. 
	
	Für den Homöomorphismus nutze
	\mapdef{\set[\big]{Y \subset \Omega(A) \text{ abg.} }}{J \lhd A}{Y}{C_0 \enbrace[\big]{\Omega(A) \setminus Y}}{}
	und 
	\[
		C_0 \enbrace*{\Omega(A) \setminus Y} \cong \set*{f \in C_0(\Omega(A)) \given f|_Y=0} = \bigcap_{y \in Y} \set*{f \in C_0(\Omega(A)) \given f(y)=0}
	\]
	Damit kann man die Stetigkeit beweisen.
\end{beweis}
% section 7 (end)
\newpage

\section{Beispiele!} % (fold)
\label{sec:beispiele}
\emph{(Mit Ausrufezeichen)}\todo{Revision 8}

\begin{satz}[{name=[endlich dimensionale C*-Algebren sind direkte Summen von Matrixgruppen]}]
	Eine endlichdimensionale $C^*$-Algebra $A$ ist von der Form 
	\[
		M_{r_1} \oplus \ldots \oplus M_{r_s} \qquad r_1,\ldots,r_s \in \mathbb{N}
	\]
\end{satz}
\begin{beweis}
	Sei zunächst $A$ einfach und $\pi \colon A \to \mathcal{B}(\mathcal{H})$ irreduzibel.
	Für $0 \neq \xi \in \mathcal{H}$ ist
	\(
		\mathcal{H} \StackText{\ref{satz:716}}{=} \pi(A)\xi
	\)
	endlichdimensional und es gilt $\pi(A)=\mathcal{B}(\mathcal{H})$ wieder nach \autoref{satz:716}.
	Da $A$ einfach ist, muss $\pi$ bereits ein Isomorphismus sein und es folgt $\mathcal{B}(\mathcal{H}) \cong \mathcal{B}(\mathbb{C}^r)=M_r$.
	
	Falls $A$ nicht einfach ist, so enthält $A$ ein Ideal $\set*{0}\neq J \neq A$ mit minimaler Dimension.
	Dann ist $J$ eine endlichdimensionale $C^*$-Algebra und enthält keine nichttrivialen Ideale, ist also einfach.
	Mit dem ersten Teil folgt also $J \cong M_{r_1}$ für ein $r_1 \in \mathbb{N}$.
	$J$ ist insbesondere unital, sodass $J= J \cdot \ind_J = A \cdot \ind_J= \ind_J \cdot A$ gilt und damit 
	\[
		A=A \enbrace*{\ind_J+ (\ind_{A^\simm}- \ind_J)} = A  \ind_J + A \enbrace*{\ind_{A^\simm} - \ind_J} = \Underbracket{\ind_J  A \ind_J}{\cong M_{r_1}} \oplus \Underbracket{\enbrace*{\ind_{A^\simm} - \ind_J}A \enbrace*{\ind_{A^\simm} - \ind_J}}{\lhd A}
	\]
	Das rechte Ideal hat eine kleinere Dimension als $A$ und damit sind wir fertig.
\end{beweis}

\begin{beispiel}
	Für eine $C^*$-Algebra $A$ sei
	\mapdef{\iota_n^A \colon M_n(A)}{M_{n+1}(A)}{X}{\begin{psmallmatrix}
		X & 0\\
		0 & 0
	\end{psmallmatrix}}
	Schon gesehen: Für $A=\mathbb{C}$ ist $\Dlim(M_n,\iota_n^\mathbb{C}) \cong \mathcal{K} \enbrace*{\ell^2(\mathbb{N})} = \mathcal{K}$.
	Für ein beliebiges $A$ schreiben wir
	\[
		\mathcal{K} \otimes A := \dlim \enbrace[\big]{M_n(A),\iota_n^A}
	\]
	Die Notation werden wir später rechtfertigen.
	Es gilt $\mathcal{K} \cong \mathcal{K} \otimes M_r \cong \mathcal{K} \otimes \mathcal{K}$ und $\mathcal{K} \otimes M_r \cong M_r(\mathcal{K})$.
	Für die primitiven Ideale gilt
	\[
		\Prim A\approx \Prim(\mathcal{K} \otimes A)
	\]
	Zwei $C^*$-Algebren $A$ und $B$ heißen \Index{stabil isomorph}, falls $\mathcal{K} \otimes A \cong \mathcal{K} \otimes B$. 
	Tensorieren mit den Kompakten Operatoren liefert \enquote{nichtkommutativen Platz} für direkte Summen von Moduln.
	Dies liefert eine abelsche Halbgruppe und daraus erhält man mit der Grothendiek-Konstruktion eine abelsche Gruppe. 
	Das ist die Grundlage der K-Theorie.
\end{beispiel}

\begin{beispiel}
	Betrachte $A := C\enbrace*{[0,1],M_2}$ und die Unteralgebren 
	\begin{align}
		B &:= \set*{f \in A \given f(1) = \begin{psmallmatrix} \lambda & 0\\ 0 & \lambda \end{psmallmatrix}, \lambda \in \mathbb{C}}\\
		C &:= \set*{f \in A \given f(1) = \begin{psmallmatrix} \lambda_1 & 0\\ 0 & \lambda_2 \end{psmallmatrix}, \lambda_1, \lambda_2 \in \mathbb{C}} \\
		D &:= \set*{f \in A \given f(1)= \begin{psmallmatrix} \lambda_1 & 0\\ 0 & \lambda_2 \end{psmallmatrix}, f(0)= \begin{psmallmatrix} \lambda_1 & 0\\ 0 & \lambda_1 \end{psmallmatrix}, \lambda_1, \lambda_2 \in \mathbb{C}}
	\end{align}
	Wir können und $\Prim A$ als Charaktere auf $A$ vorstellen und in diesem Fall sogar als Intervall, wobei die Punkte Evaluierungen $\ev_f \colon A \to M_2$ entsprechen.
	In $\Prim B$ gilt zusätzlich $\ev_1 \colon B \to \mathbb{C} \cdot \ind_{M_2} \cong \mathbb{C}$.
	Für $\Prim C$ enthält am rechten Rand des Intervalls zwei Punkte, einmal die Auswertung, die $\lambda_1$ ausgibt, und die, die $\lambda_2$ ausgibt.
	Damit ist $\Prim C$ nicht hausdorffsch, aber immerhin $T_0$.
	Für $\Prim D$ erhält man einen Kreis, der an einer Stell einen verdoppelten Punkt hat, der nur zu einer Seite hin offen ist.
\end{beispiel}

\begin{beispiel}
	Sei $A := \set[\big]{f \in C_0\enbrace*{[0,\infty),\mathcal{K}} \given f(t) \in \iota_{n,\infty}(M_n) \text{ für }t \in [n-1,n)}$.
	\missingfigure{hier kommt eine etwas esoterische Zeichnung hin \ldots }
	$\Prim A \approx [0,\in )$ ist hausdorffsch.
	Für jedes $t_0 \in \Prim A$ existieren eine Umgebung $U$ von $t_0$ und $f \in A$, sodass $f(t)$ eine Projektion von Rang 1 ist für alle $t \in U$.
	$A$ hat eine \enquote{stetige Spur}.
	Für $A,B$ mit stetiger Spur und $\Prim A \approx \Prim B$, so gilt 
	\[
		\mathcal{K} \otimes A \cong \mathcal{K} \otimes B \iff \Underbracket{H^3(\Prim A)}{\ni \delta(A)}\cong \Underbracket{H^3(\Prim B)}{\ni\delta(B)} 
	\]
	Dixmier-Douady-Klasse
\end{beispiel}

\begin{beispiel}
	\leavevmode
	\begin{enumerate}[(i)]
		\item $M_r$ ist erzeugt von den Matrixeinheiten $e_{ij} \in M_r$, $i,j=1,\ldots ,r$.
		Die $e_{ij}$ erfüllen Relationen $R_{M_r}=\set{e_{ij} e_{kl}=\delta_{jk} e_{il}, e_{ij}^*= e_{ji} \given  i,j,k,l \in 1,\ldots ,r}$.
		Falls $A$ eine weitere $C^*$-Algebra ist mit Elementen  $f_{ij}$, $i,j=1,\ldots,r$, welche die Relationen $R_{M_r}$ erfüllen, so existiert genau ein $^*$-Homomorphismus $\pi \colon M_r \to M$ mit $\pi(e_{ij})=f_{ij}$.
		Die Erzeuger $e_{ij}$ mit den Relationen $R_{M_r}$ erfüllen also eine universelle Eigenschaft.
		($M_r$ einfach $\implies \pi$ ist frei, falls $f_{ij}\neq 0$)
		\item Sei $X=\set*{x_\gamma \given \gamma \in \Gamma}$ eine Menge.
		Eine Relation mit Erzeugern in $X$ ist ein Ausdruck der Form $p(x_{\gamma_1}, \ldots ,x_{\gamma_n}, x_{\gamma_1}^*, \ldots , x_{\gamma_n}^*)=0$, wo $p$ ein Polynom in $2n$ nichtkommutierenden Variablen ist.
		Allgemeiner betrachtet man auch Relationen $[p,\varepsilon]$ der Form $\norm*{p(x_1,\ldots ,x_n,x_1^*,\ldots ,x_n^*)} \le \varepsilon$, $\varepsilon>0$.
		
		Eine Darstellung von $(X,R)$ (wo $R$ eine Menge von Relationen ist) ist ein $^*$-Homomorphismus $\pi \colon ^*\text{-}\Alg(X) \to \mathcal{B}(\mathcal{H})$, der jede Relation erfüllt, das heißt 
		\[
			\norm*{p\enbrace[\big]{\pi(x_{\gamma_1}) , \ldots , \pi(x_{\gamma_n}), \pi{x_{\gamma_1}}^* ,\ldots , \pi(x_{\gamma_n})^*}} \le \varepsilon
		\]
		$(X,R)$ heißt zulässig, falls es für jedes $x_j$ ein $c_\gamma \ge 0$ gibt mit $\norm*{\pi(x_\gamma)} \le c_\gamma$ für jede Darstellung $\pi$.
		Für zulässige $(X,R)$ definieren wir eine Halbnorm auf $^*$-$\Alg(X)$ durch 
		\[
			\norm*{x} := \sup_{\pi \text{ Darst.}} \norm*{\pi(x)}_{\mathcal{B}(\mathcal{H})}
		\]
		Definiere $C^*(X\mid R) := \overline{^*\text{-}\Alg(X)}^{\norm*{\cdot}}$ (Vervollständigung bezüglich einer Halbnorm).
		\item Beispiele:
		\begin{enumerate}[a)]
			\item $M_r \cong C^* \enbrace*{\set*{e_{ij}} \mid R_{M_r}}$.
			\item 
			\begin{align}
				C^* \enbrace*{\set*{\ind,u} \mid uu^*=\ind, u^*u=\ind, \ind \ind=\ind, \ind u=u, u \ind=u,\ind^*=\ind} \cong C \enbrace*{S^1} = C^*(\id_{S^1})
			\end{align}
			\item $A_\theta := C^* \enbrace*{\set*{u,v,1}\mid u,v \text{ unitär},uv=e^{2 \pi i \theta} vu }$ für $\theta \in [0,1]$ bezeichnet man als \Index{Drehungsalgebra}.
			Für $\theta=0$ erhält man die stetigen Funktionen auf dem Torus.
			\item Die Relation $x=x^*$ ist nicht zulässig (da die Norm den Wert unendlich haben kann). Dies lässt sich aber durch eine zusätzliche Relation reparieren
			\item Die Relationen $x=x^*$, $p=p^*$ und $xp-px=i \cdot \ind$ sind auch nicht zulässig, denn sie lassen sich nur mit unbeschränkten Operatoren darstellen.
			\item Es gilt $C^* \enbrace*{\set*{p,q} \mid p=p^*, p=p^2, f=f^*,f=f^2} \cong \set*{f \in C([0,1],M_2) \given f(0)=\diag(\lambda,\mu), f(1)= \diag(\kappa,0), \lambda,\mu, \kappa \in \mathbb{C}}$
			\missingfigure{wieder eine abstrakte Zeichnung}
			\item Sei $G$ eine diskrete Gruppe. 
			Die volle Gruppen-$C^*$-Algebra ist gegeben durch
			\[
				C^*(G) := C^* \enbrace*{X\mid R}\quad  \text{ mit } \quad X=\set*{u_g \given g \in G}  \quad R= \set*{u_g \text{ unitär}, u_g^*=u_{g^{-1}}, u_g u_h = u_{gh}}
			\]
			Alternativ kann man auch den Gruppenring betrachten:
			\[
				\mathbb{C}[G] := \set[\bigg]{\sum_{g \in G}\alpha_g \cdot g \given \text{ nur endlich viele } \alpha_g\neq 0} = \set[\big]{\alpha \colon G \to C \given \alpha \text{ hat endlichen Träger} }
			\]
			mit Faltung, das heißt $(\alpha * \beta)(h) = \sum_{g \in G} \alpha(g) \beta(g^{-1}h)$ und Involution $\alpha^*(g)=\overline{\alpha(g^{-1})}$.
			Die unitären Darstellungen $\tilde{\pi} \colon G \to \mathcal{B}(\mathcal{H})$ stehen in Bijektion zu den $^*$-Darstellungen $\pi \colon \mathbb{C}[G] \to \mathcal{B}(\mathcal{H})$.
			Die $C^*$-Norm auf $\mathbb{C}[G]$ ist gegeben durch
			\[
				\norm*{x} = \sup_{\substack{\tilde{\pi} \text{ unitäre} \\ \text{Darst.}}} \norm*{\pi(x)}_{\mathcal{B}(\mathcal{H})}
			\]
			Es folgt $C^*(G) \cong \overline{\mathbb{C}[G]}^{\norm*{\cdot}}$. $\norm*{\cdot }$ ist tatsächlich eine Norm, denn $\mathbb{C}[G]$ besitzt eine injektive Darstellung:
		\end{enumerate}
	\end{enumerate}
\end{beispiel}

\begin{beispiel}
	Sei $G$ wieder eine diskrete Gruppe
	\[
		\ell^2(G) = \set*{\xi \colon G \to \mathbb{C} \given \norm*{\xi}_2 := \enbrace*{\sum\nolimits_G \abs*{\xi(g)}^2 }^{\sfrac{1}{2}} < \infty}
	\]
	Mit der Orthonormalbasis $\set*{\delta_g \given g \in G}$. Dann ist die linksreguläre Darstellung $\tilde{\lambda} \colon G \to \mathcal{B} \enbrace*{\ell^2(G)}$ gegeben durch $\tilde{\lambda}(g)(\delta_h) = \delta_{gh}$, das heißt $\enbrace*{\tilde{\lambda}(g)(\xi)}(h) = \xi \enbrace*{g ^{-1}h}$.
	$\tilde{\lambda}$ ist eine unitäre Darstellung und mit der Bijektion erhält man eine injektive $^*$-Darstellung $\lambda \colon \mathbb{C}[G] \to \mathcal{B}(\mathcal{H})$.
	Wir definieren die reduzierte Gruppen-$C^*$-Algebra durch
	\[
		C^*_{\red}(G) := C^* \enbrace*{\tilde{\lambda}(G)}= C^* \enbrace*{\lambda \enbrace*{\mathbb{C}[G]}} \subset \mathcal{B} \enbrace*{\ell^2(G)}
	\]
	und weiter $L(G) := \overline{C^*_{\red}(G)}^\wo$. Diese wird als Gruppen-von-Neumann-Algebra bezeichnet.
	Man kann zeigen, dass $C^*(G) \twoheadrightarrow C^*_{\red}(G)$ genau dann ein Isomorphismus ist, wenn $G$ amenabel ist.
	Die freie Gruppe in $n$ variablen $\mathcal{F}_n$ ist nicht amenabel und man kann zeigen
	\[
		C^*_{\red}(\mathcal{F}_n) \not\cong C^*_{\red}(\mathcal{F}_m) \quad m \neq n
	\]
	Ein offenes Problem ist die analoge Frage 
	\[
		\text{¿ } L(\mathcal{F}_n) \cong L (\mathcal{F}_m) \text{ ?}  
	\]
\end{beispiel}

\begin{beispiel}
	Sei $X$ ein kompakter Hausdorffraum und $\overline{\alpha}$ ein Homöomorphismus $X\to X$.
	Dann erhält man $\alpha \in \Aut(C(X))$ und man kann durch
	\[
		\alpha^n(f) = f \circ \overline{\alpha}^{-n} \qquad n \in \mathbb{Z}
	\]
	Sei $\pi \colon C(X) \to \mathcal{B}(\mathcal{H})$ eine treue Darstellung.
	Setzt man $\tilde{H} := \ell^2(\mathbb{Z},\mathcal{H})$, so erhält man $\tilde{\pi} \colon C(X) \to \mathcal{B}(\tilde{\mathcal{H}})$ durch
	\[
		\tilde{\pi}(f) = \begin{pmatrix}
			\ddots & & && \\
			& \pi(\alpha^{-1} f) & & &\\
			& & \pi(f) & &\\
			& & & \pi(a f) & \\
			& & & & \ddots
		\end{pmatrix} \in \mathcal{B}(\tilde{\mathcal{H}})
 	\]
	Sei $u \in \mathcal{B}(\tilde{\mathcal{H}})$ der bilaterale Shift.
	\[
		C(X) \rtimes_\alpha \mathbb{Z} := C^* \enbrace*{\tilde{\pi}(C(X)),u} \in \mathcal{B}\enbrace*{\tilde{\mathcal{H}}}
	\]
	$C(S^1) \rtimes_{\rot_\theta} \mathbb{Z} \cong A_\theta$.
\end{beispiel}
% section beispiele (end)











\cleardoubleoddemptypage
\pagenumbering{Alph}
\setcounter{page}{1}
\cleardoubleoddemptypage
\appendix

\section{Anhang} % (fold)
\label{sec:anhang}
Hier findet man einige ausgewählte Beweise und Ergänzungen, auf die der Vorlesung hingewiesen wurde und die teilweise im Rahmen von Übungsaufgaben bearbeitet wurden.
Sie stammen also auch hauptsächlich aus den Übungen und deswegen sei hier noch einmal darauf hingewiesen, dass für die Richtigkeit keine Garantie übernommen wird!


\subsection{Beweis von \autoref{satz:37} (vii)} % (fold)
\label{sub:beweis_37_vii}
\emph{Dies war außerdem die Übungsaufgabe 1 b) von Blatt 5.} \smallskip\\
Wenn $A$ unital ist und $a \in A$ positiv und invertierbar, dann folgt aus $a \le b$ stets, dass auch $b$ invertierbar ist und $0 \le b^{-1} \le a^{-1}$ gilt.
\begin{beweis}
	Da $a$ invertierbar ist, ist $0$ nicht im Spektrum von $a$ enthalten.
	Da $\sigma(a)$ kompakt ist, ist die äquivalent dazu, dass ein $\lambda>0$ existiert mit $\sigma(a) \subset [\lambda,\infty)$.
	Weiter gilt
	Wir haben damit folgende Äquivalenzen
	\[
		\sigma(a) \subset [\lambda,\infty) \iff \sigma(a-\lambda) \subset [0,\infty) \iff a- \lambda\ge 0 \iff a \ge \lambda
	\]
	Da $b\ge a \ge \lambda$, gilt also auch $0 \notin \sigma(b)$ und $b$ ist invertierbar.
	
	Das $\sigma(a^{-1}) = \set*{\lambda^{-1} \given \lambda \in  \sigma(a)}$, sind $a^{-1}$ und $b ^{-1}$ wieder positiv.
	Insbesondere existieren die Quadratwurzeln von den Inversen. 
	Mit \autoref{satz:37} (v) folgt aus $a \le b$ also 
	\[
		\ind=a^{-\sfrac{1}{2}} a a^{- \sfrac{1}{2}} \le a^{-\sfrac{1}{2}} b a^{-\sfrac{1}{2}} =: c
	\]
	Nach Gelfand folgt aus $c \ge \ind$, dass $c$ invertierbar ist und $c^{-1} \le \ind$ gilt. Wieder mit \autoref{satz:37} (v) erhalten wir
	\[
		c^{-1} = a^{\sfrac{1}{2}}b^{-1}a^{\sfrac{1}{2}} \le \ind \implies b^{-1} \le \enbrace{a^{\sfrac{1}{2}}}^{-1} \enbrace{a^{\sfrac{1}{2}}}^{-1} =a^{-1} \qedhere
	\]
\end{beweis}
% subsection beweis_von_autoref_satz_37_vii (end)


\subsection{Alternativer Beweis von \autoref{lem:311}} % (fold)
\label{sub:alternativer_beweis_lemma_311}
Sei $A$ eine $C^*$-Algebra. Dann ist $\Lambda := \set*{a \in A_+ \given \norm*{a} < 1}$ gerichtet bezüglich \enquote{$\le$}.
\begin{beweis}[{\cite[S.~78]{Murphy}}]
	Wenn $a \in A_+$, dann ist $\ind +a$ invertierbar in $A^\simm$, da $0$ nicht im Spektrum enthalten sein kann.
	Weiter gilt $a(\ind +a)^{-1} = \ind - (\ind +a)^{-1}$.
	Wir zeigen nun
	\[
		a,b \in A_+ \text{ und } a \le b \implies a(\ind+a)^{-1} \le b (\ind +b)^{-1} 
	\]
	Aus $a \le b$ folgt nämlich $\ind +a \le \ind +b$ und damit $(\ind +b)^{-1} \le (\ind +a)^{-1}$ nach \autoref{satz:37} (vii).
	Daraus erhalten wir die Ungleichung $\ind - (\ind +a)^{-1} \le \ind - \enbrace*{\ind +b}^{-1}$, nach der Gleichung für $a(\ind +a)^{-1}$ oben also die Behauptung.
	Durch Betrachtung des Gelfandisomorphismus für $C^*(a, \ind)$ sieht man, dass $a(\ind +a)^{-1} \in \Lambda$ ist.
	
	Seien nun $a,b \in \Lambda$ beliebig.
	Dann setzen wir $a' := a \enbrace*{\ind -a}^{-1}$ und $b' := b \enbrace*{\ind -b}^{-1}$. Hier überlegt man sich die Invertierbarkeit und Positivität entweder auch mit Gelfand oder der von-Neumann-Reihe (siehe \ref{satz15_neumann}).
	Genauso ist auch $\ind + a'$ invertierbar und damit auch $\ind + a' +b'$ (wieder mit \ref{satz:37}). 
	Dann gilt nach Gelfand
	\[
		c := (a' + b') \enbrace*{\ind + a' + b'}^{-1} \in \Lambda
	\]
	Da offensichtlich $a' \le a' + b'$ gilt, folgt aus der oben gezeigten Implikation $a'(\ind +a')^{-1} \le c$.
	Es gilt nun
	\[
		a'(\ind +a')^{-1} = a (\ind -a)^{-1} \enbrace*{\ind + a(\ind-a)^{-1}}^{-1} \StackText{\ref{bem:13}}{=} a (\ind -a)^{-1} \enbrace*{(\ind-a)^{-1}}^{-1} = a
	\]
	womit wir $a \le c$ gezeigt haben. 
	Die Relation $b \le c$ folgt genauso, also ist $\Lambda$ gerichtet.
\end{beweis}
% subsection alternativer_beweis_autoref_lemma_311 (end)


\subsection{Konstruktion der eindeutigen Abbildung für den Limes von $C^*$-Algebren} % (fold)
\label{sub:konstruktion_limes}
\emph{Mit der Notation aus \autoref{satz:ger_lim}.}\smallskip\\
Wir definieren zunächst eine geeignete Abbildung 
\mapdef{\theta \colon C}{B}{(a_n)_{\mathbb{N}}}{\sum_\mathbb{N} \psi_n \enbrace[\big]{a_n - \varphi_{n-1}(a_{n-1})}}{}
Diese Abbildung ist wohldefiniert, denn nach Definition von $C$ ist die Summe endlich.
Weiter gilt offensichtlich $\theta \circ \varphi_n' = \psi_n$.
Man muss nun noch zeigen, dass $\theta$ die Elemente in $\bigoplus_{\mathbb{N}} A_n$ auf $0$ abbildet, um zu zeigen, dass man $\theta$ auf Repräsentanten in $q(C)$ anwenden kann.\todo{jetzt wo ich anfange das aufzuschreiben, scheint mir das ganze gar nicht mehr so klar \ldots}
% subsection konstruktion_der_eindeutigen_abbildung_fur_den_limes_von_c_algebren (end)


\subsection{Die GNS-Darstellung ist zyklisch} % (fold)
\label{sub:gns_zyklisch}
\emph{Dies war Bestandteil von Blatt 8, Aufgabe 2 und komplettiert die Ausführungen zur GNS-Darstellung in \autoref{gns} und \autoref{satz:512}).} \smallskip \\
Wir zeigen, dass der Grenzwert $\xi_\varphi = \lim_\lambda e_\lambda + N_\varphi$ für eine beliebige approximative Eins $(e_\lambda)_\Lambda \subset A^1_+$ existiert, eindeutig ist und ein zyklischer Vektor für die GNS-Darstellung ist (vgl. \autoref{gns}).
Weiter wird gezeigt, dass folgende Identität für $a \in A$ gilt:
\begin{equation}
	\varphi(a) = \skal[\big]{\xi_\varphi}{\pi_\varphi(a)\xi_\varphi}
	\label{eq:zykl} \tag{\dag}
\end{equation}
\begin{beweis}
	Da $(e_\lambda)_\Lambda$ aufsteigend ist, gilt $e_\mu - e_\lambda \ge 0$ für $\lambda \le \mu$. 
	Weiter gilt $e_\mu - e_\lambda \le \ind_{A^\simm}$.
	Nach \autoref{satz:37} (v) mit $c=(e_\mu -e_\lambda)^{\sfrac{1}{2}}$ folgt $(e_\mu -e_\lambda)^2 \le e_\mu -e_\lambda$.
	Es gilt nun
	\[
		\norm*{e_\mu - e_\lambda + N_\varphi}^2 = \varphi \enbrace[\big]{(e_\mu-e_\lambda)^2} \le \varphi(e_\mu -e_\lambda)= \varphi(e_\mu)- \varphi(e_\lambda)
		\grenzw{\mbox{\scriptsize\ref{satz:54}}} 0
	\]
	da postive Operatoren Ungleichungen erhalten (warum?).
	Also ist $(e_\lambda +N_\varphi)_\Lambda$ ein Cauchy-Netz und konvergiert gegen $\xi_\varphi$, da $\mathcal{H}_\varphi$ vollständig ist.
	Um zu zeigen, dass $\xi_\varphi$ zyklisch ist, reicht es $\pi_\varphi(a) \xi_\varphi = a + N_\varphi$ für $a \in A$ zu zeigen. Es gilt
	\[
		\pi_\varphi(A) \xi_\varphi = \lim_\lambda \dot{\pi}_\varphi(a) \enbrace*{e_\lambda +N_\varphi} = \lim_\lambda \enbrace*{a e_\lambda + N_\varphi} = a+ N_\varphi
	\]
	Wegen Linearität genügt es nun \eqref{eq:zykl} für positive Elemente $a^*a$ zu zeigen:\marginnote{hier ist es wichtig, dass das Skalarprodukt in der zweiten Koordinate linear ist!} 
	Da $\xi_\varphi$ zyklisch ist, gilt
	\[
		\varphi(a^*a) = \skal*{a+ N_\varphi}{a+N_\varphi} = \skal[\big]{\pi_\varphi(a) \xi_\varphi}{\pi_\varphi(a)\xi_\varphi} = \skal[\big]{\xi_\varphi}{\pi_\varphi(a^*a)\xi_\varphi}\qedhere
	\]
\end{beweis} 
% subsection gns_konstruktion_ist_zyklisch (end)

\subsection{Monotones Netz von Projektionen} % (fold)
\label{sub:mon_proj}
\emph{Dies war Bestandteil von Blatt 9, Aufgabe 4 und komplettiert \autoref{satz:62}.} \smallskip \\
Sei $(p_\lambda)_\Lambda$ ein aufsteigendes/absteigendes Netz von Projektionen. 
Sei $q$ die Projektion auf $\overline{\bigcup_\lambda p_\lambda(\mathcal{H})}$ und $r$ die Projektion auf $\bigcap_\lambda p_\lambda(\mathcal{H})$. Dann gilt
\begin{enumerate}[(i),itemsep=0pt]
	\item Falls $(p_\lambda)_\Lambda$ aufsteigend ist, dann gilt $p_\lambda \grenzw{\so} q$.
	\item Falls $(p_\lambda)_\Lambda$ absteigend ist, dann gilt $p_\lambda \grenzw{\so} r$.
\end{enumerate}
\begin{beweis}
	\leavevmode
	\begin{enumerate}[(i)]
		\item Wir müssen zeigen, dass $p_\lambda \xi\to q \xi$ für alle $\xi \in \mathcal{H}$.
		Für Projektionen gilt\todo{Beweis davon könnte man von Thanh oder \citeauthor{Murphy} übernehmen}
		\[
			p_\lambda \le p_\mu \iff p_\lambda(\mathcal{H}) \subset p_\mu(\mathcal{H})
		\]
		Jedes $\xi \in \mathcal{H}$ lässt sich eindeutig schreiben als $\xi=\xi_1 + \xi_2$ mit $\xi_1 \in \overline{\bigcup_\lambda p_\lambda(\mathcal{H})}$ und $\xi_2 \in \bigcap_\lambda p_\lambda(\mathcal{H})^\bot$.
		Wähle eine Folge $(x_n)_\mathbb{N} \to \xi_1$.
		Sei nun $\varepsilon>0$. Wähle $n_0 \in \mathbb{N}$ mit $\norm*{\xi_1-x_n}< \sfrac{\varepsilon}{2}$ für $n\ge n_0$ und passend dazu $\lambda_0 \in \Lambda$ mit $x_{n_0} \in p_{\lambda_0}\enbrace{\mathcal{H}}$.
		Dann ist $x_{n_0} \in p_\lambda \enbrace*{\mathcal{H}}$ für $\lambda\ge \lambda_0$ und es gilt
		\begin{align}
			\norm*{p_\lambda \xi -q \xi} = \norm*{p_\lambda \xi_1 + p_\lambda \xi_2 - \xi_1} = \norm*{p_\lambda \xi_1 - \xi_1} &= \norm*{p_\lambda \enbrace*{\xi_1 - x_{n_0} + x_{n_0}} - \xi_1} \\
			&\le \norm*{p_\lambda \enbrace*{\xi_1 - x_{n_0}}} + \norm*{x_{n_0} - \xi_1} \\
			&\le \Underbracket{\norm*{p_\lambda}}{=1} \cdot \norm*{\xi_1 - x_{n_0}} + \norm*{x_{n_0} - \xi_1} \\
			&< \frac{\varepsilon}{2} + \frac{\varepsilon}{2} = \varepsilon
		\end{align}
		Damit ist $p_\lambda \grenzw{\so} q$ gezeigt.
		\item Da $p_\lambda$ ein Projektion ist, ist $\ind - p_\lambda$ die Projektion auf $p_\lambda(\mathcal{H})^\bot$.
		Weiter ist dann $(\ind - p_\lambda)_\Lambda$ ein aufsteigendes Netz von Projektionen und nach (i) konvergiert dies stark gegen die Projektion auf $\overline{\bigcup_\lambda p_\lambda(\mathcal{H})^\bot}$.
		Damit konvergiert dann aber auch $p_\lambda$ stark gegen die Projektion auf
		\[
			\overline{\bigcup_\lambda p_\lambda(\mathcal{H})^\bot}^\bot = \bigcap_\lambda p_\lambda(\mathcal{H})
		\]
		Also gilt $p_\lambda \grenzw{\so} r$ wie gewünscht.\qedhere
	\end{enumerate}
\end{beweis}
% subsection mon_proj (end)

\subsection{Normabschätzung im Beweis der Polarzerlegung, \ref{satz:611}} % (fold)
\label{sub:normabschatzung_im_beweis_der_polarzerlegung}
Für $x \in A \subset \mathcal{B}(\mathcal{H})$ setzen wir $u_n := x\enbrace[\big]{\sfrac{1}{n}\cdot \ind+ \abs*{x}}^{-1}$.
Dann ist
\[
	u_n^* u_n = \enbrace[\big]{\sfrac{1}{n}\cdot \ind + \abs*{x}}^{-1} x^* x \enbrace[\big]{\sfrac{1}{n}\cdot \ind + \abs*{x}}^{-1}
\]
Weiter ist $\enbrace*{\sfrac{1}{n}\cdot \ind + \abs*{x}}^2 = \sfrac{1}{n^2} \cdot \ind + \sfrac{2}{n} \cdot \abs*{x} + x^*x$ und damit offensichtlich $x^*x \le \enbrace*{\sfrac{1}{n}\cdot \ind + \abs*{x}}^2$.
Eine Anwendung von \autoref{satz:37}(v) mit $c=\enbrace[\big]{\sfrac{1}{n}\cdot \ind+ \abs*{x}}^{-1}$ liefert dann $u^*_n u_n \le \ind$.
Insgesamt folgt also mit der $C^*$-Gleichung $u_n \in A^1$ für $n \in \mathbb{N}$.
% subsection normabschatzung_im_beweis_der_polarzerlegung_ref (end)

\subsection{Aufgabe 5 von Blatt 12} % (fold)
\label{sub:aufg5_12}
\emph{Diese Übungsaufgabe komplettiert \autoref{korr:714}.}\smallskip\\
Sei $B \subset A$ eine $C^*$-Unteralgebra und $\rho \in PS(B)$ ein reiner Zustand.
Dann ist $F := \set*{\varphi \in S(A) \given \varphi|_B = \rho}$ eine Seite von $S := \set*{\psi \in A^* \given \psi \text{ positiv}, \norm*{\psi}\le 1}$.
\begin{beweis}
	Wie in \autoref{satz:711} erläutert, ist $S$ konvex, $\w^*$-kompakt und nichtleer ($0 \in S$).
	Es seien nun $\theta \in F$, $\varphi_1, \varphi_2 \in S$ und $t \in (0,1)$ mit\marginnote{da $\norm*{\theta}=1$ und $\norm*{\varphi_i} \le 1$, folgt $\varphi_i \neq 0$}
	\[
		\theta = t \cdot \varphi_1 + (1-t) \cdot \varphi_2
	\]
	Nach Definition der Operatornorm gilt $\norm*{\varphi_i|_B} \le \norm*{\varphi_i} \le 1$ für $i=1,2$.
	Einschränken auf $B$ liefert
	\begin{equation}
		\rho = \theta\big|_B = t \cdot \varphi_1\big|_B + (1-t) \cdot \varphi_2\big|_B \label{eq:sub:aufg5_12} \tag{\#}
	\end{equation}
	Betrachtet man diese Gleichung in der Norm erhält man
	\[
		1 = \norm*{\rho} = \norm*{t \cdot \varphi_1\big|_B + (1-t) \cdot \varphi_2\big|_B} \le t \cdot \norm*{\varphi_1\big|_B} + (1-t) \cdot \norm*{\varphi_2\big|_B} \le 1
	\]
	Es folgt $\norm*{\varphi_1|_B} = \norm*{\varphi_2|_B}=1$ und damit auch $\norm*{\varphi_1}=\norm*{\varphi_2}=1$. $\varphi_1$ und $\varphi_2$ sind also Zustände.
	Aus \eqref{eq:sub:aufg5_12} folgt weiter $t \cdot \varphi_1|_B \le \rho$ und $(1-t) \varphi_2|_B \le \rho$.
	Da $\rho$ ein reiner Zustand ist, existieren $\lambda_1, \lambda_2 \in [0,1]$ mit
	\[
		t \cdot \varphi_1\big|_B = \lambda_1 \cdot \rho \quad \text{ und } \quad (1-t) \cdot \varphi_2\big|_B = \lambda_2 \cdot \rho 
	\]
	Da $\norm*{\rho}=\norm*{\varphi_1|_B}=\norm*{\varphi_2|_B}=1$ und $t, (1-t) \in (0,1)$, folgt $\varphi_1|_B = \varphi_2|_B=\rho$.
	Damit sind $\varphi_1$ und $\varphi_2$ bereits in $F$ enthalten und die Behauptung folgt.
\end{beweis}
% subsection aufg5_12 (end)


\subsection{Das Bild irreduzibler Darstellungen ist stark dicht} % (fold)
\label{sub:das_bild_irreduzibler_darstellungen}
\emph{Dies war ein kleiner Exkurs im Beweis von \hyperref[satz:716]{Kadisons Transitivitätssatz, \ref*{satz:716}}. 
Tatsächlich ist die Aussage bezüglich der Dichtheit auch eine weitere Äquivalenz zu \autoref{satz:73} (i).}\smallskip \\
Sei $(\mathcal{H},\pi)$ eine irreduzible Darstellung einer $C^*$-Algebra $A$.
Dann ist $\pi(A)$ $\so$-dicht in $\mathcal{B}(\mathcal{H})$
\emph{}
\begin{beweis}
	Da $(\mathcal{H},\pi)$ irreduzibel ist, folgt mit \autoref{satz:73}, dass $\pi(A)'= \mathbb{C} \cdot \ind_{\mathcal{B}(\mathcal{H})}$, also $\pi(A)''=\mathcal{B}(\mathcal{H})$.
	Falls $\ind_{\mathcal{B}(\mathcal{H})} \in \pi(A)$, so ist nach \autoref{lem:66}
	\[
		\overline{\pi(A)}^{\raisemath{-1pt}{\so}}= \pi(A)''=\mathcal{B}(\mathcal{H})
	\]
	Im nichtunitalen Fall betrachten wir eine approximative Eins $(e_\lambda)_\Lambda$ für $A$.
	Dann konvergiert $\pi(e_\lambda)$ gegen eine Projektion $p \in \mathcal{B}(\mathcal{H})$ und $p \mathcal{H} \subset \mathcal{H}$ ist ein abgeschlossener $\pi(A)$-invarianter Unterraum.
	Da $\pi$ irreduzibel ist, ist $p \mathcal{H}=\mathcal{H}$ und $p=\ind_{\mathcal{B}(\mathcal{H})}$.
	Weiter ist 
	\[
		\ind_{\mathcal{B}(\mathcal{H})} \in  \overline{\pi(A)}^\so = \enbrace[\big]{\overline{\pi(A)}^\so}'' \supseteq \pi(A)'' = \mathcal{B}(\mathcal{H}) \qedhere
	\]
	% \[
	% 	\begin{tikzcd}[sep=small]
	% 		\ind_{\mathcal{B}(\mathcal{H})} \in  \overline{\pi(A)}^\so \rar[equal] & \enbrace*{\overline{\pi(A)}^\so}''  \rar[equal] & \mathcal{B}(\mathcal{H}) \\
	% 		\mbox{ }\hspace{1em}\pi(A) \uar[phantom,sloped,"\subset","\text{\tiny dicht}" yshift=-.7em] \rar[phantom,sloped,description,"\subset"] & \pi(A)'' \uar[phantom,sloped,description,"\subset"] \urar[equal]
	% 	\end{tikzcd}
	% \]
\end{beweis}
% subsection das_bild_irreduzibler_darstellungen (end)

\printindex
\printbibliography
\listoffigures
\todototoc
\listoftodos[To-do's und andere Baustellen]
\end{document}